% Options for packages loaded elsewhere
\PassOptionsToPackage{unicode}{hyperref}
\PassOptionsToPackage{hyphens}{url}
\PassOptionsToPackage{dvipsnames,svgnames,x11names}{xcolor}
\documentclass[
  11pt,
  openany]{book}
\usepackage{xcolor}
\usepackage[margin=1in]{geometry}
\usepackage{amsmath,amssymb}
\setcounter{secnumdepth}{5}
\usepackage{iftex}
\ifPDFTeX
  \usepackage[T1]{fontenc}
  \usepackage[utf8]{inputenc}
  \usepackage{textcomp} % provide euro and other symbols
\else % if luatex or xetex
  \usepackage{unicode-math} % this also loads fontspec
  \defaultfontfeatures{Scale=MatchLowercase}
  \defaultfontfeatures[\rmfamily]{Ligatures=TeX,Scale=1}
\fi
\usepackage{lmodern}
\ifPDFTeX\else
  % xetex/luatex font selection
\fi
% Use upquote if available, for straight quotes in verbatim environments
\IfFileExists{upquote.sty}{\usepackage{upquote}}{}
\IfFileExists{microtype.sty}{% use microtype if available
  \usepackage[]{microtype}
  \UseMicrotypeSet[protrusion]{basicmath} % disable protrusion for tt fonts
}{}
\makeatletter
\@ifundefined{KOMAClassName}{% if non-KOMA class
  \IfFileExists{parskip.sty}{%
    \usepackage{parskip}
  }{% else
    \setlength{\parindent}{0pt}
    \setlength{\parskip}{6pt plus 2pt minus 1pt}}
}{% if KOMA class
  \KOMAoptions{parskip=half}}
\makeatother
\usepackage{color}
\usepackage{fancyvrb}
\newcommand{\VerbBar}{|}
\newcommand{\VERB}{\Verb[commandchars=\\\{\}]}
\DefineVerbatimEnvironment{Highlighting}{Verbatim}{commandchars=\\\{\}}
% Add ',fontsize=\small' for more characters per line
\newenvironment{Shaded}{}{}
\newcommand{\AlertTok}[1]{\textcolor[rgb]{1.00,0.00,0.00}{\textbf{#1}}}
\newcommand{\AnnotationTok}[1]{\textcolor[rgb]{0.38,0.63,0.69}{\textbf{\textit{#1}}}}
\newcommand{\AttributeTok}[1]{\textcolor[rgb]{0.49,0.56,0.16}{#1}}
\newcommand{\BaseNTok}[1]{\textcolor[rgb]{0.25,0.63,0.44}{#1}}
\newcommand{\BuiltInTok}[1]{\textcolor[rgb]{0.00,0.50,0.00}{#1}}
\newcommand{\CharTok}[1]{\textcolor[rgb]{0.25,0.44,0.63}{#1}}
\newcommand{\CommentTok}[1]{\textcolor[rgb]{0.38,0.63,0.69}{\textit{#1}}}
\newcommand{\CommentVarTok}[1]{\textcolor[rgb]{0.38,0.63,0.69}{\textbf{\textit{#1}}}}
\newcommand{\ConstantTok}[1]{\textcolor[rgb]{0.53,0.00,0.00}{#1}}
\newcommand{\ControlFlowTok}[1]{\textcolor[rgb]{0.00,0.44,0.13}{\textbf{#1}}}
\newcommand{\DataTypeTok}[1]{\textcolor[rgb]{0.56,0.13,0.00}{#1}}
\newcommand{\DecValTok}[1]{\textcolor[rgb]{0.25,0.63,0.44}{#1}}
\newcommand{\DocumentationTok}[1]{\textcolor[rgb]{0.73,0.13,0.13}{\textit{#1}}}
\newcommand{\ErrorTok}[1]{\textcolor[rgb]{1.00,0.00,0.00}{\textbf{#1}}}
\newcommand{\ExtensionTok}[1]{#1}
\newcommand{\FloatTok}[1]{\textcolor[rgb]{0.25,0.63,0.44}{#1}}
\newcommand{\FunctionTok}[1]{\textcolor[rgb]{0.02,0.16,0.49}{#1}}
\newcommand{\ImportTok}[1]{\textcolor[rgb]{0.00,0.50,0.00}{\textbf{#1}}}
\newcommand{\InformationTok}[1]{\textcolor[rgb]{0.38,0.63,0.69}{\textbf{\textit{#1}}}}
\newcommand{\KeywordTok}[1]{\textcolor[rgb]{0.00,0.44,0.13}{\textbf{#1}}}
\newcommand{\NormalTok}[1]{#1}
\newcommand{\OperatorTok}[1]{\textcolor[rgb]{0.40,0.40,0.40}{#1}}
\newcommand{\OtherTok}[1]{\textcolor[rgb]{0.00,0.44,0.13}{#1}}
\newcommand{\PreprocessorTok}[1]{\textcolor[rgb]{0.74,0.48,0.00}{#1}}
\newcommand{\RegionMarkerTok}[1]{#1}
\newcommand{\SpecialCharTok}[1]{\textcolor[rgb]{0.25,0.44,0.63}{#1}}
\newcommand{\SpecialStringTok}[1]{\textcolor[rgb]{0.73,0.40,0.53}{#1}}
\newcommand{\StringTok}[1]{\textcolor[rgb]{0.25,0.44,0.63}{#1}}
\newcommand{\VariableTok}[1]{\textcolor[rgb]{0.10,0.09,0.49}{#1}}
\newcommand{\VerbatimStringTok}[1]{\textcolor[rgb]{0.25,0.44,0.63}{#1}}
\newcommand{\WarningTok}[1]{\textcolor[rgb]{0.38,0.63,0.69}{\textbf{\textit{#1}}}}
\usepackage{longtable,booktabs,array}
\usepackage{calc} % for calculating minipage widths
% Correct order of tables after \paragraph or \subparagraph
\usepackage{etoolbox}
\makeatletter
\patchcmd\longtable{\par}{\if@noskipsec\mbox{}\fi\par}{}{}
\makeatother
% Allow footnotes in longtable head/foot
\IfFileExists{footnotehyper.sty}{\usepackage{footnotehyper}}{\usepackage{footnote}}
\makesavenoteenv{longtable}
\ifLuaTeX
\usepackage[bidi=basic]{babel}
\else
\usepackage[bidi=default]{babel}
\fi
\babelprovide[main,import]{english}
% get rid of language-specific shorthands (see #6817):
\let\LanguageShortHands\languageshorthands
\def\languageshorthands#1{}
\ifLuaTeX
  \usepackage[english]{selnolig} % disable illegal ligatures
\fi
\setlength{\emergencystretch}{3em} % prevent overfull lines
\providecommand{\tightlist}{%
  \setlength{\itemsep}{0pt}\setlength{\parskip}{0pt}}
% Raw LaTeX preamble for the ChakraDB book (included verbatim by pandoc via
% --include-in-header, so it is NOT markdown-processed).
\usepackage{graphicx}
\usepackage{float}
\usepackage{xcolor}
\usepackage{pagecolor}
\usepackage{tikz}
\usepackage{fancyhdr}
\usepackage{sectsty}
\usepackage[most]{tcolorbox}
\usepackage{epigraph}
\setlength{\epigraphwidth}{0.60\textwidth}
\setlength{\epigraphrule}{0pt}
\setlength{\beforeepigraphskip}{0.6\baselineskip}
\setlength{\afterepigraphskip}{2.0\baselineskip}
\renewcommand{\epigraphflush}{flushright}
\renewcommand{\epigraphsize}{\small\itshape}

% Theme: Harvard Crimson with saffron / yellow / white accents.
\definecolor{crimson}{HTML}{A51C30}    % Harvard crimson (cover + headings)
\definecolor{crimsondark}{HTML}{7B1420} % deeper crimson for subheadings
\definecolor{saffron}{HTML}{F4C430}    % saffron / marigold accent
\definecolor{saffrondeep}{HTML}{E8A317}% deeper saffron
\definecolor{goldyellow}{HTML}{FFD24A} % bright yellow highlight
\colorlet{chakra}{crimson}             % back-compat alias used elsewhere
\chapterfont{\color{crimson}}
\subsectionfont{\color{crimsondark}}
\sectionfont{\color{crimsondark}}

\pagestyle{fancy}
\fancyhf{}
\fancyhead[LE,RO]{\thepage}
\fancyhead[RE]{\itshape\nouppercase{\leftmark}}
\fancyhead[LO]{\itshape\nouppercase{\rightmark}}
\renewcommand{\headrulewidth}{0.4pt}

% A unicode-capable mono font so pseudocode glyphs (alpha, <=, ->, cap, ...)
% render in code listings. XeTeX only; the book is built with tectonic (XeTeX).
\usepackage{iftex}
\ifXeTeX
  \usepackage{fontspec}
  \setmonofont{DejaVu Sans Mono}[Scale=0.82]
\fi

% Render math symbols that appear in prose as proper LaTeX math (the body font
% lacks them). Code blocks are verbatim, so their glyphs come from the mono font
% (DejaVu Sans Mono) and are unaffected by these mappings.
\ifXeTeX
  \usepackage{newunicodechar}
  \newunicodechar{≤}{\ensuremath{\leq}}
  \newunicodechar{≥}{\ensuremath{\geq}}
  \newunicodechar{≈}{\ensuremath{\approx}}
  \newunicodechar{↔}{\ensuremath{\leftrightarrow}}
  \newunicodechar{⟺}{\ensuremath{\iff}}
  \newunicodechar{∎}{\ensuremath{\blacksquare}}
  \newunicodechar{✅}{\textcolor{crimson}{\checkmark}}
  \newunicodechar{❌}{\textcolor{crimsondark}{\ensuremath{\times}}}
\fi

% Keep long code lines from overflowing the right margin.
\usepackage{fvextra}
\DefineVerbatimEnvironment{Highlighting}{Verbatim}{breaklines,breakanywhere,commandchars=\\\{\}}
\usepackage{bookmark}
\IfFileExists{xurl.sty}{\usepackage{xurl}}{} % add URL line breaks if available
\urlstyle{same}
\hypersetup{
  pdftitle={ChakraDB --- The Definitive Guide},
  pdflang={en},
  colorlinks=true,
  linkcolor={RoyalBlue},
  filecolor={Maroon},
  citecolor={Blue},
  urlcolor={MidnightBlue},
  pdfcreator={LaTeX via pandoc}}

\author{}
\date{}

\begin{document}
\frontmatter

% Custom cover for the ChakraDB book (pandoc --include-before-body).
% Theme: solid Harvard Crimson background, saffron/yellow/white accents.
\begingroup
\thispagestyle{empty}
\pagecolor{crimson}
\color{white}
\centering
\vspace*{2.2cm}

% --- Chakra wheel emblem (TikZ), in saffron ---
\begin{tikzpicture}[line width=0.5pt]
  \foreach \r in {0.45, 1.15, 1.7} { \draw[saffron] (0,0) circle (\r); }
  \foreach \a in {0,30,...,330} {
    \draw[saffron] (\a:0.45) -- (\a:1.7);
    \fill[goldyellow] (\a:1.7) circle (0.055);
  }
  \fill[goldyellow] (0,0) circle (0.09);
  % two small satellite wheels
  \begin{scope}[shift={(-2.7,0.6)}, scale=0.42]
    \draw[saffron] (0,0) circle (1.0);
    \foreach \a in {0,45,...,315} { \draw[saffron] (0,0) -- (\a:1.0); \fill[goldyellow] (\a:1.0) circle (0.09); }
    \fill[goldyellow] (0,0) circle (0.14);
  \end{scope}
  \begin{scope}[shift={(2.7,-0.6)}, scale=0.42]
    \draw[saffron] (0,0) circle (1.0);
    \foreach \a in {0,45,...,315} { \draw[saffron] (0,0) -- (\a:1.0); \fill[goldyellow] (\a:1.0) circle (0.09); }
    \fill[goldyellow] (0,0) circle (0.14);
  \end{scope}
  \draw[saffron, dashed, line width=0.4pt] (-1.85,0.35) -- (-1.7,0.15);
  \draw[saffron, dashed, line width=0.4pt] (1.85,-0.35) -- (1.7,-0.15);
\end{tikzpicture}

\vspace{1.4cm}
{\color{saffron}\scshape\large An Embedded-Database Monograph}\\[0.9cm]
{\color{white}\fontsize{46}{50}\selectfont\bfseries ChakraDB}\\[0.35cm]
{\color{saffron}\Large\itshape The Definitive Guide}\\[1.0cm]
{\color{white}\large\itshape The embedded {\color{goldyellow}\bfseries HTAP} database with built-in graph}\\[0.15cm]
{\color{white}\large\itshape capabilities — write, analyze, and traverse live data.}\\[2.0cm]

{\color{saffron}\rule{5cm}{0.8pt}}\\[0.8cm]
{\color{white}\large The ChakraDB Authors}\\[0.5cm]
{\color{saffron}\scshape Architecture $\cdot$ Algorithms $\cdot$ Graph $\cdot$ Comparisons}\quad{\color{goldyellow}$\bullet$}\quad{\color{saffron}\scshape Working Draft, 2026}

\vfill
\endgroup
\clearpage
\nopagecolor
\color{black}

{
\hypersetup{linkcolor=black}
\setcounter{tocdepth}{2}
\tableofcontents
}
\mainmatter
\part{Preface}

\chapter{About This Book}\label{about-this-book}

\epigraph{If I have seen further it is by standing on the shoulders of Giants.}{--- Isaac Newton}

This is the complete guide to \textbf{ChakraDB} --- an embedded HTAP
database with built-in graph capabilities. It is written for three
audiences at once, and it is organized so each can find its level:

\begin{itemize}
\tightlist
\item
  \textbf{Users} who want to build applications --- start at Part V
  (\emph{Getting Started}) and Part VI (\emph{Developer Guide \&
  Tutorials}).
\item
  \textbf{Architects} evaluating the engine --- read Part I
  (\emph{Introduction}), Part II (\emph{Architecture}), and Part IX
  (\emph{Comparative Studies}).
\item
  \textbf{Contributors and the curious} who want to know \emph{how} it
  works down to the algorithm --- Part III (\emph{Algorithms}) is the
  deep end.
\end{itemize}

\section{How the book is organized}\label{how-the-book-is-organized}

\begin{longtable}[]{@{}
  >{\raggedright\arraybackslash}p{(\linewidth - 2\tabcolsep) * \real{0.5000}}
  >{\raggedright\arraybackslash}p{(\linewidth - 2\tabcolsep) * \real{0.5000}}@{}}
\toprule\noalign{}
\begin{minipage}[b]{\linewidth}\raggedright
Part
\end{minipage} & \begin{minipage}[b]{\linewidth}\raggedright
What it covers
\end{minipage} \\
\midrule\noalign{}
\endhead
\bottomrule\noalign{}
\endlastfoot
I --- Introduction & What ChakraDB is, the problem it solves, its design
principles and cost model. \\
II --- Architecture & The engine's structure: storage, MVCC, durability,
compaction, the query router. \\
III --- Algorithms & Each core algorithm in detail, with pseudocode and
complexity. \\
IV --- Graph & The property-graph layer: adjacency, CSR snapshots, and
graph algorithms. \\
V --- Getting Started & Install, first database, Python, first graph. \\
VI --- Developer Guide & SQL, types, constraints, transactions, COPY,
backup, observability, tutorials. \\
VII --- Operations & Durability tuning, limits, monitoring, disaster
recovery. \\
VIII --- Case Studies & Worked end-to-end applications. \\
IX --- Comparative Studies & Head-to-head with DuckDB, SQLite, Neo4j,
PostgreSQL, with methodology. \\
X --- Reference & The Rust and Python APIs, configuration, errors, file
formats. \\
Appendices & Glossary, design-decision log, bibliography. \\
\end{longtable}

\section{Conventions}\label{conventions}

\begin{itemize}
\tightlist
\item
  \textbf{Diagrams} are rendered with
  \href{https://mermaid.js.org/}{Mermaid}. Every architecture and
  algorithm chapter has at least one.
\item
  \textbf{Code} is real. Rust snippets compile against the public API;
  SQL runs on the shipped engine.
\item
  \textbf{Callouts} flag the important asymmetries: \textgreater{} This
  is a \emph{design commitment} --- a place where ChakraDB deliberately
  pays a \textgreater{} cost so that it can be cheap somewhere that
  matters more.
\end{itemize}

\section{Relationship to the source
docs}\label{relationship-to-the-source-docs}

This book synthesizes and expands the design documents that live
alongside the code --- \texttt{requirements.md} (the architecture
specification), \texttt{sql-reference.md},
\texttt{clickbench-findings.md}, and the milestone records. Where the
book and a source doc disagree, the code is the arbiter; the book aims
to track it.

\section{Status}\label{status}

ChakraDB is a working engine --- durable, crash-tested, benchmarked
against DuckDB --- with an actively evolving surface. Chapters that
describe not-yet-shipped work (for example, parts of the graph roadmap)
say so explicitly.

\part{Introduction}

\chapter{Introduction to ChakraDB}\label{introduction-to-chakradb}

\epigraph{No man ever steps in the same river twice, for it is not the same river and he is not the same man.}{--- Heraclitus}

Most databases are built for data that has stopped moving. You load it,
then you query it. ChakraDB is built for the opposite premise --- the
one Heraclitus noticed twenty-five centuries ago: \textbf{the data is
always changing, and the interesting question is what is true of it
\emph{now}, while it is still arriving.} Everything in this book follows
from taking that premise seriously.

\section{The one-sentence summary}\label{the-one-sentence-summary}

\textbf{ChakraDB is an embedded database that serves a continuous stream
of durable writes while running analytical queries and graph traversals
that never block --- over the same live data, in one process, in an open
columnar format.}

Unpack that sentence and you have the whole system:

\begin{itemize}
\tightlist
\item
  \emph{embedded} --- a library linked into your process, like SQLite or
  DuckDB, not a server you operate.
\item
  \emph{continuous durable writes} --- ACID, WAL-backed, and
  \textbf{concurrent}: many writers at once.
\item
  \emph{analytical queries and graph traversals that never block} ---
  readers see a consistent snapshot and are never stopped by a writer.
\item
  \emph{the same live data} --- no ETL, no second system, no staleness
  between the operational store and the analytical one.
\item
  \emph{open columnar format} --- data lives on disk as Apache Arrow,
  readable by any Arrow tool.
\end{itemize}

\section{What makes it different}\label{what-makes-it-different}

Every database chooses a shape of work to be excellent at, and pays for
that choice somewhere else. ChakraDB's choice is the \emph{seam} between
transactional and analytical work, extended to graph --- the place most
systems refuse to stand.

Consider where the competition sits:

\begin{itemize}
\tightlist
\item
  \textbf{DuckDB} gives blazing analytics over data you loaded earlier.
  But it holds a \textbf{single-writer file lock} --- a second writer is
  refused at the OS level (\texttt{IO\ Error:\ Could\ not\ set\ lock}).
  It is superb at static analytics and structurally unable to serve a
  live write stream.
\item
  \textbf{SQLite} is a rock-solid embedded transactional store, but its
  analytics are row-at-a-time, and one writer serializes the entire
  database.
\item
  \textbf{Neo4j} gives graph traversals, but it is a server, and
  mutation takes locks that contend with the reads doing the traversal.
\item
  \textbf{NetworkX / igraph} give graph algorithms --- over a \emph{dead
  static copy} you loaded into memory, stale the instant it is built.
\end{itemize}

ChakraDB's wager is that the valuable modern workload is precisely the
one that falls \emph{between} these: \textbf{data that is still
arriving, that you want to analyze and traverse at the same time.} That
is HTAP --- Hybrid Transactional/Analytical Processing --- plus graph,
in an embedded library.

\begin{center}\includegraphics[width=0.92\linewidth,height=0.80\textheight,keepaspectratio]{diagrams/diagram-1.pdf}\end{center}

The differentiator is \textbf{concurrency, not raw scan speed.} ChakraDB
does not aim to out-scan DuckDB on a static file; it aims to serve a
workload DuckDB \emph{structurally cannot} --- concurrent writers with
non-blocking analytical and graph reads.

\subsection{The four properties, at
once}\label{the-four-properties-at-once}

Individually, existing engines have three of the four properties below.
The gap ChakraDB targets is having all four together:

\begin{longtable}[]{@{}
  >{\raggedright\arraybackslash}p{(\linewidth - 8\tabcolsep) * \real{0.1304}}
  >{\centering\arraybackslash}p{(\linewidth - 8\tabcolsep) * \real{0.2174}}
  >{\centering\arraybackslash}p{(\linewidth - 8\tabcolsep) * \real{0.2174}}
  >{\centering\arraybackslash}p{(\linewidth - 8\tabcolsep) * \real{0.2174}}
  >{\centering\arraybackslash}p{(\linewidth - 8\tabcolsep) * \real{0.2174}}@{}}
\toprule\noalign{}
\begin{minipage}[b]{\linewidth}\raggedright
\end{minipage} & \begin{minipage}[b]{\linewidth}\centering
Embedded
\end{minipage} & \begin{minipage}[b]{\linewidth}\centering
ACID + MVCC
\end{minipage} & \begin{minipage}[b]{\linewidth}\centering
Concurrent writes + non-blocking scans
\end{minipage} & \begin{minipage}[b]{\linewidth}\centering
Open on-disk format
\end{minipage} \\
\midrule\noalign{}
\endhead
\bottomrule\noalign{}
\endlastfoot
DuckDB & ✅ & ✅ & ❌ single writer & ⚠ via extensions \\
SQLite & ✅ & ✅ & ❌ one writer & ❌ \\
Neo4j & ❌ server & ✅ & ⚠ lock contention & ❌ \\
ArcticDB & ✅ & ❌ & ✅ & ✅ \\
\textbf{ChakraDB} & ✅ & ✅ & ✅ & ✅ Arrow IPC parts \\
\end{longtable}

\subsection{\ldots and now, graph}\label{and-now-graph}

On top of that HTAP core, ChakraDB adds a \textbf{built-in graph layer}
--- not a bolted-on second store, but a thin layer that reuses exactly
the properties above. A graph edge whose key encodes
\texttt{(src,\ dst)} is stored \emph{sorted by source}, so a node's
neighbors are a contiguous key range (clustered adjacency for free); a
graph algorithm builds its working set from \textbf{one MVCC snapshot},
so it runs to completion over a consistent graph while edges keep
streaming in (live graph analytics); and sorted-by-source edges are
already grouped by source, so building the CSR representation every
algorithm wants is a single linear scan (the storage format is the
algorithm's input format). The result is a client experience where BFS,
shortest path, PageRank, connected components, and triangle counting are
one method call away, over data you are still writing to.

\section{The problem: why another
database?}\label{the-problem-why-another-database}

The case for ChakraDB is a case against \textbf{fragmentation}. The
modern data stack has splintered a single logical need --- ``know what
is true of my live data'' --- across three or four systems: an OLTP
store for the writes, a warehouse for the analytics, a graph database
for the relationships, and an ETL fabric to keep them approximately in
sync. Each hop adds latency, staleness, operational surface, and a
consistency model to reconcile.

The fragmentation exists because of a genuine technical tension --- the
shapes of work want opposite physical layouts:

\begin{itemize}
\tightlist
\item
  \textbf{Transactional} work wants row-at-a-time mutation, point
  lookups, and a durable log. It is latency-bound and write-heavy.
\item
  \textbf{Analytical} work wants columnar, compressed, vectorized scans
  over large row counts. It is throughput-bound and read-heavy.
\item
  \textbf{Graph} work wants pointer-chasing adjacency and whole-graph
  iteration.
\end{itemize}

For decades the accepted wisdom --- Stonebraker's ``one size does not
fit all'' --- was that these needs demand separate engines. ChakraDB
does not dispute that a single \emph{physical layout} cannot serve all
three. It disputes that they need separate \emph{systems}. The insight
is that a \textbf{log-structured, Arrow-native store with MVCC} can
present the right layout to each: a row buffer and a sorted key index
for the transactional path, immutable columnar parts for the analytical
path, and sorted-by-source edges for the graph path --- all over
\textbf{one snapshot clock}, so a reader of any kind sees one consistent
instant while writers proceed unblocked.

The specific thing that becomes possible, that no single embedded engine
offered before, is: \textbf{ingest continuously, and analyze and
traverse the result at the same time, with no ETL and no lock fight.}
The \href{../case-studies/fraud.md}{case studies} show why that is not a
luxury --- for real-time fraud scoring, live recommendations, and
streaming dashboards, the staleness and the lock contention of the
fragmented stack are the whole problem.

\section{Design principles}\label{design-principles}

Five commitments shape every decision in this book. They are worth
stating up front because when a trade-off arises, these are how it is
resolved.

\begin{enumerate}
\def\labelenumi{\arabic{enumi}.}
\item
  \textbf{Concurrency is the wedge.} ChakraDB competes on serving
  writes-plus-analytics-plus-graph on live data --- a workload DuckDB
  structurally cannot --- \emph{not} on out-scanning it on a static
  file. When a choice trades a little cold-scan speed for a lot of
  live-concurrency, it takes that trade.
\item
  \textbf{Buy execution, build storage.} The vectorized analytical
  engine is a solved, multi-year problem; ChakraDB \emph{buys} it
  (Apache DataFusion) and spends its own effort where the
  differentiation is --- storage, MVCC, durability, the transactional
  interpreter, and the graph layer. (See
  \href{../engine/query-router.md}{The Query Layer}.)
\item
  \textbf{Cold data pays nothing.} The cost of concurrency is borne only
  by recently-modified data. A part that predates every live snapshot is
  scanned with no per-row version check at all --- so a billion-row
  table with a thousand recent mutations pays version costs on a
  thousand rows, not a billion. (See \href{../engine/mvcc.md}{MVCC}.)
\item
  \textbf{Immutability over invalidation.} Sealed parts are never edited
  in place; updates and deletes are new versions and deletion-vector
  marks. Immutable data is safe to share lock-free with any number of
  readers, is trivially cacheable, and makes an open on-disk format
  possible.
\item
  \textbf{Publish the harness, not just the number.} No neutral
  benchmark measures the concurrency axis ChakraDB competes on, so every
  performance claim in this book ships with the code that produced it
  (Part IX). A number without a reproducible source is a hypothesis, not
  a measurement.
\end{enumerate}

\section{The cost model}\label{the-cost-model}

Fast writes, fast scans, and high concurrency are in genuine tension;
every system claiming all three has chosen \emph{where to absorb the
cost}. Stating ChakraDB's choice explicitly:

\begin{longtable}[]{@{}
  >{\raggedright\arraybackslash}p{(\linewidth - 4\tabcolsep) * \real{0.3333}}
  >{\raggedright\arraybackslash}p{(\linewidth - 4\tabcolsep) * \real{0.3333}}
  >{\raggedright\arraybackslash}p{(\linewidth - 4\tabcolsep) * \real{0.3333}}@{}}
\toprule\noalign{}
\begin{minipage}[b]{\linewidth}\raggedright
The goal
\end{minipage} & \begin{minipage}[b]{\linewidth}\raggedright
The cost it creates
\end{minipage} & \begin{minipage}[b]{\linewidth}\raggedright
Where ChakraDB pays it
\end{minipage} \\
\midrule\noalign{}
\endhead
\bottomrule\noalign{}
\endlastfoot
Fast writes & data accumulates as unmerged parts → scans slow &
\textbf{compaction} (a designed, back-pressured subsystem), never the
read path \\
Fast scans & needs merged, sorted, columnar data → write amplification &
absorbed by compaction, off the write path \\
High concurrency & needs versioning → readers may pay version cost per
row & paid \textbf{only} by recently-modified data; cold parts pay
\texttt{O(1)} \\
Open on-disk format & needs a committed format + manifest → commit
overhead & Arrow IPC parts + a manifest; publication cadence is a
tunable knob \\
\end{longtable}

Three absorption points define the engine:

\begin{itemize}
\tightlist
\item
  \textbf{The cost of fast writes is paid by background-free,
  caller-driven compaction} --- and if compaction cannot keep up, the
  engine applies \emph{explicit} ingest backpressure rather than
  silently degrading scans. This is a hard commitment, not an
  aspiration.
\item
  \textbf{The cost of concurrency is paid only by recently-modified
  data} --- the MVCC fast path (principle 3).
\item
  \textbf{The cost of the open format is paid in visibility latency, not
  write latency} --- internal transactions commit fast to a local log;
  the durable columnar state is materialized on a separate cadence.
\end{itemize}

If any of these trades is unacceptable for your workload, the
architecture is the wrong fit --- so they are stated here, at the front,
to be challenged.

\section{Who uses ChakraDB}\label{who-uses-chakradb}

The engine fits workloads that are simultaneously \emph{operational} and
\emph{analytical}, on a single node, embedded in an application:

\begin{itemize}
\tightlist
\item
  \textbf{Real-time fraud and risk} --- score transactions as they
  arrive, and traverse the entity graph they form, over one consistent
  view.
\item
  \textbf{Live recommendations} --- personalized PageRank /
  common-neighbors over a graph that updates continuously.
\item
  \textbf{Streaming dashboards} --- ingest events and serve aggregations
  without an ETL hop or a read/write lock fight.
\item
  \textbf{Exact-money ledgers} --- \texttt{DECIMAL} arithmetic that
  never rounds, with ACID transactions.
\end{itemize}

\section{The technology stack}\label{the-technology-stack}

ChakraDB is written in Rust; the core forbids \texttt{unsafe} code. Data
is Apache \textbf{Arrow} end to end --- in memory, on disk (Arrow IPC),
and across the DataFusion boundary (zero-copy). Vectorized analytical
execution is \textbf{bought} from Apache DataFusion; storage, MVCC,
durability, the transactional interpreter, and the graph layer are
\textbf{built}. The next part,
\href{../engine/overview.md}{Architecture}, shows how those pieces fit,
and why the ``buy execution, build storage'' split is the whole
strategy.

\part{The Engine}

\chapter{System Overview}\label{system-overview}

\epigraph{Architecture is the decisions that you wish you could get right early in a project.}{--- Ralph Johnson}

ChakraDB is a stack of narrow layers with one guiding split: \textbf{buy
execution, build storage.} The vectorized analytical engine is Apache
DataFusion; the storage engine, MVCC, durability, transactional
interpreter, and graph layer are ChakraDB's own. The layers meet at
deliberately thin interfaces.

\begin{center}\includegraphics[width=0.92\linewidth,height=0.80\textheight,keepaspectratio]{diagrams/diagram-2.pdf}\end{center}

\section{The layers, top to bottom}\label{the-layers-top-to-bottom}

\textbf{API.} Three entry points: the Rust
\texttt{Database}/\texttt{Storage}/\texttt{Graph} types, a
standards-based Python DB-API 2.0 (PEP 249) driver, and the SQL front
end (\texttt{SqlEngine}).

\textbf{Execution --- the HTAP router.} Every SQL statement is routed to
whichever engine fits its shape. Cheap transactional shapes --- writes,
key point-lookups, metadata \texttt{COUNT(*)}/\texttt{MIN}/\texttt{MAX},
and selective range scans --- run on a hand-written \textbf{interpreter}
that owns the write path. Analytical shapes --- large scans, joins,
windows, subqueries --- run on \textbf{DataFusion}. The \textbf{graph
engine} is a third consumer, building CSR working sets from the same
parts. See \href{query-router.md}{The HTAP Router}.

\textbf{Storage Engine.} A three-tier log-structured store: an in-memory
\textbf{L0} buffer absorbs writes at memory speed; when it fills it is
sealed, sorted by key, into an immutable columnar \textbf{Arrow part};
parts accumulate and are periodically merged by compaction. Every part
carries a small resident \textbf{index} --- a Bloom filter, per -column
min/max zonemaps, version stamps, and a deletion vector. See
\href{storage.md}{The Storage Engine}.

\textbf{Durability.} A write-ahead log with group commit makes writes
durable before they are acknowledged; a manifest records the catalog and
the set of live parts; recovery replays the WAL past the last
checkpoint. Everything reaches disk through a single \texttt{trait\ Io},
which has a real POSIX implementation and an in-memory fault -injecting
one for tests. See \href{durability.md}{Durability}.

\section{The two invariants that make it
fast}\label{the-two-invariants-that-make-it-fast}

Two design invariants recur throughout the architecture, and they are
worth holding in mind before the details:

\begin{quote}
\textbf{1. The sorted key column \emph{is} the index.} Because a sealed
part is written sorted by its key, a row's ordinal position \emph{is}
its offset. There is no separate key→location map to store or maintain
--- the index cost is \textasciitilde1.25 bytes per row and flat with
table size. This is the \href{storage.md}{primary-key index}.
\end{quote}

\begin{quote}
\textbf{2. Cold data pays zero concurrency cost.} MVCC visibility is a
pure function of a snapshot number and a part's version stamps. A part
that predates every live snapshot is ``fully visible'' and is scanned
with no per-row version check at all. A billion-row table with a
thousand recent mutations pays version costs on a thousand rows, not a
billion. This is \href{mvcc.md}{MVCC}.
\end{quote}

\section{The data flow of a write and a
read}\label{the-data-flow-of-a-write-and-a-read}

A \textbf{write} appends to the WAL (durable), updates the L0 buffer and
the MVCC clock, and acknowledges --- all without touching any reader. A
\textbf{read} takes a snapshot number, then either funnels through the
index to a single row (point lookup) or scans the parts visible to that
snapshot (analytical / graph), holding no lock that a writer could
contend on. The following chapters trace both paths in detail.

\chapter{The Storage Engine}\label{the-storage-engine}

\epigraph{Show me your flowcharts and conceal your tables, and I shall continue to be mystified. Show me your tables, and I won't usually need your flowcharts; they'll be obvious.}{--- Fred Brooks}

ChakraDB's storage is \textbf{log-structured} and \textbf{Arrow-native}.
Writes land in memory and are sealed into immutable columnar parts;
reads see a merged view across the tiers. This chapter is the shape of
the store; the algorithms that run over it --- visibility, merge,
pruning --- get their own chapters in Part III.

\section{Three tiers}\label{three-tiers}

\begin{center}\includegraphics[width=0.92\linewidth,height=0.80\textheight,keepaspectratio]{diagrams/diagram-3.pdf}\end{center}

\begin{enumerate}
\def\labelenumi{\arabic{enumi}.}
\item
  \textbf{L0 --- the write buffer.} New row versions land here at memory
  speed. L0 is a small in-memory structure keyed for point lookup and
  range scan. It holds the newest, not-yet-sealed writes.
\item
  \textbf{Sealed parts --- the columnar body.} When L0 fills (a
  configurable threshold), it is \textbf{sealed}: its rows are sorted by
  the table's key and written as an immutable \textbf{Apache Arrow}
  record batch. On the durable path a part persists as an Arrow
  \textbf{IPC} stream --- an open format any Arrow reader can open. A
  part never changes after sealing; updates and deletes are expressed as
  \emph{new} versions and \emph{deletion-vector} marks, not in-place
  edits.
\item
  \textbf{Compaction --- keeping scans fast.} Parts accumulate as writes
  flow. A background-free, caller-driven \textbf{compaction} merges
  parts, drops rows no live snapshot can see, and collapses version
  stamps --- trading write amplification for scan speed. See
  \href{storage.md}{Compaction}.
\end{enumerate}

\begin{quote}
\textbf{The absorption point.} Fast writes create unmerged deltas that
slow scans; ChakraDB pays that debt in \textbf{compaction}, not in the
read path, and applies explicit \textbf{backpressure} if compaction
cannot keep up --- never silent scan degradation.
\end{quote}

\section{What a part carries: the resident
index}\label{what-a-part-carries-the-resident-index}

The reason ChakraDB can keep row \emph{data} on disk but stay fast is
that each part keeps a small \textbf{index resident in memory} --- about
\textbf{1.25 bytes per row}, flat with table size:

\begin{center}\includegraphics[width=0.92\linewidth,height=0.80\textheight,keepaspectratio]{diagrams/diagram-4.pdf}\end{center}

\begin{itemize}
\tightlist
\item
  \textbf{Bloom filter} --- answers ``could key \emph{k} be in this
  part?'' with no disk touch, so a point lookup skips parts that
  certainly lack the key.
\item
  \textbf{Zonemaps} --- per-column \texttt{(min,\ max)}. A
  \texttt{WHERE} range or a graph adjacency scan skips any part whose
  range cannot overlap. This is \href{storage.md}{zonemap pruning}.
\item
  \textbf{Version stamps} --- the \texttt{created}/\texttt{deleted} CSN
  per row (or one uniform stamp when the whole part shares one), the
  input to \href{mvcc.md}{MVCC visibility}.
\item
  \textbf{Deletion vector} --- which ordinals are tombstoned, so a scan
  skips deleted rows without rewriting the part.
\item
  \textbf{The sorted key run} \emph{is} the index: because the part is
  sorted by key, the ordinal position is the row offset, and a lookup is
  a Bloom probe plus a binary search --- no separate key→location map
  exists. See \href{storage.md}{The Primary-Key Index}.
\end{itemize}

\section{Arbitrary schemas, any-type
keys}\label{arbitrary-schemas-any-type-keys}

A table has any number of columns of any supported type. Its
\textbf{primary key} is a single column of any type (integer, text,
float, boolean, date, decimal), or --- if no key is declared --- a
hidden auto-increment \texttt{\_rowid} (a \emph{keyless} table). One
idea keeps the engine simple: \textbf{every table has exactly one key
column}; ``keyless'' is just a table whose key is hidden.

\section{Why not a general buffer
pool?}\label{why-not-a-general-buffer-pool}

Parts are immutable and whole-part granular, so ChakraDB deliberately
has \textbf{no LRU, no page-replacement cache}. The policy is ``fault a
part's data in on first touch, keep it until the part is dropped.'' That
is the honest amount of machinery for immutable parts --- and it is why
the \emph{resident index}, not disk, is the scaling ceiling (see
\href{overview.md}{Limits}).

\section{The write path in one
picture}\label{the-write-path-in-one-picture}

\begin{center}\includegraphics[width=0.92\linewidth,height=0.80\textheight,keepaspectratio]{diagrams/diagram-5.pdf}\end{center}

Notice what the write path does \emph{not} do: it never takes a lock a
reader holds. Readers run against a snapshot number and see a consistent
view; writers advance the clock and append. That decoupling is the
concurrency wedge, and the next chapter --- MVCC --- is how it stays
correct.

\section{The Sorted-Part Key Index}\label{the-sorted-part-key-index}

This is the hardest and most consequential piece of the storage engine:
how a row is located by key without paying for a key→location map. The
answer --- following Apache Doris / StarRocks --- is that \textbf{the
sorted key column \emph{is} the index.}

\section{The idea}\label{the-idea}

A sealed part is written \textbf{sorted by its primary key}. Therefore a
row's \emph{ordinal position} in the part is its \emph{offset}: to find
where key \texttt{k} lives, binary-search the sorted key column. There
is no separate structure mapping keys to row numbers, so the resident
index cost is not the \textasciitilde12 bytes/row an explicit map would
need --- it is the Bloom filter plus min/max bounds, about \textbf{1.25
bytes/row, flat with table size.}

\begin{Shaded}
\begin{Highlighting}[]
\NormalTok{flowchart LR}
\NormalTok{    K["lookup key = k"] {-}{-}\textgreater{} B\{"min ≤ k ≤ max?\textless{}br/\textgreater{}(part bounds)"\}}
\NormalTok{    B {-}{-}\textgreater{}|no| SKIP["skip part"]:::x}
\NormalTok{    B {-}{-}\textgreater{}|yes| BL\{"Bloom: might contain k?"\}}
\NormalTok{    BL {-}{-}\textgreater{}|no| SKIP}
\NormalTok{    BL {-}{-}\textgreater{}|yes| BS["binary{-}search the\textless{}br/\textgreater{}sorted key column"]:::ok}
\NormalTok{    BS {-}{-}\textgreater{} HIT["ordinal → row"]}
\NormalTok{    classDef x fill:\#f5d6d6; classDef ok fill:\#d6f5d6;}
\end{Highlighting}
\end{Shaded}

\section{The lookup funnel}\label{the-lookup-funnel}

A point lookup consults a cascade of cheap filters before it ever
binary-searches, newest tier first (a later write must win):

\begin{quote}
\textbf{ALGORITHM 1 --- Point lookup by key}

\begin{Shaded}
\begin{Highlighting}[]
\NormalTok{Input:  key k; snapshot S}
\NormalTok{Output: the visible row for k, or NONE}
\NormalTok{1  if L0 has a version of k visible to S:            ▷ newest writes first}
\NormalTok{2      return that version}
\NormalTok{3  for each sealed part P, newest to oldest:}
\NormalTok{4      if k \textless{} P.min\_key or k \textgreater{} P.max\_key: continue    ▷ ALG 11: bounds skip}
\NormalTok{5      if not P.bloom.might\_contain(k):    continue    ▷ ALG 15: Bloom skip}
\NormalTok{6      o ← BinarySearchKey(P, k)                       ▷ ALGORITHM 2}
\NormalTok{7      if o found and version at o is visible to S:}
\NormalTok{8          return the row at ordinal o}
\NormalTok{9  return NONE}
\end{Highlighting}
\end{Shaded}
\end{quote}

Each part that certainly lacks \texttt{k} is dropped by a \textbf{bounds
comparison} (min/max) or a \textbf{Bloom probe} --- neither touches the
column data on disk. Only a part that \emph{might} hold \texttt{k} is
searched, and the search is over the small resident key run.

\section{Binary search within a part}\label{binary-search-within-a-part}

\begin{quote}
\textbf{ALGORITHM 2 --- Binary search the sorted key column}

\begin{Shaded}
\begin{Highlighting}[]
\NormalTok{Input:  part P (sorted by key), key k}
\NormalTok{Output: an ordinal o with P.key(o) = k, or NOTFOUND}
\NormalTok{1  lo ← 0;  hi ← P.len − 1}
\NormalTok{2  while lo ≤ hi:}
\NormalTok{3      mid ← (lo + hi) / 2}
\NormalTok{4      c ← total\_cmp(P.key(mid), k)                    ▷ total order over Values}
\NormalTok{5      if c \textless{} 0: lo ← mid + 1}
\NormalTok{6      elif c \textgreater{} 0: hi ← mid − 1}
\NormalTok{7      else: return mid                                ▷ found}
\NormalTok{8  return NOTFOUND}
\end{Highlighting}
\end{Shaded}
\end{quote}

Keys are compared with \texttt{total\_cmp}, a total order over all value
types (integers compared exactly --- never routed through \texttt{f64},
which would conflate integers beyond 2⁵³ and corrupt an integer key).
For a duplicate-tolerant scan the search extends left/right from
\texttt{mid} to the run of equal keys.

\section{Any-type keys, and the hidden
rowid}\label{any-type-keys-and-the-hidden-rowid}

The key column may be any type --- integer, text, float, boolean, date,
decimal --- because \texttt{total\_cmp} orders them all. A table with
\textbf{no} declared key gets a hidden auto-increment \texttt{\_rowid};
it is still a single sorted key column, just invisible to
\texttt{SELECT\ *}. So ``keyless'' costs nothing extra: there is no
second index for the rowid --- the sorted rowid column \emph{is} the
index, exactly as for a user key.

\section{Why this shape}\label{why-this-shape}

\begin{quote}
\textbf{Proposition 1 (Index cost is flat).} The resident per-row index
cost of the sorted-part scheme is independent of the number of rows.

\emph{Proof sketch.} The resident structures are the Bloom filter (a
fixed bits-per-key budget), the per-part min/max bounds (constant per
part), the per-row version stamps, and the deletion vector. None grows
super-linearly with row count, and crucially there is \textbf{no}
key→location map --- the map is replaced by ``ordinal = offset,'' which
stores nothing. Measured at \textasciitilde1.25 B/row, it stays flat as
tables grow (\texttt{m0-bench}). ∎
\end{quote}

The consequence is the whole storage strategy: because the index is
nearly free and flat, ChakraDB can hold row \emph{data} on disk and keep
only the index resident, and the resident index --- not disk --- becomes
the scaling ceiling (see \href{overview.md}{Limits}).

\section{The Merge / Compaction
Algorithm}\label{the-merge-compaction-algorithm}

Writes create parts; parts accumulate; scans slow down. Compaction is
the debt payment: it merges parts into fewer, larger, sorted parts,
physically drops rows no live reader can see, and collapses redundant
version stamps. It is \textbf{caller-driven} --- there is no background
thread --- so it runs when asked, and only up to a safe
\href{mvcc.md}{GC horizon}.

\section{The k-way merge}\label{the-k-way-merge}

Because every input part is sorted by key, merging them into one sorted
part is a classic k-way merge over a heap of cursors:

\begin{quote}
\textbf{ALGORITHM 8 --- Merge parts}

\begin{Shaded}
\begin{Highlighting}[]
\NormalTok{Input:  parts P₁..P\_k (each sorted by key); reclamation horizon H}
\NormalTok{Output: one merged, sorted part}
\NormalTok{1  heap ← \{ (P\_i.key(0), i, 0) for each non{-}empty P\_i \}   ▷ min{-}heap on key}
\NormalTok{2  out ← empty builder}
\NormalTok{3  while heap not empty:}
\NormalTok{4      (k, i, o) ← heap.pop\_min()                          ▷ smallest current key}
\NormalTok{5      keep ← RowIsLive(P\_i, o, H)                          ▷ see ALGORITHM 9}
\NormalTok{6      if keep: out.append(P\_i.row(o))                      ▷ carry the surviving row}
\NormalTok{7      if o+1 \textless{} P\_i.len: heap.push((P\_i.key(o+1), i, o+1))  ▷ advance that cursor}
\NormalTok{8  return out.seal()                                        ▷ sorted Arrow part}
\end{Highlighting}
\end{Shaded}
\end{quote}

The output is produced in sorted key order in a single streaming pass
--- \texttt{O(N\ log\ k)} for \texttt{N} total rows and \texttt{k}
inputs --- with no intermediate sort. Newer parts take precedence when
keys tie, so an updated row's newest version wins.

\begin{center}\includegraphics[width=0.92\linewidth,height=0.80\textheight,keepaspectratio]{diagrams/diagram-7.pdf}\end{center}

\section{Reclamation: which rows
survive}\label{reclamation-which-rows-survive}

The horizon \texttt{H} is the oldest CSN any live snapshot may observe
(the \href{mvcc.md}{GC watermark}). A row can be physically dropped only
when no snapshot \texttt{≥\ H} could still see it:

\begin{quote}
\textbf{ALGORITHM 9 --- Row liveness under a horizon}

\begin{Shaded}
\begin{Highlighting}[]
\NormalTok{Input:  a row at ordinal o in part P; horizon H}
\NormalTok{Output: true if the row must be kept}
\NormalTok{1  if the row is not tombstoned: return true             ▷ a live version — keep}
\NormalTok{2  d ← its deletion CSN}
\NormalTok{3  if d ≤ H: return false                                 ▷ deleted before/at the}
\NormalTok{4  return true                                            ▷   horizon → reclaimable}
\end{Highlighting}
\end{Shaded}
\end{quote}

Line 3 is the crux: a row deleted at \texttt{d\ ≤\ H} is invisible to
\emph{every} live snapshot (all are
\texttt{≥\ H\ \textgreater{}}\ldots{} \texttt{≥\ d}), so it is safe to
drop. A row deleted \emph{after} \texttt{H} is still visible to a reader
on an old snapshot and \textbf{must be kept}. Getting this bound wrong
is the silent-wrong-answer hazard the \href{mvcc.md}{GC watermark}
chapter is devoted to.

After merging, the surviving versions' stamps are \textbf{collapsed}: if
all rows of the merged part now share one creation CSN, the part stores
a single \texttt{Uniform(csn)} instead of a per-row array --- restoring
the cheap \href{mvcc.md}{full-part fast path} for future scans.

\section{Selection and
incrementality}\label{selection-and-incrementality}

Two more practicalities keep compaction from being wasteful:

\begin{itemize}
\tightlist
\item
  \textbf{Selection.} Not every part is merged every time. A policy
  chooses a set of candidates (for example, similar-sized adjacent
  parts), so compaction does bounded work per invocation rather than
  rewriting the whole table.
\item
  \textbf{Incremental persistence.} When the merged result is
  checkpointed, unchanged parts are skipped and a part that only
  \emph{gained tombstones} since the last checkpoint has just those
  deletion-vector entries appended --- not a full rewrite. This fixed an
  \texttt{O(n²)} that a naive ``rewrite everything each checkpoint''
  would incur.
\end{itemize}

\section{The backpressure contract}\label{the-backpressure-contract}

Compaction trades write amplification for scan speed. If it cannot keep
up, the engine does \textbf{not} silently let scans degrade --- it
applies explicit ingest \textbf{backpressure} (a throttle, then a stall)
keyed on the part count, surfaced in
\href{../getting-started/getting-started.md}{\texttt{Storage::stats()}}.
This is the ``cost of fast writes is paid in compaction, and made
visible'' commitment from the \href{../introduction/cost-model.md}{cost
model}.

\begin{quote}
\textbf{Proposition 6 (Merge preserves the visible database).} For any
snapshot \texttt{S\ ≥\ H}, the set of rows visible to \texttt{S} is
identical before and after a merge with horizon \texttt{H}.

\emph{Proof sketch.} A merge changes physical layout, not logical
content, except that it drops rows with \texttt{deleted\ ≤\ H} (ALG 9).
Such a row is invisible to any \texttt{S\ ≥\ H} because
\texttt{S\ ≥\ H\ ≥\ deleted} fails \texttt{S\ \textless{}\ deleted}
(\href{mvcc.md}{ALGORITHM 3}). Every row visible to \texttt{S} has
\texttt{deleted\ \textgreater{}\ H}, is therefore kept, and its
\texttt{(created,\ deleted)} stamps are preserved (collapse only
rewrites the storage form, not the values). Hence \texttt{S}'s visible
set is unchanged. The requirement \texttt{S\ ≥\ H} is exactly what the
GC watermark guarantees for every live reader. ∎
\end{quote}

\section{Zonemap Part Pruning}\label{zonemap-part-pruning}

A selective query should not read parts that cannot contain a match.
ChakraDB skips them using \textbf{zonemaps} --- the per-column
\texttt{(min,\ max)} bounds every part carries. This is DuckDB-style
rowgroup pruning, and it is also the mechanism behind
\href{../graph/model.md}{graph adjacency}.

\section{The idea}\label{the-idea-1}

Each sealed part stores, per column, the minimum and maximum value over
its rows. For a predicate on that column, if the part's
\texttt{{[}min,\ max{]}} interval cannot satisfy the predicate, no row
in the part can --- so the whole part is skipped without touching its
data.

\begin{Shaded}
\begin{Highlighting}[]
\NormalTok{flowchart LR}
\NormalTok{    Q["WHERE x \textgreater{}= 500"]}
\NormalTok{    P1["Part A\textless{}br/\textgreater{}x ∈ [1, 120]"]:::skip}
\NormalTok{    P2["Part B\textless{}br/\textgreater{}x ∈ [300, 900]"]:::hit}
\NormalTok{    P3["Part C\textless{}br/\textgreater{}x ∈ [950, 999]"]:::hit}
\NormalTok{    Q {-}{-}\textgreater{} P1 \& P2 \& P3}
\NormalTok{    classDef skip fill:\#f5d6d6,stroke:\#c00; classDef hit fill:\#d6f5d6,stroke:\#0a0;}
\end{Highlighting}
\end{Shaded}

Part A (\texttt{max\ =\ 120\ \textless{}\ 500}) is skipped; B and C are
scanned.

\section{The exclusion test}\label{the-exclusion-test}

The predicate is compiled to a conservative \emph{excludes} test: given
a part's bounds, can this predicate be satisfied by \emph{no} row in the
part?

\begin{quote}
\textbf{ALGORITHM 11 --- Predicate excludes a part}

\begin{Shaded}
\begin{Highlighting}[]
\NormalTok{Input:  predicate φ; part bounds (per column) [min\_c, max\_c]}
\NormalTok{Output: true if NO row in the part can satisfy φ (safe to skip)}
\NormalTok{1  match φ:}
\NormalTok{2    (col c) = v :  return v \textless{} min\_c  or  v \textgreater{} max\_c        ▷ value out of range}
\NormalTok{3    (col c) \textless{} v :  return min\_c ≥ v                        ▷ all values ≥ v}
\NormalTok{4    (col c) ≤ v :  return min\_c \textgreater{} v}
\NormalTok{5    (col c) \textgreater{} v :  return max\_c ≤ v                        ▷ all values ≤ v}
\NormalTok{6    (col c) ≥ v :  return max\_c \textless{} v}
\NormalTok{7    φ₁ AND φ₂ :  return excludes(φ₁) or excludes(φ₂)       ▷ either kills it}
\NormalTok{8    φ₁ OR  φ₂ :  return excludes(φ₁) and excludes(φ₂)      ▷ both must kill it}
\NormalTok{9    otherwise :  return false                              ▷ unsure ⇒ keep (safe)}
\end{Highlighting}
\end{Shaded}
\end{quote}

The test is \textbf{conservative}: it returns \texttt{true} only when it
can \emph{prove} the part holds no match, and \texttt{false} (keep the
part) whenever it is unsure. So pruning never changes an answer --- it
only avoids work.

The scan then drops every fully-materialised part the predicate
excludes:

\begin{Shaded}
\begin{Highlighting}[]
\NormalTok{segments ← scan\_segments(snapshot)}
\NormalTok{segments.retain(seg → not (seg is a sealed part P and excludes(φ, P.bounds)))}
\end{Highlighting}
\end{Shaded}

\section{Two ways it is used}\label{two-ways-it-is-used}

\textbf{As a SQL accelerator.} A selective \texttt{WHERE\ x\ =\ k} or
\texttt{WHERE\ x\ BETWEEN\ a\ AND\ b} prunes to the parts whose bounds
overlap --- the interpreter's second-stage router sends such queries
here rather than to a full vectorized scan (see
\href{query-router.md}{Query Routing}).

\textbf{As a graph adjacency index.} A graph edge key encodes
\texttt{(src,\ dst)} src-major, so ``neighbors of X'' is a key-range
scan \texttt{key\ ∈\ {[}(X,0),\ (X+1,0))}. Zonemap pruning on the key
column touches only the parts holding \texttt{src\ =\ X}. This is
\href{../graph/model.md}{clustered adjacency} --- the graph traversal
primitive falls directly out of ALGORITHM 11.

\section{The scaling property}\label{the-scaling-property}

\begin{quote}
\textbf{Proposition 7 (O(1) selective range scans).} For a key-range
query over a table sorted by that key, the number of parts scanned
depends on the range's selectivity, not the table size.

\emph{Proof sketch.} Parts are sorted by key, so each part's key-bounds
interval is disjoint and ordered. A key range \texttt{{[}lo,\ hi)}
overlaps only the contiguous run of parts whose intervals intersect it
--- a count that grows with \texttt{hi\ −\ lo}, not with the number of
parts. Every other part is excluded by ALG 11 (lines 2--6). Measured: a
needle range scan stays \textasciitilde1 ms as the table grows 100×
(ClickBench Q13, Part IX). ∎
\end{quote}

\section{The correctness guarantee,
restated}\label{the-correctness-guarantee-restated}

Pruning removes only parts it \emph{proves} empty of matches, so the
produced answer is identical to a full scan. The unit tests exercise the
boundaries directly --- equality at the min/max edges, swapped operands
(\texttt{v\ =\ col}), \texttt{AND}/\texttt{OR} combinations, and missing
bounds --- and an end-to-end suite scans a multi-part table to confirm a
pruned range returns exactly the matching rows.

\section{Bloom Filters \& the Lookup
Funnel}\label{bloom-filters-the-lookup-funnel}

A point lookup should not read a part that does not contain the key. The
min/max bounds catch some parts; a \textbf{Bloom filter} catches the
rest --- it answers ``could key \texttt{k} be here?'' from a few bits in
memory, with no disk touch and no false negatives.

\section{What a Bloom filter
guarantees}\label{what-a-bloom-filter-guarantees}

A Bloom filter is a bit array with \texttt{h} hash functions. To
\emph{insert} \texttt{k}, set the \texttt{h} bits it hashes to; to
\emph{test} \texttt{k}, check those \texttt{h} bits.

\begin{itemize}
\tightlist
\item
  If any bit is \texttt{0}, \texttt{k} is \textbf{definitely absent} ---
  a true negative.
\item
  If all \texttt{h} bits are \texttt{1}, \texttt{k} is \textbf{probably
  present} --- possibly a false positive.
\end{itemize}

Crucially there are \textbf{no false negatives}: a key that was inserted
always tests positive. That is exactly the property a lookup filter
needs --- it may occasionally fail to skip a part, but it never wrongly
skips one that holds the key.

\begin{quote}
\textbf{ALGORITHM 15 --- Bloom membership test}

\begin{Shaded}
\begin{Highlighting}[]
\NormalTok{Input:  key k; part Bloom filter B with h hash functions}
\NormalTok{Output: DEFINITELY\_ABSENT, or MAYBE\_PRESENT}
\NormalTok{1  seed ← value\_seed(k)                              ▷ reduce any Value type to 64 bits}
\NormalTok{2  for i in 0..h:}
\NormalTok{3      bit ← hash(seed, i) mod |B|}
\NormalTok{4      if B[bit] = 0: return DEFINITELY\_ABSENT        ▷ a single 0 bit settles it}
\NormalTok{5  return MAYBE\_PRESENT}
\end{Highlighting}
\end{Shaded}
\end{quote}

\texttt{value\_seed} maps a key of any type to a 64-bit seed: an integer
maps to itself (so the integer path is unchanged and collision-free in
that dimension), and text, float, boolean, and decimal keys get a
deterministic reduction.

\section{Its place in the funnel}\label{its-place-in-the-funnel}

The Bloom test is the second filter in the
\href{storage.md}{point-lookup funnel}, after the min/max bounds and
before the binary search:

\begin{Shaded}
\begin{Highlighting}[]
\NormalTok{flowchart LR}
\NormalTok{    K["lookup k"] {-}{-}\textgreater{} BND\{"min ≤ k ≤ max?"\}}
\NormalTok{    BND {-}{-}\textgreater{}|no| S1["skip"]:::x}
\NormalTok{    BND {-}{-}\textgreater{}|yes| BLM\{"Bloom: maybe present?"\}}
\NormalTok{    BLM {-}{-}\textgreater{}|absent| S2["skip"]:::x}
\NormalTok{    BLM {-}{-}\textgreater{}|maybe| BS["binary search\textless{}br/\textgreater{}(reads the key run)"]:::ok}
\NormalTok{    classDef x fill:\#f5d6d6; classDef ok fill:\#d6f5d6;}
\end{Highlighting}
\end{Shaded}

The ordering is deliberate --- cheapest test first. The bounds check is
two comparisons; the Bloom probe is a handful of bit reads; only a part
that passes both pays for a binary search over its resident key run. A
part that certainly lacks the key is dropped before any of its data is
examined.

\section{Sizing and cost}\label{sizing-and-cost}

The filter trades a small, fixed number of bits per key for a low
false-positive rate. In ChakraDB it is part of the \textasciitilde1.25
B/row resident index budget (\href{storage.md}{Proposition 1}): the
Bloom bits, the min/max bounds, the version stamps, and the deletion
vector together stay flat with table size, which is what lets the engine
keep data on disk and the index in memory.

\begin{quote}
\textbf{Proposition 11 (No false negatives ⇒ lookups are correct).} The
Bloom skip in ALGORITHM 15 never causes a lookup to miss a key that is
present.

\emph{Proof sketch.} If key \texttt{k} is in the part, it was inserted,
so all \texttt{h} of its bits are \texttt{1}, so ALG 15 returns
\texttt{MAYBE\_PRESENT} and the funnel proceeds to the binary search
that finds it. The filter can only err the other way --- a
\texttt{MAYBE\_PRESENT} for an absent key --- which costs a wasted
binary search, not a wrong answer. This is asserted directly: over many
keys, the filter has zero false negatives (\texttt{bloom} tests). ∎
\end{quote}

\section{Why it matters for
concurrency}\label{why-it-matters-for-concurrency}

Point lookups are the transactional read shape, and they run on the
interpreter precisely because this funnel makes them \texttt{O(log\ n)}
with almost no I/O. That keeps the transactional path fast and lock-free
while analytical scans run on the same snapshot --- the HTAP split the
\href{query-router.md}{router} exploits.

\chapter{MVCC \& Snapshot Isolation}\label{mvcc-snapshot-isolation}

\epigraph{Time is what keeps everything from happening at once.}{--- Ray Cummings}

Multi-Version Concurrency Control is the reason readers never block
writers. This chapter is the \emph{model}; the \href{mvcc.md}{visibility
algorithm} and the \href{mvcc.md}{GC watermark} get the full treatment
in Part III.

\section{The commit sequence number
(CSN)}\label{the-commit-sequence-number-csn}

There is one piece of global state every writer touches: a monotonic
64-bit counter, the \textbf{commit sequence number}. Allocating a CSN is
a single atomic increment. A \textbf{snapshot} is just a CSN --- a
number naming an instant.

Every row version carries two stamps:

\begin{itemize}
\tightlist
\item
  \texttt{created} --- the CSN at which this version came to exist.
\item
  \texttt{deleted} --- the CSN at which it was superseded or removed (or
  \texttt{NEVER\_DELETED}).
\end{itemize}

\section{The visibility predicate}\label{the-visibility-predicate}

A version is visible to a snapshot \texttt{S} if and only if it was
created at or before \texttt{S} and not yet deleted as of \texttt{S}:

\begin{verbatim}
visible(S, created, deleted)  ⟺  created ≤ S < deleted
\end{verbatim}

That is the whole rule. It is a pure function of three integers --- no
locks, no undo log, no read view to construct.

\begin{center}\includegraphics[width=0.92\linewidth,height=0.80\textheight,keepaspectratio]{diagrams/diagram-10.pdf}\end{center}

At \texttt{S=3}, only \texttt{v1} satisfies
\texttt{created\ ≤\ 3\ \textless{}\ deleted} (1 ≤ 3 \textless{} 5). At
\texttt{S=7}, only \texttt{v2} (5 ≤ 7 \textless{} 9). Exactly one
version of a live key is visible to any snapshot --- which is what makes
an \texttt{UPDATE} a create-plus-delete rather than a mutation.

\section{Why cold data is free}\label{why-cold-data-is-free}

The predicate has two fast-path corollaries that make ChakraDB's scans
cheap:

\begin{quote}
\textbf{A part created entirely at or before \texttt{S} and with no
deletions after \texttt{S} is ``fully visible'' --- every one of its
rows passes, so the scan skips the per-row check entirely.}
\end{quote}

Concretely, a part stores a uniform \texttt{created} stamp and the
minimum \texttt{deleted} CSN across its rows. If
\texttt{created\_max\ ≤\ S} and \texttt{S\ \textless{}\ min\_deleted},
the whole part is visible with a single comparison. A cold, unmodified
part --- the vast majority of a large table --- pays \textbf{zero}
per-row version cost on a scan. Only recently-modified parts and the L0
buffer pay the per-row predicate. This is the ``cost of concurrency is
paid only by data that was recently modified'' principle, from Neumann
et al.'s MVCC design, made concrete.

\section{Snapshot isolation, across
tables}\label{snapshot-isolation-across-tables}

The CSN clock is global to a database, so a snapshot is consistent
\textbf{across every table at once.} An application reading two tables
at snapshot \texttt{S} observes both as of the same instant --- no torn
cross-table view. This is what lets the graph layer build a consistent
adjacency from an edges table while a nodes table is also being written.

\section{Readers do not block writers --- and the
reverse}\label{readers-do-not-block-writers-and-the-reverse}

\begin{itemize}
\tightlist
\item
  A \textbf{reader} takes a snapshot number and scans the parts visible
  to it. It grabs the part list under a brief read lock, then scans
  \textbf{lock-free}; the parts are immutable, so nothing a writer does
  can change what the reader is reading.
\item
  A \textbf{writer} advances the clock and appends a new version. It
  never waits on a reader.
\end{itemize}

The one serialization point is \emph{per table}: two writers to the
\emph{same} table take a short write lock for the L0 insert. Different
tables write concurrently. (This is the deliberate trade in the
\href{../introduction/cost-model.md}{cost model}.)

\section{Garbage collection needs a
watermark}\label{garbage-collection-needs-a-watermark}

Because old versions stay until compaction reclaims them, compaction
must not drop a version some live snapshot can still see. ChakraDB
tracks a \textbf{live-snapshot registry}: a long-running reader (a query
or a transaction) \emph{pins} its snapshot, and compaction's reclamation
horizon is held back to the oldest pinned snapshot. A reader on an old
snapshot keeps its view even as newer writes are compacted away. The
full mechanism --- and why it is correct under concurrency --- is in
\href{mvcc.md}{The GC Watermark}.

\section{Durability of the clock}\label{durability-of-the-clock}

The CSN is monotonic and never reset. On recovery, the counter is raised
above the highest CSN any replayed record used, so a recovered database
never reissues a stamp --- two rows can never collide on a version
number. The counter is 64-bit: \textasciitilde1.8×10¹⁹ total mutations
over a database's entire lifetime.

\section{MVCC Visibility}\label{mvcc-visibility}

Visibility is the pure function at the heart of snapshot isolation. It
decides, for a snapshot \texttt{S} and a row version stamped
\texttt{(created,\ deleted)}, whether that version is the one \texttt{S}
should see. Everything about ChakraDB's non-blocking reads follows from
how cheap this function is.

\section{The predicate}\label{the-predicate}

\begin{quote}
\textbf{ALGORITHM 3 --- Version visibility}

\begin{Shaded}
\begin{Highlighting}[]
\NormalTok{Input:  snapshot S (a CSN); version stamps created, deleted}
\NormalTok{Output: true if the version is visible to S}
\NormalTok{1  return created ≤ S  and  S \textless{} deleted            ▷ deleted = ∞ if never deleted}
\end{Highlighting}
\end{Shaded}
\end{quote}

That is the whole rule: a version is visible iff it was created at or
before \texttt{S} and not yet deleted as of \texttt{S}. It is three
integer comparisons --- no lock, no undo log, no read-view to
materialize.

An \texttt{UPDATE} is modeled as \emph{delete-old + create-new}: the old
version's \texttt{deleted} is set to the new CSN, and the new version's
\texttt{created} is that same CSN.

\begin{quote}
\textbf{Proposition 2 (Exactly one live version).} For any live key and
any snapshot \texttt{S}, exactly one version satisfies
\texttt{created\ ≤\ S\ \textless{}\ deleted}.

\emph{Proof sketch.} The versions of a key form a chain:
\texttt{v₁(c₁,\ d₁),\ v₂(c₂,\ d₂),\ …} with \texttt{d\_i\ =\ c\_\{i+1\}}
(each delete coincides with the next create), the last version having
\texttt{deleted\ =\ ∞}. The intervals \texttt{{[}c\_i,\ d\_i)} therefore
\textbf{partition} the CSN line from \texttt{c₁} to \texttt{∞}. Any
\texttt{S\ ≥\ c₁} lands in exactly one interval; any
\texttt{S\ \textless{}\ c₁} lands in none (the key did not exist yet). ∎
\end{quote}

\begin{Shaded}
\begin{Highlighting}[]
\NormalTok{flowchart LR}
\NormalTok{    subgraph chain["Version chain of one key — intervals partition the CSN line"]}
\NormalTok{      direction LR}
\NormalTok{      A["[1, 5)"] {-}{-}\textgreater{} B["[5, 9)"] {-}{-}\textgreater{} C["[9, ∞)"]}
\NormalTok{    end}
\NormalTok{    S3(["S=3 → [1,5)"]):::a {-}{-}\textgreater{} A}
\NormalTok{    S7(["S=7 → [5,9)"]):::b {-}{-}\textgreater{} B}
\NormalTok{    S12(["S=12 → [9,∞)"]):::c {-}{-}\textgreater{} C}
\NormalTok{    classDef a fill:\#bde0fe; classDef b fill:\#a2d2ff; classDef c fill:\#8ecae6;}
\end{Highlighting}
\end{Shaded}

\section{Why cold data is free}\label{why-cold-data-is-free-1}

The predicate has a batched fast path that is the reason a large scan is
cheap. Instead of testing every row, a scan tests the \emph{whole part}
first:

\begin{quote}
\textbf{ALGORITHM 4 --- Full-part visibility fast path}

\begin{Shaded}
\begin{Highlighting}[]
\NormalTok{Input:  snapshot S; part P with uniform created stamp c\_max}
\NormalTok{        and minimum deletion m\_del over its rows}
\NormalTok{Output: how to scan P under S}
\NormalTok{1  if c\_max ≤ S and S \textless{} m\_del:                     ▷ every row created ≤ S,}
\NormalTok{2      scan ALL rows of P — no per{-}row check        ▷ none deleted ≤ S}
\NormalTok{3  elif S \textless{} c\_min:                                  ▷ part entirely in the future}
\NormalTok{4      skip P                                        ▷ nothing visible}
\NormalTok{5  else:}
\NormalTok{6      for each row r in P:                          ▷ the slow path}
\NormalTok{7          if ALGORITHM 3 holds and r not tombstoned: emit r}
\end{Highlighting}
\end{Shaded}
\end{quote}

Line 1 is the corollary that matters: a part whose newest creation is
\texttt{≤\ S} and whose earliest deletion is \texttt{\textgreater{}\ S}
is \textbf{fully visible} --- every row passes, so the scan emits the
batch with a single comparison and \emph{zero} per-row work.

\begin{quote}
\textbf{Proposition 3 (Cold-scan cost).} A scan at \texttt{S} pays the
per-row visibility check only on parts modified in the window the
snapshot straddles; cold, unmodified parts pay \texttt{O(1)} per part.

\emph{Proof sketch.} A cold part has a uniform \texttt{created\ ≤\ S}
(it was written long ago) and no deletions after \texttt{S}, so it takes
the fast path (ALG 4, line 1) --- one comparison for the entire part.
Only parts with a \texttt{created} or a \texttt{deleted} inside
\texttt{(S\_min,\ S\_max{]}} fall to the slow path. Hence a billion-row
table with a thousand recent mutations checks a thousand rows, not a
billion --- the ``cost of concurrency is paid only by recently-modified
data'' principle, from Neumann et al.~∎
\end{quote}

\section{Where the stamps live}\label{where-the-stamps-live}

Per-part, ChakraDB stores the version stamps compactly: a part whose
rows all share one CSN stores a single \texttt{Uniform(csn)} rather than
a stamp per row (the common case for a bulk-sealed part), and a
\textbf{deletion vector} records which ordinals are tombstoned and at
which CSN. So the fast path's \texttt{c\_max} and \texttt{m\_del} are
\texttt{O(1)} to read, and the slow path consults the deletion vector
rather than rewriting the part.

\section{The consequence for
concurrency}\label{the-consequence-for-concurrency}

Because visibility is a pure function of a snapshot number and immutable
per-row stamps, a reader needs \textbf{no lock}: it fixes \texttt{S},
grabs the (immutable) part list, and scans. A concurrent writer advances
the clock and appends new versions with higher CSNs --- invisible to
\texttt{S}. This is the mechanism behind ``readers never block
writers,'' and it is why the same snapshot feeds analytics,
transactions, and the \href{../graph/snapshot.md}{graph CSR}
consistently.

\section{The GC Watermark \& Live-Snapshot
Registry}\label{the-gc-watermark-live-snapshot-registry}

Compaction reclaims space by dropping row versions no one can see any
more. Get the ``no one can see'' test wrong and you get the worst kind
of bug: a silent wrong answer for a reader on an old snapshot. This
chapter is how ChakraDB gets it right --- the correctness guarantee
behind the concurrency wedge.

\section{The hazard}\label{the-hazard}

Compaction, when it merges parts, physically reclaims rows whose
deletion CSN is at or before a \textbf{horizon}. The rule is:

\begin{quote}
The horizon must be \textbf{≤ the oldest CSN any live snapshot may still
observe.} Reclaim anything newer, and a reader holding an older snapshot
loses a row it should still see.
\end{quote}

The danger case, if the horizon were just ``now'':

\begin{center}\includegraphics[width=0.92\linewidth,height=0.80\textheight,keepaspectratio]{diagrams/diagram-12.pdf}\end{center}

At \texttt{S=100}, key \texttt{K} (deleted at 150) \emph{should} be
visible. A horizon of ``now'' would have reclaimed it.

\section{The fix: a live-snapshot
registry}\label{the-fix-a-live-snapshot-registry}

ChakraDB tracks the oldest live reader and holds the horizon back to it.
The CSN generator keeps a small registry --- a map from CSN to a count
of live \emph{pins}:

\begin{itemize}
\tightlist
\item
  A read that must outlive a clock advance --- a query or a transaction
  --- \textbf{pins} its snapshot: it registers its CSN and gets an RAII
  guard.
\item
  \textbf{\texttt{gc\_horizon()}} = the smallest pinned CSN, or the
  current clock if none are pinned.
\item
  Compaction reclaims only up to \texttt{gc\_horizon()}.
\end{itemize}

\begin{Shaded}
\begin{Highlighting}[]
\NormalTok{pin():                       gc\_horizon():               on guard drop:}
\NormalTok{  lock registry                lock registry               lock registry}
\NormalTok{  csn = current()              h = min(registry keys)      registry[csn] {-}= 1}
\NormalTok{  registry[csn] += 1              or current() if empty     if 0: remove csn}
\NormalTok{  return guard(csn)            return h                    unlock}
\end{Highlighting}
\end{Shaded}

\begin{center}\includegraphics[width=0.92\linewidth,height=0.80\textheight,keepaspectratio]{diagrams/diagram-13.pdf}\end{center}

\section{The subtle part: no window}\label{the-subtle-part-no-window}

The correctness hinges on one detail: \textbf{reading the current CSN
and registering the pin happen under the same lock that
\texttt{gc\_horizon} uses.} Otherwise a compaction could compute a
horizon \emph{above} a pin that is about to register at a lower CSN ---
and the race is back.

Because both operations are linearized by the one registry lock, a
compaction that runs \emph{after} a pin registers sees that pin (so its
horizon ≤ the pin), and a compaction that runs \emph{before} the pin
registers used a horizon ≤ the pin's snapshot anyway (the clock only
advances). Either way the horizon never exceeds a live snapshot. The
proof is small precisely because the critical section is small.

\section{What pins, and what doesn't}\label{what-pins-and-what-doesnt}

Only reads that could \textbf{outlive a clock advance} need to pin:

\begin{itemize}
\tightlist
\item
  The \textbf{SQL interpreter} pins for the duration of a
  \texttt{SELECT}/\texttt{UPDATE}/\texttt{DELETE}.
\item
  The \textbf{DataFusion bridge} pins for the whole query.
\item
  A \textbf{transaction} pins from \texttt{BEGIN} until commit/rollback.
\item
  A \textbf{\texttt{GraphView}} pins while it builds its CSR.
\end{itemize}

A transient read at \texttt{current} needs no pin --- nothing it can see
is reclaimable, because the horizon can never exceed \texttt{current}.
So the hot path stays lock-free; the registry is touched once per long
read, not per row.

\section{Why compaction is safe to be
lazy}\label{why-compaction-is-safe-to-be-lazy}

Compaction is \textbf{caller-driven} --- there is no background thread.
Combined with the registry, that means space is reclaimed only when you
ask, and only up to the oldest live reader.
\texttt{Storage::compact\_all()} uses \texttt{gc\_horizon()}
automatically; the low-level \texttt{Database::compact\_all(horizon)}
takes an explicit horizon for callers who know a safe one.

\section{Regression-tested}\label{regression-tested}

The guarantee is pinned down by a test: a reader pins a snapshot, a
writer deletes a key at a later CSN, compaction runs --- and the pinned
reader \textbf{still sees the row}; after the pin drops, the row becomes
reclaimable. The test also checks the horizon tracks the oldest of
several concurrent pins (\texttt{tests/gc\_watermark.rs}). It would fail
against the naive ``horizon = now,'' so it genuinely guards the fix.

\section{Transactions}\label{transactions}

\texttt{BEGIN\ …\ COMMIT} gives ChakraDB ACID multi-statement
transactions with \textbf{snapshot isolation} and
\textbf{first-committer-wins} conflict detection. The design is a
private \textbf{overlay} that is invisible until commit, so a crash or
\texttt{ROLLBACK} simply discards it.

\section{The overlay model}\label{the-overlay-model}

A transaction is a private catalog initialized from a committed
snapshot, plus a change-set to replay on commit:

\begin{center}\includegraphics[width=0.92\linewidth,height=0.80\textheight,keepaspectratio]{diagrams/diagram-14.pdf}\end{center}

\begin{itemize}
\tightlist
\item
  \textbf{Reads} inside the transaction see the committed state as of
  \texttt{BEGIN} plus the transaction's own writes (read-your-writes),
  and \emph{nothing} uncommitted from other connections.
\item
  \textbf{Writes} go only into the overlay and the change-set --- never
  the real store or the WAL --- so rolling back is free and a crash
  loses nothing.
\item
  \textbf{COMMIT} replays the change-set to the real store, which logs
  it as \textbf{one} crash-atomic WAL record.
\end{itemize}

\section{Snapshot isolation}\label{snapshot-isolation}

The transaction pins its \texttt{BEGIN} snapshot (which also holds the
\href{mvcc.md}{GC watermark} back, so the versions it may read are not
compacted away). Every statement in the transaction reads that one
consistent instant, extended by its own writes. Other connections'
commits are invisible.

\section{First-committer-wins}\label{first-committer-wins}

At \texttt{COMMIT}, ChakraDB checks whether any key the transaction
wrote was changed by another committed transaction since \texttt{BEGIN}
--- i.e.~whether the key's current committed value differs from what the
transaction saw at \texttt{BEGIN}. If so, the commit fails with a
\textbf{write conflict} and the transaction is aborted; retry it.

\begin{center}\includegraphics[width=0.92\linewidth,height=0.80\textheight,keepaspectratio]{diagrams/diagram-15.pdf}\end{center}

Synthetic-rowid inserts never conflict --- each is a brand-new row with
a fresh key.

\section{Crash-atomic commit}\label{crash-atomic-commit}

Because the whole change-set is handed to the durable backend as one
batch, it is logged as a single \texttt{WalRecord::Txn}. A WAL record is
framed with a length and CRC, so recovery applies it \textbf{all or
nothing}: a crash mid-commit leaves a torn frame that is discarded, and
the transaction never partially appears. This is verified by truncating
the log at every byte and asserting the row count is only ever
``baseline'' or ``baseline + the whole transaction.''

\section{Scope}\label{scope}

Statements in a transaction run on the \textbf{single-table interpreter}
(the overlay is materialized from the committed snapshot on first
touch). Joins and subqueries belong \emph{outside} a transaction --- use
them on the autocommit path where DataFusion is available. DDL
(\texttt{CREATE\ TABLE}) inside a transaction is applied immediately and
is not rolled back in v1.

\section{Example}\label{example}

\begin{Shaded}
\begin{Highlighting}[]
\ControlFlowTok{BEGIN}\NormalTok{;}
  \KeywordTok{UPDATE}\NormalTok{ accounts }\KeywordTok{SET}\NormalTok{ balance }\OperatorTok{=}\NormalTok{ balance }\OperatorTok{{-}} \DecValTok{100} \KeywordTok{WHERE} \KeywordTok{id} \OperatorTok{=} \DecValTok{1}\NormalTok{;}
  \KeywordTok{UPDATE}\NormalTok{ accounts }\KeywordTok{SET}\NormalTok{ balance }\OperatorTok{=}\NormalTok{ balance }\OperatorTok{+} \DecValTok{100} \KeywordTok{WHERE} \KeywordTok{id} \OperatorTok{=} \DecValTok{2}\NormalTok{;}
\KeywordTok{COMMIT}\NormalTok{;   }\CommentTok{{-}{-} both, atomically, or neither}
\end{Highlighting}
\end{Shaded}

Outside an explicit transaction, every statement autocommits.

\chapter{Durability: WAL, Group Commit \&
Recovery}\label{durability-wal-group-commit-recovery}

\epigraph{The palest ink is better than the best memory.}{--- Chinese proverb}

A durable ChakraDB (\texttt{Storage}) guarantees that an acknowledged
write survives a crash. The mechanism is a classic \textbf{write-ahead
log} with \textbf{group commit}, a \textbf{manifest} for the catalog and
part set, and a \textbf{recovery} pass that replays the log. This
chapter is the shape; the \href{durability.md}{WAL algorithm} and
\href{durability.md}{recovery} chapters have the byte-level detail.

\section{Write-ahead logging}\label{write-ahead-logging}

Every mutation is appended to the WAL \textbf{before} it is
acknowledged. A record is framed with a length and a CRC, so a crash
mid-append leaves a \emph{torn} final frame that recovery detects by
checksum and discards --- the log's valid prefix is always recoverable.

\begin{center}\includegraphics[width=0.92\linewidth,height=0.80\textheight,keepaspectratio]{diagrams/diagram-16.pdf}\end{center}

\section{Three durability modes}\label{three-durability-modes}

The trade between latency and the guarantee is a knob:

\begin{longtable}[]{@{}
  >{\raggedright\arraybackslash}p{(\linewidth - 4\tabcolsep) * \real{0.3333}}
  >{\raggedright\arraybackslash}p{(\linewidth - 4\tabcolsep) * \real{0.3333}}
  >{\raggedright\arraybackslash}p{(\linewidth - 4\tabcolsep) * \real{0.3333}}@{}}
\toprule\noalign{}
\begin{minipage}[b]{\linewidth}\raggedright
Mode
\end{minipage} & \begin{minipage}[b]{\linewidth}\raggedright
Guarantee
\end{minipage} & \begin{minipage}[b]{\linewidth}\raggedright
Cost
\end{minipage} \\
\midrule\noalign{}
\endhead
\bottomrule\noalign{}
\endlastfoot
\texttt{Sync} & every write \texttt{fsync}'d before ack & strongest, one
fsync/write \\
\texttt{Group} & concurrent writers share one \texttt{fsync} & strong,
amortized \\
\texttt{Async} & acked before \texttt{fsync} & fastest; a crash may lose
the last unflushed writes \\
\end{longtable}

\textbf{Group commit} is the default sweet spot: when many threads
commit at once, their records are appended and made durable with a
\emph{single} \texttt{fsync}, so throughput scales with concurrency
instead of paying a sync per write.

\section{The manifest and
checkpointing}\label{the-manifest-and-checkpointing}

The WAL grows without bound if never trimmed. A \textbf{checkpoint}
makes the in-memory state durable in parts, records the live part set
and each table's schema in the \textbf{manifest}, and advances a
\texttt{checkpoint\_csn} --- after which the log \emph{before} that
point is reclaimable. Checkpointing is incremental: unchanged parts are
skipped, and parts that only gained tombstones get just those appended.
See \href{durability.md}{Checkpointing}.

\begin{center}\includegraphics[width=0.92\linewidth,height=0.80\textheight,keepaspectratio]{diagrams/diagram-17.pdf}\end{center}

\section{Recovery}\label{recovery}

On open, recovery reconstructs the exact acknowledged state:

\begin{enumerate}
\def\labelenumi{\arabic{enumi}.}
\tightlist
\item
  Read the \textbf{manifest} --- the catalog (tables + schemas) and the
  live part set as of the last checkpoint. Part \emph{indexes} are
  loaded resident; part \emph{data} is faulted lazily on first touch
  (fast reopen).
\item
  Replay the \textbf{WAL} past \texttt{checkpoint\_csn}, applying each
  record to the in-memory tables. A torn final frame is discarded.
\item
  Raise the CSN clock above the highest replayed stamp, so no version
  number is ever reissued.
\end{enumerate}

\begin{center}\includegraphics[width=0.92\linewidth,height=0.80\textheight,keepaspectratio]{diagrams/diagram-18.pdf}\end{center}

\section{How well it's tested}\label{how-well-its-tested}

Durability is the property you cannot get wrong, so it is tested
adversarially:

\begin{itemize}
\tightlist
\item
  \textbf{10,000 randomized crash trials} verify that every acknowledged
  write survives, in all three durability modes
  (\texttt{crash\_consistency}).
\item
  \textbf{Durable SQL} adds tens of thousands more crash trials ---
  millions of acknowledged writes verified across integer, text, and
  keyless schemas (\texttt{durable\_sql\_crash}).
\item
  A committed multi-statement transaction is one WAL record, so
  truncating the log at \emph{every byte} leaves the transaction either
  fully applied or fully absent --- never partial
  (\texttt{torn\_commit\_record\_is\_all\_or\_nothing}).
\end{itemize}

\section{The single I/O seam}\label{the-single-io-seam}

Everything reaches disk through one \texttt{trait\ Io} (see
\href{storage.md}{The I/O Abstraction}) --- a real POSIX backend for
production and an in-memory, fault-injecting backend (\texttt{MemIo})
that makes those thousands of crash trials possible without touching a
real disk. The durability logic is identical over both.

\section{Write-Ahead Logging \& Group
Commit}\label{write-ahead-logging-group-commit}

Durability is the promise that an acknowledged write survives a crash.
ChakraDB keeps it with a write-ahead log: every mutation is framed,
checksummed, and made durable \emph{before} the write is acknowledged.
Group commit amortizes the \texttt{fsync} so throughput scales with
concurrency.

\section{Record framing}\label{record-framing}

Each WAL record is wrapped in a self-describing frame so a torn write is
detectable:

\begin{Shaded}
\begin{Highlighting}[]
\NormalTok{ ┌──────────┬──────────┬──────────────────────────┐}
\NormalTok{ │ len (u32)│ crc (u32)│ payload (len bytes)       │}
\NormalTok{ └──────────┴──────────┴──────────────────────────┘}
\end{Highlighting}
\end{Shaded}

The payload is one of: \texttt{Insert\{table,\ csn,\ row\}},
\texttt{Delete\{table,\ csn,\ key\}},
\texttt{Seal\{table,\ csn,\ part\_id\}}, \texttt{Checkpoint\{csn\}}, or
\texttt{Txn\{ops\}} (a whole committed transaction as one record). The
length lets recovery find the next frame; the CRC lets it reject a frame
torn by a crash mid-append.

\section{The append}\label{the-append}

\begin{quote}
\textbf{ALGORITHM 5 --- WAL append and acknowledge}

\begin{Shaded}
\begin{Highlighting}[]
\NormalTok{Input:  record R; durability mode M ∈ \{Sync, Group, Async\}}
\NormalTok{Output: acknowledgement (the write is durable per M)}
\NormalTok{1  frame ← [len(R)] ++ [crc32(R)] ++ encode(R)}
\NormalTok{2  acquire the append lock                          ▷ assigns byte offsets in order}
\NormalTok{3  append frame to the log at the current end}
\NormalTok{4  release the append lock}
\NormalTok{5  case M of}
\NormalTok{6    Sync:  fsync()                                 ▷ durable before returning}
\NormalTok{7    Group: wait for the group fsync (ALGORITHM 6)  ▷ shared with concurrent writers}
\NormalTok{8    Async: return now                              ▷ a crash may lose the last frames}
\NormalTok{9  return OK}
\end{Highlighting}
\end{Shaded}
\end{quote}

The append lock is held only for the byte-copy that assigns offsets; the
\texttt{fsync} (the expensive part) happens outside it, so writers do
not serialize on the disk.

\section{Group commit}\label{group-commit}

When many threads commit at once, forcing one \texttt{fsync} per write
wastes the disk's throughput --- a single \texttt{fsync} makes
\emph{everything appended so far} durable. Group commit exploits that:

\begin{quote}
\textbf{ALGORITHM 6 --- Group commit}

\begin{Shaded}
\begin{Highlighting}[]
\NormalTok{Input:  a stream of concurrent committers, each having appended its frame}
\NormalTok{Output: one fsync makes a whole batch durable}
\NormalTok{1  each committer records the log end e\_i after its append, then waits}
\NormalTok{2  one committer becomes the LEADER:}
\NormalTok{3      target ← current log end                     ▷ covers every frame ≤ target}
\NormalTok{4      fsync()                                       ▷ a single disk sync}
\NormalTok{5      publish durable{-}through = target}
\NormalTok{6  every waiter with e\_i ≤ target wakes and returns OK}
\end{Highlighting}
\end{Shaded}
\end{quote}

Because an \texttt{fsync} covers all bytes written before it, \texttt{k}
concurrent commits need one sync instead of \texttt{k}. The measured
effect is that syncs-per-append drops sharply as concurrency rises ---
group commit is why durable throughput scales with writers rather than
collapsing to disk-sync latency.

\begin{center}\includegraphics[width=0.92\linewidth,height=0.80\textheight,keepaspectratio]{diagrams/diagram-19.pdf}\end{center}

\section{The three modes}\label{the-three-modes}

\begin{longtable}[]{@{}
  >{\raggedright\arraybackslash}p{(\linewidth - 6\tabcolsep) * \real{0.2500}}
  >{\raggedright\arraybackslash}p{(\linewidth - 6\tabcolsep) * \real{0.2500}}
  >{\raggedright\arraybackslash}p{(\linewidth - 6\tabcolsep) * \real{0.2500}}
  >{\raggedright\arraybackslash}p{(\linewidth - 6\tabcolsep) * \real{0.2500}}@{}}
\toprule\noalign{}
\begin{minipage}[b]{\linewidth}\raggedright
Mode
\end{minipage} & \begin{minipage}[b]{\linewidth}\raggedright
Line
\end{minipage} & \begin{minipage}[b]{\linewidth}\raggedright
Guarantee
\end{minipage} & \begin{minipage}[b]{\linewidth}\raggedright
When to use
\end{minipage} \\
\midrule\noalign{}
\endhead
\bottomrule\noalign{}
\endlastfoot
\texttt{Sync} & 6 & every write \texttt{fsync}'d before ack & strictest;
low write rate \\
\texttt{Group} & 7 & concurrent writers share one \texttt{fsync} &
default; high concurrency \\
\texttt{Async} & 8 & acked before \texttt{fsync} & fastest; a crash may
lose the last unflushed writes \\
\end{longtable}

\section{The crash-atomicity of a
transaction}\label{the-crash-atomicity-of-a-transaction}

A committed transaction is written as \textbf{one} \texttt{Txn} record
--- a single frame. Since a frame is all-or-nothing under the CRC,
recovery applies the whole transaction or none of it.

\begin{quote}
\textbf{Proposition 4 (Commit atomicity).} A crash during the commit of
a transaction leaves the database in a state with either all of the
transaction's writes or none.

\emph{Proof sketch.} The transaction's writes are one framed record. If
the crash occurs before the frame's bytes (and their covering
\texttt{fsync}) are durable, recovery sees no valid frame there (the CRC
fails on the torn tail) and discards it --- \emph{none}. If the frame
and its \texttt{fsync} completed, recovery replays the whole frame ---
\emph{all}. There is no third outcome because the frame is the unit of
durability. This is verified empirically by truncating the log at
\textbf{every} byte and asserting the recovered row count is only ever
the pre-transaction value or the post-transaction value
(\texttt{torn\_commit\_record\_is\_all\_or\_nothing}). ∎
\end{quote}

Recovery --- how the log is replayed to reconstruct the acknowledged
state --- is \href{durability.md}{the next chapter}.

\section{Crash Recovery}\label{crash-recovery}

Recovery reconstructs the exact acknowledged state after a crash: the
catalog from the manifest, the data from the parts, and everything since
the last checkpoint by replaying the WAL --- discarding any torn tail.

\section{The procedure}\label{the-procedure}

\begin{quote}
\textbf{ALGORITHM 7 --- Recovery}

\begin{Shaded}
\begin{Highlighting}[]
\NormalTok{Input:  a database directory (manifest, part files, wal.log)}
\NormalTok{Output: an open database at the last acknowledged state}
\NormalTok{1  M ← read manifest                                 ▷ catalog + live parts + checkpoint\_csn}
\NormalTok{2  for each table T in M:}
\NormalTok{3      rebuild T\textquotesingle{}s schema; register its live parts    ▷ indexes resident, data lazy}
\NormalTok{4  max ← M.checkpoint\_csn}
\NormalTok{5  for each frame F in wal.log, in order:}
\NormalTok{6      if not crc\_ok(F): break                        ▷ torn tail — stop cleanly (ALG 7a)}
\NormalTok{7      if F.csn ≤ M.checkpoint\_csn: continue           ▷ already durable in parts}
\NormalTok{8      apply F to the in{-}memory tables                 ▷ Insert/Delete/Txn; skip unknown tables}
\NormalTok{9      max ← builtin\_max(max, F.csn)}
\NormalTok{10 set the CSN clock floor to max + 1                  ▷ never reissue a stamp}
\NormalTok{11 return the open database}
\end{Highlighting}
\end{Shaded}
\end{quote}

Three things make this correct and fast.

\textbf{Only the tail is replayed.} A checkpoint has already written
everything at or below \texttt{checkpoint\_csn} into parts (recorded in
the manifest), so recovery skips those frames (line 7) and replays only
what came after --- usually a small tail.

\textbf{Parts load lazily.} Line 3 registers each part's \emph{index}
(Bloom, bounds, version stamps, deletion vector) resident, but leaves
its \emph{column data} on disk until first touched. Reopening a large
database is therefore near-instant --- flat in the number of parts, not
their bytes.

\textbf{Unknown tables are skipped.} If a frame references a table not
in the manifest (it was dropped, or a \texttt{TRUNCATE} gave it a fresh
id), line 8 ignores it. This is what makes \texttt{DROP\ TABLE} and
\texttt{TRUNCATE} durable without a special log record --- see
\href{../getting-started/getting-started.md}{the SQL surface}.

\section{The torn tail}\label{the-torn-tail}

A crash mid-append leaves a partial final frame. Recovery detects it by
checksum and stops --- the log's valid prefix is always recoverable:

\begin{quote}
\textbf{ALGORITHM 7a --- Torn-tail detection (line 6, expanded)}

\begin{Shaded}
\begin{Highlighting}[]
\NormalTok{Input:  the byte at the current read position in the log}
\NormalTok{Output: the next valid frame, or STOP}
\NormalTok{1  if fewer than 8 bytes remain: STOP                 ▷ no room for len+crc}
\NormalTok{2  len ← read u32;  crc ← read u32}
\NormalTok{3  if fewer than len bytes remain: STOP               ▷ payload truncated}
\NormalTok{4  payload ← read len bytes}
\NormalTok{5  if crc32(payload) ≠ crc: STOP                       ▷ torn or corrupt frame}
\NormalTok{6  return decode(payload)}
\end{Highlighting}
\end{Shaded}
\end{quote}

Everything after the first bad frame is discarded --- those writes were
never acknowledged (their covering \texttt{fsync} had not completed), so
dropping them is correct.

\begin{Shaded}
\begin{Highlighting}[]
\NormalTok{flowchart LR}
\NormalTok{    subgraph log["wal.log after a crash"]}
\NormalTok{      F1["frame 1 ✓"] {-}{-}\textgreater{} F2["frame 2 ✓"] {-}{-}\textgreater{} F3["frame 3 ✓"] {-}{-}\textgreater{} T["torn frame ✗"]}
\NormalTok{    end}
\NormalTok{    F1 \& F2 \& F3 {-}{-}\textgreater{}|replayed| DB[("recovered state")]}
\NormalTok{    T {-}.{-}\textgreater{}|discarded (crc fails)| X["(dropped)"]}
\NormalTok{    style T fill:\#f5d6d6}
\end{Highlighting}
\end{Shaded}

\section{Why it recovers the exact acknowledged
state}\label{why-it-recovers-the-exact-acknowledged-state}

\begin{quote}
\textbf{Proposition 5 (Recovery completeness).} After recovery, the
database contains every acknowledged write and no unacknowledged one.

\emph{Proof sketch.} An \emph{acknowledged} write is, by the durability
mode (\href{durability.md}{ALGORITHM 5}), one whose frame and covering
\texttt{fsync} completed --- so its frame is intact and precedes the
torn tail, and it is replayed (or already in a part below the
checkpoint). An \emph{unacknowledged} write either never reached the log
or sits in the torn tail; the CRC rejects the latter, so it is dropped.
Hence the recovered set is exactly the acknowledged set. This is
stress-tested by tens of thousands of randomized crash trials across all
durability modes and schemas (\texttt{crash\_consistency},
\texttt{durable\_sql\_crash}). ∎
\end{quote}

\section{Clock safety}\label{clock-safety}

Line 10 raises the CSN clock above the highest replayed stamp. Without
it, a recovered database could hand out a CSN a replayed row already
used, and two rows would collide on a version number --- corrupting
\href{mvcc.md}{visibility}. The floor makes recovered CSNs safe: the
next allocation is strictly greater than anything on disk.

\chapter{The Query Layer: The HTAP
Router}\label{the-query-layer-the-htap-router}

\epigraph{If the only tool you have is a hammer, it is tempting to treat everything as if it were a nail.}{--- Abraham Maslow}

ChakraDB runs \textbf{two} execution engines and routes each statement
to whichever fits. This is the concrete shape of ``buy execution, build
storage'': DataFusion is bought for vectorized analytics; a compact
interpreter is built for the transactional and point-query path it wins.

\section{Why two engines}\label{why-two-engines}

Different query shapes have opposite ideal engines:

\begin{itemize}
\tightlist
\item
  A \textbf{key point-lookup} (\texttt{WHERE\ id\ =\ 42}) wants an index
  funnel --- a Bloom probe and a binary search --- not a vectorized scan
  operator. The interpreter wins by orders of magnitude.
\item
  A \textbf{\texttt{COUNT(*)}} or \textbf{bare
  \texttt{MIN}/\texttt{MAX}} wants to be answered from metadata (row
  counts, zonemaps) with \emph{no scan at all}. The interpreter does
  exactly that.
\item
  A \textbf{large \texttt{GROUP\ BY} / join / window} wants vectorized,
  spill-capable operators. DataFusion wins.
\item
  A \textbf{write}
  (\texttt{INSERT}/\texttt{UPDATE}/\texttt{DELETE}/\texttt{COPY}) must
  own the WAL and the snapshot clock. Only the interpreter writes.
\end{itemize}

Routing to a single engine would lose one side of this. So ChakraDB has
a \textbf{cost-based router} that inspects the planned statement.

\section{The routing decision}\label{the-routing-decision}

\begin{center}\includegraphics[width=0.92\linewidth,height=0.80\textheight,keepaspectratio]{diagrams/diagram-21.pdf}\end{center}

The first stage asks whether the interpreter can even plan the
statement; joins, subqueries, and window functions it cannot, so they go
straight to DataFusion. For the rest, a second cost stage checks the
cheap-shape shortcuts --- metadata answers, point lookups, and highly
selective ranges that \href{storage.md}{zonemap pruning} makes cheaper
on the interpreter than a full vectorized scan --- and everything else
takes the vectorized path.

\section{The interpreter half}\label{the-interpreter-half}

A hand-written, allocation-conscious evaluator over the sorted parts.
It:

\begin{itemize}
\tightlist
\item
  resolves columns to indices at plan time (no per-row name lookup),
\item
  evaluates predicates and projections columnar, over the Arrow batches
  in place,
\item
  answers \texttt{COUNT(*)}/\texttt{MIN}/\texttt{MAX} from part
  metadata,
\item
  funnels point lookups through bounds → Bloom → binary search,
\item
  prunes parts on selective ranges,
\item
  and is the \textbf{only} path that mutates state and writes the WAL.
\end{itemize}

It is also the \textbf{zero-heavy-dependency fallback}: build without
the \texttt{datafusion} feature and the engine still runs single-table
SQL on the interpreter alone.

\section{The DataFusion half}\label{the-datafusion-half}

DataFusion receives the raw SQL and executes it over an Arrow view of
the current MVCC snapshot --- \textbf{zero-copy}, since sealed parts
already \emph{are} Arrow record batches, handed over by reference. It
brings the large SQL surface (joins, subqueries, windows, a rich
function library) and vectorized, spill-capable operators. ChakraDB pins
the snapshot for the query's duration, so DataFusion sees a consistent
read while writers keep committing.

\begin{quote}
\textbf{A subtlety worth stating.} Because the interpreter owns writes
and DDL, a statement that fails to \emph{plan} on the interpreter ---
say, a constraint violation --- is a real error and is surfaced, not
silently retried on DataFusion. Only a \emph{SELECT} the interpreter
cannot plan falls through to the vectorized engine.
\end{quote}

\section{Transactions run on the
interpreter}\label{transactions-run-on-the-interpreter}

Inside a \texttt{BEGIN\ …\ COMMIT}, statements run on the interpreter
against a private overlay (read-your-writes; nothing hits the WAL until
commit). Joins and subqueries belong \emph{outside} a transaction ---
the transactional path is single-table by design. See
\href{mvcc.md}{Transactions}.

\section{Why this is the right split}\label{why-this-is-the-right-split}

Building a world-class vectorized engine is a multi-year effort on the
axis ChakraDB \emph{deliberately does not compete on} (raw scan speed).
Buying DataFusion spends those years on storage, MVCC, durability, and
graph instead --- while the interpreter captures the transactional
shapes DataFusion is simply the wrong tool for. The
\href{../comparisons/comparisons.md}{DuckDB comparison} shows where each
lands.

\section{Query Routing (Cost Model)}\label{query-routing-cost-model}

The HTAP router decides, for each statement, whether to run it on the
hand-written \textbf{interpreter} or on \textbf{DataFusion}. The
decision is cheap --- it inspects the \emph{planned} statement, not the
data --- and it is the concrete cost model that lets one engine serve
both transactional and analytical shapes well.

\section{The two-stage decision}\label{the-two-stage-decision}

\begin{quote}
\textbf{ALGORITHM 12 --- Route a statement}

\begin{Shaded}
\begin{Highlighting}[]
\NormalTok{Input:  SQL statement text}
\NormalTok{Output: the engine to execute it on}
\NormalTok{1  plan ← interpreter.plan(statement)}
\NormalTok{2  if plan failed:                                   ▷ join / subquery / window}
\NormalTok{3      if statement is a SELECT: return DataFusion    ▷ it can plan these}
\NormalTok{4      else: raise the error                          ▷ a write/DDL error is real}
\NormalTok{5  if statement is a write or DDL: return Interpreter ▷ owns WAL + snapshot clock}
\NormalTok{6  if plan is bare COUNT(*) with no filter:  return Interpreter  ▷ metadata, no scan}
\NormalTok{7  if plan is bare MIN(c)/MAX(c):            return Interpreter  ▷ zonemaps, no scan}
\NormalTok{8  if plan is a key point lookup:            return Interpreter  ▷ index funnel}
\NormalTok{9  if plan is a selective range that prunes ≥ 80\% of rows across ≥ 4 parts:}
\NormalTok{10     return Interpreter                             ▷ pruned scan beats vectorized}
\NormalTok{11 return DataFusion                                  ▷ scans, joins, aggregation}
\end{Highlighting}
\end{Shaded}
\end{quote}

The first stage (lines 1--4) is a capability check: joins, subqueries,
and windows the interpreter cannot plan, so a \texttt{SELECT} using them
falls through to DataFusion, while a \emph{write} that fails to plan
(e.g.~a constraint violation) is a real error and is surfaced --- never
silently retried on the read-only engine.

The second stage (lines 5--11) is the cost model: the shapes the
interpreter \emph{wins} are captured explicitly, and everything else
takes the vectorized path.

\begin{center}\includegraphics[width=0.92\linewidth,height=0.80\textheight,keepaspectratio]{diagrams/diagram-22.pdf}\end{center}

\section{Why each branch}\label{why-each-branch}

\begin{itemize}
\tightlist
\item
  \textbf{Writes / DDL (line 5).} Only the interpreter mutates state and
  appends to the WAL; it owns the \href{mvcc.md}{snapshot clock}.
\item
  \textbf{Bare \texttt{COUNT(*)} (line 6).} Answered from part
  row-counts --- no scan at all.
\item
  \textbf{Bare \texttt{MIN}/\texttt{MAX} (line 7).} Answered from the
  per-part zonemaps --- no scan.
\item
  \textbf{Point lookup (line 8).} The \href{storage.md}{index funnel} is
  \texttt{O(log\ n)}; a vectorized scan operator would be far slower for
  a single row.
\item
  \textbf{Selective range (lines 9--10).} \href{storage.md}{Zonemap
  pruning} makes the interpreter touch only a few parts; below
  \textasciitilde20\% selectivity the pruned interpreter scan beats a
  full vectorized scan. The threshold is a cost estimate, not a scan ---
  it counts, from the zonemaps alone, how many rows survive pruning.
\item
  \textbf{Everything else (line 11).} Large scans, \texttt{GROUP\ BY},
  joins, windows, and subqueries want DataFusion's vectorized,
  spill-capable operators.
\end{itemize}

\section{The estimate is
metadata-only}\label{the-estimate-is-metadata-only}

Line 9's ``prunes ≥ 80\%'' is computed without reading a row: for each
part, the \texttt{excludes} test (\href{storage.md}{ALGORITHM 11}) plus
the part's row count says how many rows survive. Summing over parts
gives the surviving fraction. So the routing decision itself is
\texttt{O(number\ of\ parts)}, not \texttt{O(rows)}.

\begin{quote}
\textbf{Proposition 8 (Routing is answer-preserving).} For any
statement, the interpreter and DataFusion produce the same result;
routing changes only cost.

\emph{Proof sketch.} Both engines read the same MVCC snapshot and the
same immutable parts. The interpreter's fast paths are exact ---
\texttt{COUNT(*)} from true row counts, \texttt{MIN}/\texttt{MAX} from
exact zonemaps, point lookup from the sorted key, pruned scan from a
conservative \texttt{excludes} that never drops a match (ALG 11).
DataFusion evaluates the same relational semantics over the same Arrow
data. Hence the two agree on every routed statement; the
HTAP-equivalence property tests assert exactly this across randomized
queries. ∎
\end{quote}

\section{The pin}\label{the-pin}

Whichever engine runs, the statement first \textbf{pins} its snapshot
for the duration (so compaction cannot reclaim a version it will read
--- the \href{mvcc.md}{GC watermark}). Transient reads at the current
clock need no pin; long reads (a full scan, a DataFusion query, a
transaction) hold one.

\chapter{Exact Decimal Arithmetic}\label{exact-decimal-arithmetic}

\epigraph{Round numbers are always false.}{--- Samuel Johnson}

Money must not round. ChakraDB's \texttt{DECIMAL(p,\ s)} is exact
fixed-point --- an \texttt{i128} unscaled mantissa with a scale, backed
by Arrow \texttt{Decimal128}, never \texttt{f64}. This chapter is how
values are parsed, compared, and summed without ever touching a float.

\section{Representation}\label{representation}

A decimal is \texttt{(mantissa,\ scale)}: the number is
\texttt{mantissa\ /\ 10\^{}scale}. So \texttt{12.34} is
\texttt{Decimal(1234,\ 2)}, and \texttt{-0.01} is
\texttt{Decimal(-1,\ 2)}. \texttt{DECIMAL(p,\ s)} bounds the
\textbf{precision} \texttt{p} (total significant digits,
\texttt{p\ ≤\ 38} for \texttt{i128}) and the \textbf{scale} \texttt{s}.
Values are stored at the column's scale.

\section{Exact parsing from text}\label{exact-parsing-from-text}

The exactness starts at ingestion: a decimal literal is parsed from its
\textbf{source text}, never via \texttt{f64}. \texttt{9.99} becomes
\texttt{Decimal(999,\ 2)} directly --- not the nearest double
\texttt{9.9900000000000002…} rounded back.

\begin{quote}
\textbf{ALGORITHM 13 --- Parse and fit a decimal literal}

\begin{Shaded}
\begin{Highlighting}[]
\NormalTok{Input:  the literal\textquotesingle{}s text t; target type DECIMAL(p, s)}
\NormalTok{Output: Decimal(m, s), or REJECT}
\NormalTok{1  (digits, point) ← split t on \textquotesingle{}.\textquotesingle{}                  ▷ integer part, fraction}
\NormalTok{2  from ← len(fraction);  raw ← integer ++ fraction as i128   ▷ exact, base{-}10}
\NormalTok{3  m ← Rescale(raw, from, s)                          ▷ align to the column scale}
\NormalTok{4  if |m| ≥ 10\^{}p: REJECT                              ▷ exceeds declared precision}
\NormalTok{5  return Decimal(m, s)}
\end{Highlighting}
\end{Shaded}
\end{quote}

Line 4 enforces precision: \texttt{1000} into \texttt{DECIMAL(3,0)} (max
\texttt{999}) is rejected, not silently stored --- a real correctness
guard for money.

Rescaling to a target scale is exact when growing, and rounds
half-away-from-zero when shrinking (SQL \texttt{CAST} rounding):

\begin{quote}
\begin{Shaded}
\begin{Highlighting}[]
\NormalTok{Rescale(m, from, to):}
\NormalTok{  if to = from: return m}
\NormalTok{  if to \textgreater{} from: return m · 10\^{}(to−from)              ▷ exact}
\NormalTok{  div ← 10\^{}(from−to);  half ← div/2                   ▷ shrinking → round}
\NormalTok{  return (m + sign(m)·half) / div}
\end{Highlighting}
\end{Shaded}
\end{quote}

\section{Exact comparison}\label{exact-comparison}

Two decimals of different scales are compared by aligning to the larger
scale in \texttt{i128} --- never through \texttt{f64}, which would
conflate values a distributed ledger must distinguish:

\begin{quote}
\begin{Shaded}
\begin{Highlighting}[]
\NormalTok{cmp(Decimal(a, sa), Decimal(b, sb)):}
\NormalTok{  s ← max(sa, sb)}
\NormalTok{  return compare\_i128( a·10\^{}(s−sa),  b·10\^{}(s−sb) )    ▷ exact; f64 only on overflow}
\end{Highlighting}
\end{Shaded}
\end{quote}

So \texttt{10.00} and \texttt{10.0} and the integer \texttt{10} all
compare equal, exactly. Decimals also order correctly against integers
(widened to scale 0).

\section{Exact aggregation and
arithmetic}\label{exact-aggregation-and-arithmetic}

\texttt{SUM} of a decimal column accumulates the \texttt{i128} mantissa
at the shared scale and stays exact --- the interpreter's accumulator
keeps a running \texttt{i128} while every value is a decimal of one
scale (falling back to \texttt{f64} only on a genuine mix).
\texttt{MIN}/\texttt{MAX} are exact via the comparison above.
Arithmetic:

\begin{itemize}
\tightlist
\item
  \texttt{+} / \texttt{−} --- align scales, add/subtract mantissas
  (checked; \texttt{f64} only on \texttt{i128} overflow).
\item
  \texttt{×} --- multiply mantissas, add scales:
  \texttt{Decimal(a,\ sa)\ ×\ Decimal(b,\ sb)\ =\ \ \ Decimal(a·b,\ sa+sb)}.
\item
  \texttt{÷}, \texttt{\%} --- fall back to \texttt{f64} (exact decimal
  division has an implementation-defined scale; documented as the one
  approximate operator).
\item
  \texttt{AVG} --- returns a float (a mean is inherently fractional).
\end{itemize}

\begin{quote}
\textbf{Proposition 9 (Exactness of storage, comparison, and sum).} For
decimals that fit \texttt{DECIMAL(p,\ s)}, round-trip storage,
comparison, and \texttt{SUM} introduce no error.

\emph{Proof sketch.} Storage keeps the exact base-10 mantissa (ALG 13,
parsed from text, no \texttt{f64}). Comparison aligns scales by exact
\texttt{i128} multiplication. Summation adds exact \texttt{i128}
mantissas at a common scale. None of these operations uses floating
point, so no rounding occurs. The famous witness ---
\texttt{0.1\ +\ 0.2}, which is \texttt{0.30000000000000004} in
\texttt{f64} --- stores and sums to exactly \texttt{0.3}
(\texttt{tests/decimal.rs}). ∎
\end{quote}

\section{Where the float boundary is}\label{where-the-float-boundary-is}

The only places a decimal meets \texttt{f64} are division/modulo,
\texttt{AVG}, and an explicit cast to \texttt{FLOAT}. Everywhere a
ledger cares --- INSERT, comparison, \texttt{SUM},
\texttt{MIN}/\texttt{MAX}, \texttt{+}/\texttt{−}/\texttt{×} --- the
arithmetic is exact \texttt{i128}. That boundary is stated plainly so no
one is surprised by a rounded cent.

\section{Temporal Encoding}\label{temporal-encoding}

\texttt{DATE} and \texttt{TIMESTAMP} are \textbf{logical types over
integer storage}: a date is a count of days since the Unix epoch, a
timestamp a count of microseconds. That choice keeps comparison,
zonemaps, keys, and MVCC working unchanged on integers, while exposing
the columns to Arrow (and DataFusion) as their native temporal types.

\section{The representation}\label{the-representation}

\begin{longtable}[]{@{}
  >{\raggedright\arraybackslash}p{(\linewidth - 6\tabcolsep) * \real{0.2500}}
  >{\raggedright\arraybackslash}p{(\linewidth - 6\tabcolsep) * \real{0.2500}}
  >{\raggedright\arraybackslash}p{(\linewidth - 6\tabcolsep) * \real{0.2500}}
  >{\raggedright\arraybackslash}p{(\linewidth - 6\tabcolsep) * \real{0.2500}}@{}}
\toprule\noalign{}
\begin{minipage}[b]{\linewidth}\raggedright
Type
\end{minipage} & \begin{minipage}[b]{\linewidth}\raggedright
Stored as
\end{minipage} & \begin{minipage}[b]{\linewidth}\raggedright
Arrow type
\end{minipage} & \begin{minipage}[b]{\linewidth}\raggedright
Literal
\end{minipage} \\
\midrule\noalign{}
\endhead
\bottomrule\noalign{}
\endlastfoot
\texttt{DATE} & \texttt{i64} days since 1970-01-01 & \texttt{Date32} &
\texttt{\textquotesingle{}YYYY-MM-DD\textquotesingle{}} or
\texttt{DATE\ \textquotesingle{}…\textquotesingle{}} \\
\texttt{TIMESTAMP} & \texttt{i64} µs since the epoch &
\texttt{Timestamp(µs)} &
\texttt{\textquotesingle{}YYYY-MM-DD{[}\ T{]}HH:MM:SS{[}.ffffff{]}\textquotesingle{}} \\
\end{longtable}

Because the physical value is an integer, a \texttt{DATE} column sorts,
prunes, and compares exactly like an integer column --- a date range
\texttt{WHERE\ d\ \textgreater{}=\ \textquotesingle{}2024-01-01\textquotesingle{}}
prunes via \href{storage.md}{zonemaps}, and a date can be a primary key.

\section{Civil ↔ days}\label{civil-days}

The conversion between a calendar date and a day number uses Howard
Hinnant's proleptic-Gregorian algorithm --- exact, branch-light, and
dependency-free (the core crate forbids \texttt{unsafe} and has no
\texttt{chrono}).

\begin{quote}
\textbf{ALGORITHM 14 --- Days from a civil date}

\begin{Shaded}
\begin{Highlighting}[]
\NormalTok{Input:  year y, month m, day d}
\NormalTok{Output: days since 1970{-}01{-}01}
\NormalTok{1  y ← y − (m ≤ 2 ? 1 : 0)                            ▷ March{-}based year}
\NormalTok{2  era ← (y ≥ 0 ? y : y−399) / 400}
\NormalTok{3  yoe ← y − era·400                                   ▷ year of era  [0,399]}
\NormalTok{4  doy ← (153·(m \textgreater{} 2 ? m−3 : m+9) + 2)/5 + d − 1       ▷ day of year  [0,365]}
\NormalTok{5  doe ← yoe·365 + yoe/4 − yoe/100 + doy               ▷ day of era   [0,146096]}
\NormalTok{6  return era·146097 + doe − 719468}
\end{Highlighting}
\end{Shaded}
\end{quote}

The inverse (\texttt{days\ →\ (y,\ m,\ d)}) is the mirror image and is
used for rendering. A timestamp splits into
\texttt{days\ ·\ 86\_400\_000\_000\ +\ micros\_of\_day}; rendering uses
floored division so pre-epoch instants format correctly.

\section{Parsing is bounded}\label{parsing-is-bounded}

A hostile literal must not overflow the integer math. Parsing bounds the
year and uses checked arithmetic, so an out-of-range date is
\emph{rejected}, never wrapped:

\begin{quote}
\begin{Shaded}
\begin{Highlighting}[]
\NormalTok{parse\_date(t):}
\NormalTok{  (y, m, d) ← split t on \textquotesingle{}{-}\textquotesingle{}}
\NormalTok{  if |y| \textgreater{} 262143 or m ∉ [1,12] or d ∉ [1,31]: REJECT   ▷ bounded → no overflow}
\NormalTok{  return days\_from\_civil(y, m, d)}
\NormalTok{parse\_timestamp(t):}
\NormalTok{  days ← parse\_date(date\_part);  µs ← parse\_time(time\_part)}
\NormalTok{  return checked(days · 86\_400\_000\_000 + µs)            ▷ None on overflow → REJECT}
\end{Highlighting}
\end{Shaded}
\end{quote}

The bound (\texttt{\textbar{}year\textbar{}\ ≤\ 262143}) comfortably
covers the range Arrow \texttt{Timestamp} can represent and keeps every
downstream multiplication inside \texttt{i64}. Before this bound,
\texttt{DATE\ \textquotesingle{}300000-01-01\textquotesingle{}}
overflowed --- a debug panic, and worse, a silent wrap to a garbage
epoch in release. Now it errors cleanly.

\section{Rendering, consistently across
engines}\label{rendering-consistently-across-engines}

A \texttt{DATE}/\texttt{TIMESTAMP} column stores an integer, but must
\emph{display} as a date string. Two paths render it, and they agree:

\begin{itemize}
\tightlist
\item
  \textbf{DataFusion} reads the native
  \texttt{Date32}/\texttt{Timestamp} Arrow array and renders it.
\item
  \textbf{The interpreter} threads each projection's output type so a
  bare temporal column renders through the civil-from-days formatter,
  not as a raw integer.
\end{itemize}

So \texttt{SELECT\ hire\_date\ FROM\ emp} shows \texttt{2024-01-15}
whether the query ran on the interpreter (a point lookup) or on
DataFusion (a scan) --- the tests assert both paths, and a hostile
literal errors instead of panicking.

\section{Why logical-over-integer}\label{why-logical-over-integer}

\begin{quote}
\textbf{Proposition 10 (Temporal correctness is inherited).} Ordering,
zonemap pruning, and MVCC visibility are correct for
\texttt{DATE}/\texttt{TIMESTAMP} with no new code.

\emph{Proof sketch.} Each is defined on the physical value.
Dates/timestamps are stored as monotonically-increasing integers (a
later instant is a larger integer), so integer ordering \emph{is}
chronological ordering; the zonemap min/max are the earliest/latest
instants; and visibility compares CSNs, untouched by the column type.
The only temporal-specific code is text↔integer conversion at the parse
and render boundaries (ALG 14). ∎
\end{quote}

\part{The Graph Engine}

\chapter{Graphs on ChakraDB}\label{graphs-on-chakradb}

\epigraph{Only connect!}{--- E. M. Forster}

ChakraDB's graph layer is not a second database bolted onto a relational
one. It is a thin layer that \emph{exposes what the engine already is}.
Three properties we built for HTAP turn out to be exactly what a graph
engine wants.

\begin{center}\includegraphics[width=0.92\linewidth,height=0.80\textheight,keepaspectratio]{diagrams/diagram-23.pdf}\end{center}

\section{The three properties}\label{the-three-properties}

\begin{enumerate}
\def\labelenumi{\arabic{enumi}.}
\item
  \textbf{Clustered adjacency, for free.} If an edge's key encodes
  \texttt{(src,\ dst)} src-major, then all of a node's out-edges are one
  \emph{contiguous key range}. ``Neighbors of X'' becomes a key-range
  scan that \href{../engine/storage.md}{zonemap pruning} answers by
  touching only the parts holding \texttt{src\ =\ X}. No secondary
  index; the primary-key sort \emph{is} the adjacency index. See
  \href{model.md}{Clustered Adjacency}.
\item
  \textbf{Live graph analytics.} A whole-graph algorithm builds its
  working set from \textbf{one MVCC snapshot}. Writers keep adding
  edges; the algorithm sees a consistent graph and runs to completion.
  This is the \href{../engine/mvcc.md}{concurrency wedge} applied to
  graphs --- the thing pure graph libraries (a dead static copy) and
  lock-based graph databases cannot do in one embedded process. See
  \href{snapshot.md}{Live Graph Analytics}.
\item
  \textbf{Arrow ≈ CSR.} Edges stored sorted by \texttt{src} are already
  \emph{grouped by source}, so building the CSR (Compressed Sparse Row)
  form every algorithm wants is a single linear, vectorized scan over
  Arrow columns --- no sort, no hash join. See \href{snapshot.md}{The
  CSR Snapshot}.
\end{enumerate}

\section{What the client writes}\label{what-the-client-writes}

\begin{Shaded}
\begin{Highlighting}[]
\KeywordTok{use} \PreprocessorTok{chakradb::}\OperatorTok{\{}\NormalTok{Database}\OperatorTok{,}\NormalTok{ Graph}\OperatorTok{\};}
\KeywordTok{use} \PreprocessorTok{std::sync::}\NormalTok{Arc}\OperatorTok{;}

\KeywordTok{let}\NormalTok{ g }\OperatorTok{=} \PreprocessorTok{Graph::}\NormalTok{open(}\PreprocessorTok{Arc::}\NormalTok{new(}\PreprocessorTok{Database::}\NormalTok{new())}\OperatorTok{,} \StringTok{"social"}\NormalTok{)}\OperatorTok{?;}
\NormalTok{g}\OperatorTok{.}\NormalTok{add\_edge(}\DecValTok{1}\OperatorTok{,} \DecValTok{2}\OperatorTok{,} \DecValTok{1.0}\NormalTok{)}\OperatorTok{?;}                 \CommentTok{// src {-}\textgreater{} dst, weight}
\NormalTok{g}\OperatorTok{.}\NormalTok{add\_edges([(}\DecValTok{2}\OperatorTok{,} \DecValTok{3}\OperatorTok{,} \DecValTok{1.0}\NormalTok{)}\OperatorTok{,}\NormalTok{ (}\DecValTok{1}\OperatorTok{,} \DecValTok{3}\OperatorTok{,} \DecValTok{1.0}\NormalTok{)])}\OperatorTok{?;}

\CommentTok{// Live adjacency (pruned range scan):}
\KeywordTok{let}\NormalTok{ nbrs }\OperatorTok{=}\NormalTok{ g}\OperatorTok{.}\NormalTok{out\_neighbors(}\DecValTok{1}\NormalTok{)}\OperatorTok{?;}         \CommentTok{// [2, 3]}

\CommentTok{// Whole{-}graph algorithms over a consistent snapshot:}
\KeywordTok{let}\NormalTok{ view   }\OperatorTok{=}\NormalTok{ g}\OperatorTok{.}\NormalTok{view()}\OperatorTok{?;}
\KeywordTok{let}\NormalTok{ ranks  }\OperatorTok{=}\NormalTok{ view}\OperatorTok{.}\NormalTok{pagerank(}\DecValTok{20}\OperatorTok{,} \DecValTok{0.85}\NormalTok{)}\OperatorTok{;}   \CommentTok{// node {-}\textgreater{} score}
\KeywordTok{let}\NormalTok{ path   }\OperatorTok{=}\NormalTok{ view}\OperatorTok{.}\NormalTok{shortest\_path(}\DecValTok{1}\OperatorTok{,} \DecValTok{3}\NormalTok{)}\OperatorTok{;}  \CommentTok{// Some([1, 3])}
\KeywordTok{let}\NormalTok{ comps  }\OperatorTok{=}\NormalTok{ view}\OperatorTok{.}\NormalTok{connected\_components()}\OperatorTok{;}
\KeywordTok{let}\NormalTok{ tris   }\OperatorTok{=}\NormalTok{ view}\OperatorTok{.}\NormalTok{triangle\_count()}\OperatorTok{;}
\end{Highlighting}
\end{Shaded}

The client never writes a traversal by hand and never manages a lock.
The engine handles adjacency (pruned scans), consistency (snapshots),
and the algorithms.

\section{Positioning}\label{positioning}

This is the \textbf{G} in ``ChakraDB --- the embedded HTAP database with
built-in graph capabilities.'' Transactions, analytics, and graph
traversals run over the \emph{same live data}, on the \emph{same
non-blocking snapshots}, in \emph{one process}. The rest of this part
builds the layer up: modeling, adjacency, CSR, and the algorithms.

\begin{quote}
\textbf{Status.} The graph layer is real and tested
(\texttt{src/graph.rs}, \texttt{tests/graph.rs}) --- directed edges,
adjacency, CSR views, and the core algorithm set. Node ids are currently
\texttt{u32}; a native composite key (on the roadmap) will lift that and
remove the key-encoding step. See \href{model.md}{Modeling a Property
Graph}.
\end{quote}

\section{Modeling a Property Graph}\label{modeling-a-property-graph}

A property graph is \textbf{nodes and edges}, each with a label/type and
properties. On ChakraDB it maps to ordinary tables --- with one twist
that makes traversal fast.

\section{The tables}\label{the-tables}

\begin{Shaded}
\begin{Highlighting}[]
\CommentTok{{-}{-} Nodes: keyed by node id → fast point lookup and id{-}range scan.}
\KeywordTok{CREATE} \KeywordTok{TABLE}\NormalTok{ nodes (}
  \KeywordTok{id}     \DataTypeTok{INT} \KeywordTok{PRIMARY} \KeywordTok{KEY}\NormalTok{,}
  \KeywordTok{label}  \DataTypeTok{VARCHAR}\NormalTok{(}\DecValTok{64}\NormalTok{),}
\NormalTok{  props  TEXT               }\CommentTok{{-}{-} JSON, or promote hot properties to real columns}
\NormalTok{);}

\CommentTok{{-}{-} Edges: keyed (src, dst) src{-}major → clustered adjacency by source.}
\KeywordTok{CREATE} \KeywordTok{TABLE}\NormalTok{ edges (}
  \KeywordTok{key}    \DataTypeTok{INT} \KeywordTok{PRIMARY} \KeywordTok{KEY}\NormalTok{,   }\CommentTok{{-}{-} encode(src, dst); see Clustered Adjacency}
\NormalTok{  src    }\DataTypeTok{INT}\NormalTok{,}
\NormalTok{  dst    }\DataTypeTok{INT}\NormalTok{,}
  \KeywordTok{type}   \DataTypeTok{VARCHAR}\NormalTok{(}\DecValTok{32}\NormalTok{),}
\NormalTok{  weight }\DataTypeTok{DOUBLE}
\NormalTok{);}
\end{Highlighting}
\end{Shaded}

The \texttt{Graph} handle manages the edge encoding for you; the schema
above is what it creates under the hood (\texttt{\{name\}\_edges}). A
companion \texttt{nodes} table for node properties is optional in v1 ---
nodes are otherwise implied by the ids that appear in edges.

\section{Hot vs.~cold properties}\label{hot-vs.-cold-properties}

Put \textbf{hot} edge properties --- the ones you filter on --- in real
typed columns. They get \href{../engine/storage.md}{zonemaps}, so a
query like ``follow-edges created after T'' prunes:

\begin{Shaded}
\begin{Highlighting}[]
\KeywordTok{SELECT}\NormalTok{ dst }\KeywordTok{FROM}\NormalTok{ edges}
\KeywordTok{WHERE} \KeywordTok{key} \OperatorTok{\textgreater{}=}\NormalTok{ (X}\OperatorTok{\textless{}\textless{}}\DecValTok{32}\NormalTok{) }\KeywordTok{AND} \KeywordTok{key} \OperatorTok{\textless{}}\NormalTok{ ((X}\OperatorTok{+}\DecValTok{1}\NormalTok{)}\OperatorTok{\textless{}\textless{}}\DecValTok{32}\NormalTok{)   }\CommentTok{{-}{-} neighbors of X (clustered)}
  \KeywordTok{AND} \KeywordTok{type} \OperatorTok{=} \StringTok{\textquotesingle{}follows\textquotesingle{}} \KeywordTok{AND}\NormalTok{ weight }\OperatorTok{\textgreater{}} \FloatTok{0.5}\NormalTok{;      }\CommentTok{{-}{-} pruned by zonemaps}
\end{Highlighting}
\end{Shaded}

Put \textbf{cold / sparse} properties in a \texttt{props} blob. This is
the classic wide-vs-narrow trade: typed columns for what you query, a
blob for the long tail.

\section{Directed vs.~undirected}\label{directed-vs.-undirected}

Edges are stored \textbf{directed}. Model an undirected edge by
inserting both \texttt{(u,v)} and \texttt{(v,u)}, or let the
\href{snapshot.md}{CSR builder} symmetrize for undirected algorithms
(connected components, triangles do this internally). For fast
\textbf{backward} traversal, keep a second table keyed
\texttt{(dst,\ src)} --- the reverse adjacency.

\section{Transactions keep it
consistent}\label{transactions-keep-it-consistent}

Adding an edge that touches multiple tables (an edge row, a reverse-edge
row, a degree counter) belongs in \textbf{one transaction}, so a reader
never sees half an edge:

\begin{Shaded}
\begin{Highlighting}[]
\ControlFlowTok{BEGIN}\NormalTok{;}
  \KeywordTok{INSERT} \KeywordTok{INTO}\NormalTok{ edges }\KeywordTok{VALUES}\NormalTok{ (encode(}\DecValTok{1}\NormalTok{,}\DecValTok{2}\NormalTok{), }\DecValTok{1}\NormalTok{, }\DecValTok{2}\NormalTok{, }\StringTok{\textquotesingle{}follows\textquotesingle{}}\NormalTok{, }\FloatTok{1.0}\NormalTok{);}
  \KeywordTok{INSERT} \KeywordTok{INTO}\NormalTok{ edges\_rev }\KeywordTok{VALUES}\NormalTok{ (encode(}\DecValTok{2}\NormalTok{,}\DecValTok{1}\NormalTok{), }\DecValTok{1}\NormalTok{, }\DecValTok{2}\NormalTok{);}
\KeywordTok{COMMIT}\NormalTok{;}
\end{Highlighting}
\end{Shaded}

\section{\texorpdfstring{What's \emph{not}
modeled}{What's not modeled}}\label{whats-not-modeled}

\begin{itemize}
\tightlist
\item
  \textbf{Composite / multi-column keys} --- not yet native, hence the
  \texttt{(src,dst)} encoding (the roadmap's top item).
\item
  \textbf{Foreign keys} --- an explicit non-goal; referential integrity
  between nodes and edges is the application's responsibility.
\item
  \textbf{Node ids beyond 2³¹} --- the packed-key encoding's current
  ceiling.
\end{itemize}

The next chapters --- \href{model.md}{Clustered Adjacency} and
\href{snapshot.md}{The CSR Snapshot} --- show why this simple mapping
traverses fast.

\section{Clustered Adjacency}\label{clustered-adjacency}

The single hard constraint ChakraDB puts on a graph is also the source
of its adjacency speed. The engine has a \textbf{single-column primary
key} and \textbf{no secondary indexes}. So ``find all edges with
\texttt{src\ =\ X}'' is fast \emph{only if the edge table is physically
sorted by \texttt{src}.} And parts are sorted by the primary key.
Therefore:

\begin{quote}
\textbf{The edge key must sort src-major.} Encode \texttt{(src,\ dst)}
into one sortable key, and a node's out-edges become a contiguous key
range that part pruning answers cheaply.
\end{quote}

\section{The key encoding}\label{the-key-encoding}

The graph layer packs a directed edge into a single integer key:

\begin{verbatim}
key(src, dst) = (src << 32) | dst
\end{verbatim}

Because the high 32 bits are \texttt{src}, keys sort by \texttt{src}
first, then \texttt{dst}. All of node \texttt{X}'s out-edges occupy the
half-open key range \texttt{{[}X·2³²,\ (X+1)·2³²)}.

\begin{center}\includegraphics[width=0.92\linewidth,height=0.80\textheight,keepaspectratio]{diagrams/diagram-24.pdf}\end{center}

\texttt{out\_neighbors(1)} scans the key range
\texttt{{[}(1,0),\ (2,0))}, which is exactly the green rows ---
contiguous, and nothing else is touched.

\section{Why the scan is cheap: part
pruning}\label{why-the-scan-is-cheap-part-pruning}

Sealed parts are sorted by key, so each part's keys occupy a
\texttt{{[}min\_key,\ max\_key{]}} interval. The neighbor scan skips any
part whose interval cannot overlap the query range:

\begin{center}\includegraphics[width=0.92\linewidth,height=0.80\textheight,keepaspectratio]{diagrams/diagram-25.pdf}\end{center}

Only Part B can hold \texttt{src\ =\ 500}, so only Part B is scanned.
This is the same \href{../engine/storage.md}{zonemap part pruning} that
accelerates a selective SQL \texttt{WHERE} --- reused as the graph
adjacency index. In the ClickBench benchmark, a key-range scan of this
shape stays effectively \textbf{O(1) as the table grows 100×} (Part IX).
For a graph, that means neighbor lookup cost tracks the node's degree,
not the graph's size.

The primitive in the engine is:

\begin{Shaded}
\begin{Highlighting}[]
\CommentTok{// src/table.rs — prunes parts by [min\_key, max\_key], scans survivors.}
\KeywordTok{pub} \KeywordTok{fn}\NormalTok{ scan\_key\_range(}\OperatorTok{\&}\KeywordTok{self}\OperatorTok{,}\NormalTok{ lo}\OperatorTok{:} \OperatorTok{\&}\NormalTok{Value}\OperatorTok{,}\NormalTok{ hi}\OperatorTok{:} \OperatorTok{\&}\NormalTok{Value}\OperatorTok{,}\NormalTok{ snap}\OperatorTok{:}\NormalTok{ Snapshot) }\OperatorTok{{-}\textgreater{}} \DataTypeTok{Vec}\OperatorTok{\textless{}}\NormalTok{Row}\OperatorTok{\textgreater{};}
\end{Highlighting}
\end{Shaded}

and the graph layer uses it:

\begin{Shaded}
\begin{Highlighting}[]
\KeywordTok{pub} \KeywordTok{fn}\NormalTok{ out\_neighbors(}\OperatorTok{\&}\KeywordTok{self}\OperatorTok{,}\NormalTok{ node}\OperatorTok{:}\NormalTok{ NodeId) }\OperatorTok{{-}\textgreater{}} \DataTypeTok{Result}\OperatorTok{\textless{}}\DataTypeTok{Vec}\OperatorTok{\textless{}}\NormalTok{NodeId}\OperatorTok{\textgreater{}\textgreater{}} \OperatorTok{\{}
    \KeywordTok{let}\NormalTok{ lo }\OperatorTok{=}\NormalTok{ encode(node}\OperatorTok{,} \DecValTok{0}\NormalTok{)}\OperatorTok{;}
    \KeywordTok{let}\NormalTok{ hi }\OperatorTok{=}\NormalTok{ encode(node }\OperatorTok{+} \DecValTok{1}\OperatorTok{,} \DecValTok{0}\NormalTok{)}\OperatorTok{;}        \CommentTok{// exclusive}
    \CommentTok{// scan\_key\_range prunes to the parts that hold this src, over a snapshot.}
\OperatorTok{\}}
\end{Highlighting}
\end{Shaded}

\section{Directed, and the reverse
direction}\label{directed-and-the-reverse-direction}

Today the layer stores \textbf{directed} out-edges in one table. Two
traversal patterns follow:

\begin{itemize}
\tightlist
\item
  \textbf{Forward} (out-neighbors): the range scan above.
\item
  \textbf{Backward} (in-neighbors): the design keeps a second table
  keyed \texttt{(dst,\ src)}, so in-neighbors are a range scan on
  \emph{that} table. (Undirected algorithms --- like connected
  components and triangle counting --- symmetrize the adjacency in the
  CSR; see \href{snapshot.md}{The CSR Snapshot}.)
\end{itemize}

\section{The limit, and the fix}\label{the-limit-and-the-fix}

The packed-integer encoding caps node ids at \texttt{\textless{}\ 2³¹}
(so keys stay positive \texttt{i64}). That is a real v1 limitation. The
clean removal is a \textbf{native composite / bytes key} --- letting a
table declare \texttt{PRIMARY\ KEY\ (src,\ dst)} or a lexicographic
\texttt{BYTES} key. That would make edges first-class, remove the
encoding entirely, and benefit non-graph multi-column keys too. It is
the highest-leverage item on the graph roadmap (see the
\href{../../graph-exploration.md}{design exploration}).

\chapter{The CSR Snapshot}\label{the-csr-snapshot}

\epigraph{Simplicity is prerequisite for reliability.}{--- Edsger W. Dijkstra}

Whole-graph algorithms --- BFS, PageRank, connected components --- do
not want to issue a range scan per node. They want the
\textbf{Compressed Sparse Row (CSR)} representation: a flat neighbor
array plus per-node offsets, giving cache-friendly \texttt{O(1)}
neighbor iteration. \texttt{Graph::view()} builds one from a single MVCC
snapshot.

\section{What CSR is}\label{what-csr-is}

\begin{center}\includegraphics[width=0.92\linewidth,height=0.80\textheight,keepaspectratio]{diagrams/diagram-26.pdf}\end{center}

To iterate node \texttt{u}'s out-neighbors, take
\texttt{neighbors{[}offsets{[}u{]}\ ..\ offsets{[}u+1{]}{]}}. For node 2
above: \texttt{offsets{[}2{]}=3,\ offsets{[}3{]}=5} →
\texttt{neighbors{[}3..5{]}\ =\ {[}0,\ 3{]}}.

\section{Why building it is a linear
scan}\label{why-building-it-is-a-linear-scan}

The edge table is already sorted by key = \texttt{(src,\ dst)}
src-major. So a single scan yields \texttt{(src,\ dst)} pairs
\emph{already grouped by source}. Building CSR is then a counting pass
and a fill pass --- no sort, no hash join:

\begin{center}\includegraphics[width=0.92\linewidth,height=0.80\textheight,keepaspectratio]{diagrams/diagram-27.pdf}\end{center}

The actual build (\texttt{src/graph.rs}):

\begin{Shaded}
\begin{Highlighting}[]
\CommentTok{// Dense node numbering over every id that appears.}
\KeywordTok{let} \KeywordTok{mut}\NormalTok{ ids}\OperatorTok{:} \DataTypeTok{Vec}\OperatorTok{\textless{}}\NormalTok{NodeId}\OperatorTok{\textgreater{}} \OperatorTok{=}\NormalTok{ pairs}\OperatorTok{.}\NormalTok{iter()}\OperatorTok{.}\NormalTok{flat\_map(}\OperatorTok{|\&}\NormalTok{(s}\OperatorTok{,}\NormalTok{d)}\OperatorTok{|}\NormalTok{ [s}\OperatorTok{,}\NormalTok{d])}\OperatorTok{.}\NormalTok{collect()}\OperatorTok{;}
\NormalTok{ids}\OperatorTok{.}\NormalTok{sort\_unstable()}\OperatorTok{;}\NormalTok{ ids}\OperatorTok{.}\NormalTok{dedup()}\OperatorTok{;}
\KeywordTok{let}\NormalTok{ index}\OperatorTok{:}\NormalTok{ HashMap}\OperatorTok{\textless{}}\NormalTok{NodeId}\OperatorTok{,}\DataTypeTok{u32}\OperatorTok{\textgreater{}} \OperatorTok{=} \CommentTok{/* id {-}\textgreater{} dense idx */}\OperatorTok{;}

\CommentTok{// CSR by counting sort on the source\textquotesingle{}s dense index.}
\KeywordTok{let} \KeywordTok{mut}\NormalTok{ offsets }\OperatorTok{=} \PreprocessorTok{vec!}\NormalTok{[}\DecValTok{0u32}\OperatorTok{;}\NormalTok{ n }\OperatorTok{+} \DecValTok{1}\NormalTok{]}\OperatorTok{;}
\ControlFlowTok{for} \OperatorTok{\&}\NormalTok{(s}\OperatorTok{,}\NormalTok{\_) }\KeywordTok{in} \OperatorTok{\&}\NormalTok{pairs }\OperatorTok{\{}\NormalTok{ offsets[index[}\OperatorTok{\&}\NormalTok{s] }\KeywordTok{as} \DataTypeTok{usize} \OperatorTok{+} \DecValTok{1}\NormalTok{] }\OperatorTok{+=} \DecValTok{1}\OperatorTok{;} \OperatorTok{\}}
\ControlFlowTok{for}\NormalTok{ i }\KeywordTok{in} \DecValTok{0}\OperatorTok{..}\NormalTok{n }\OperatorTok{\{}\NormalTok{ offsets[i}\OperatorTok{+}\DecValTok{1}\NormalTok{] }\OperatorTok{+=}\NormalTok{ offsets[i]}\OperatorTok{;} \OperatorTok{\}}         \CommentTok{// prefix sum}

\KeywordTok{let} \KeywordTok{mut}\NormalTok{ adj }\OperatorTok{=} \PreprocessorTok{vec!}\NormalTok{[}\DecValTok{0u32}\OperatorTok{;}\NormalTok{ pairs}\OperatorTok{.}\NormalTok{len()]}\OperatorTok{;}
\KeywordTok{let} \KeywordTok{mut}\NormalTok{ cursor }\OperatorTok{=}\NormalTok{ offsets}\OperatorTok{.}\NormalTok{clone()}\OperatorTok{;}
\ControlFlowTok{for} \OperatorTok{\&}\NormalTok{(s}\OperatorTok{,}\NormalTok{d) }\KeywordTok{in} \OperatorTok{\&}\NormalTok{pairs }\OperatorTok{\{}
    \KeywordTok{let}\NormalTok{ si }\OperatorTok{=}\NormalTok{ index[}\OperatorTok{\&}\NormalTok{s] }\KeywordTok{as} \DataTypeTok{usize}\OperatorTok{;}
\NormalTok{    adj[cursor[si] }\KeywordTok{as} \DataTypeTok{usize}\NormalTok{] }\OperatorTok{=}\NormalTok{ index[}\OperatorTok{\&}\NormalTok{d]}\OperatorTok{;}
\NormalTok{    cursor[si] }\OperatorTok{+=} \DecValTok{1}\OperatorTok{;}
\OperatorTok{\}}
\end{Highlighting}
\end{Shaded}

Cost is \texttt{O(V\ +\ E)} plus the densify sort. The result is a
compact \texttt{\{ids,\ index,\ offsets,\ adj\}} --- the
\texttt{GraphView}.

\section{Consistency: it is a
snapshot}\label{consistency-it-is-a-snapshot}

\texttt{view()} pins an MVCC snapshot for the duration of the scan, so
the CSR is built from \textbf{one consistent instant} of the graph.
Writers keep committing edges; the \texttt{GraphView} is an immutable
copy and does not change. That is what the test asserts:

\begin{Shaded}
\begin{Highlighting}[]
\KeywordTok{let}\NormalTok{ v }\OperatorTok{=}\NormalTok{ g}\OperatorTok{.}\NormalTok{view()}\OperatorTok{?;}
\KeywordTok{let}\NormalTok{ before }\OperatorTok{=}\NormalTok{ v}\OperatorTok{.}\NormalTok{edge\_count()}\OperatorTok{;}
\NormalTok{g}\OperatorTok{.}\NormalTok{add\_edges([(}\DecValTok{6}\OperatorTok{,}\DecValTok{7}\OperatorTok{,}\DecValTok{1.0}\NormalTok{)}\OperatorTok{,}\NormalTok{(}\DecValTok{7}\OperatorTok{,}\DecValTok{8}\OperatorTok{,}\DecValTok{1.0}\NormalTok{)])}\OperatorTok{?;}    \CommentTok{// writes after the view was taken}
\PreprocessorTok{assert\_eq!}\NormalTok{(v}\OperatorTok{.}\NormalTok{edge\_count()}\OperatorTok{,}\NormalTok{ before)}\OperatorTok{;}      \CommentTok{// the snapshot view is unchanged}
\PreprocessorTok{assert\_eq!}\NormalTok{(g}\OperatorTok{.}\NormalTok{view()}\OperatorTok{?.}\NormalTok{edge\_count()}\OperatorTok{,}\NormalTok{ before }\OperatorTok{+} \DecValTok{2}\NormalTok{)}\OperatorTok{;} \CommentTok{// a fresh view sees them}
\end{Highlighting}
\end{Shaded}

\section{Memory: the honest ceiling}\label{memory-the-honest-ceiling}

The CSR lives in RAM. \texttt{offsets} is \texttt{4·(V+1)} bytes and
\texttt{neighbors} is \texttt{4·E} bytes --- about \textbf{4 bytes per
edge} plus the id maps. A 100-million-edge graph is \textasciitilde0.4
GB of CSR; a billion edges is a few GB. Graph algorithms are
memory-bound, and this is the same ceiling as ChakraDB's
\href{../engine/overview.md}{resident index}: ``fits in a big server's
RAM'' is in scope, trillion-edge graphs are not. For graphs larger than
memory, restrict the view to a subgraph (a k-hop neighborhood, a
component) before running the algorithm.

\section{Live Graph Analytics}\label{live-graph-analytics}

This is the chapter that says \emph{why ChakraDB's graph layer is
different}, not just convenient. The difference is one word:
\textbf{live.}

\section{The problem with the
alternatives}\label{the-problem-with-the-alternatives}

To run a graph algorithm you need a consistent view of the graph. The
usual ways to get one all give something up:

\begin{itemize}
\tightlist
\item
  \textbf{Load into NetworkX / igraph.} You get a \emph{dead static
  copy}. The moment you loaded it, it began to go stale. Re-analyzing
  the latest graph means reloading the whole thing.
\item
  \textbf{A lock-based graph database (e.g.~Neo4j).} A long analytical
  traversal contends with the writers mutating the graph; you either
  block ingest or read an inconsistent, partially-updated graph.
\item
  \textbf{A separate OLAP copy (ETL the graph into a warehouse).} Now
  the analysis is minutes-to-hours stale, and you run two systems.
\end{itemize}

\section{What ChakraDB does instead}\label{what-chakradb-does-instead}

\begin{Shaded}
\begin{Highlighting}[]
\NormalTok{sequenceDiagram}
\NormalTok{    participant Writers}
\NormalTok{    participant Clock as MVCC clock}
\NormalTok{    participant Alg as Algorithm (GraphView)}
\NormalTok{    Writers{-}\textgreater{}\textgreater{}Clock: add edge @ CSN 100}
\NormalTok{    Alg{-}\textgreater{}\textgreater{}Clock: view() pins snapshot S = 100}
\NormalTok{    Note over Alg: builds CSR from S; runs PageRank...}
\NormalTok{    Writers{-}\textgreater{}\textgreater{}Clock: add edge @ CSN 101  (never blocked)}
\NormalTok{    Writers{-}\textgreater{}\textgreater{}Clock: add edge @ CSN 102}
\NormalTok{    Note over Alg: still computing over S = 100 — a consistent graph}
\NormalTok{    Alg{-}{-}\textgreater{}\textgreater{}Alg: result reflects exactly the graph at S = 100}
\end{Highlighting}
\end{Shaded}

An algorithm builds its CSR from \textbf{one MVCC snapshot} and runs to
completion over that consistent instant. Meanwhile writers keep
committing edges --- \textbf{never blocked}, because readers and writers
share nothing but the append-only clock. The \href{../engine/mvcc.md}{GC
watermark} holds the snapshot's versions alive for the algorithm's
duration even as newer data is compacted.

The result: you can run PageRank (or components, or a fraud traversal)
\textbf{every minute over the latest consistent graph, while ingest
never pauses} --- in one embedded process. That is the combination the
alternatives cannot offer.

\section{The pattern in code}\label{the-pattern-in-code}

\begin{Shaded}
\begin{Highlighting}[]
\CommentTok{// A background loop: recompute influence over the live graph, on a cadence.}
\ControlFlowTok{loop} \OperatorTok{\{}
    \KeywordTok{let}\NormalTok{ view }\OperatorTok{=}\NormalTok{ graph}\OperatorTok{.}\NormalTok{view()}\OperatorTok{?;}              \CommentTok{// consistent snapshot; ingest continues}
    \KeywordTok{let}\NormalTok{ ranks }\OperatorTok{=}\NormalTok{ view}\OperatorTok{.}\NormalTok{pagerank(}\DecValTok{20}\OperatorTok{,} \DecValTok{0.85}\NormalTok{)}\OperatorTok{;}
\NormalTok{    publish\_top\_influencers(}\OperatorTok{\&}\NormalTok{ranks)}\OperatorTok{;}       \CommentTok{// serve results}
\NormalTok{    sleep(}\PreprocessorTok{Duration::}\NormalTok{from\_secs(}\DecValTok{60}\NormalTok{))}\OperatorTok{;}
\OperatorTok{\}}
\end{Highlighting}
\end{Shaded}

Each \texttt{view()} is a fresh consistent snapshot. The writers feeding
\texttt{graph} never notice the analytics running.

\section{Why it composes with the rest of the
database}\label{why-it-composes-with-the-rest-of-the-database}

Because the graph \emph{is} tables, graph results compose with SQL and
transactions over the \textbf{same} live data:

\begin{itemize}
\tightlist
\item
  Score an entity with a graph traversal \textbf{and} join it to its
  transactional row in one consistent snapshot.
\item
  Ingest an event (transactional), update the graph edge
  (transactional), and let the next analytics pass pick it up --- no
  cross-system consistency to reason about.
\end{itemize}

This is HTAP extended to graphs: \textbf{T + A + G over one snapshot
clock.} The \href{../case-studies/fraud.md}{fraud case study} works it
end to end.

\begin{quote}
\textbf{The honest scope.} v1 \emph{recomputes} an algorithm over a
snapshot on demand; it does not yet \emph{maintain} live results
incrementally as edges change. Incremental graph algorithms (streaming
PageRank, incremental components) are a roadmap item --- a v2 that leans
even harder on this live-mutation story.
\end{quote}

\chapter{Graph Algorithms}\label{graph-algorithms}

\epigraph{A journey of a thousand miles begins with a single step.}{--- Lao Tzu}

All of ChakraDB's graph algorithms run over a \texttt{GraphView} --- the
immutable \href{snapshot.md}{CSR snapshot}. Because the view is a
consistent copy taken from one MVCC snapshot, every algorithm below runs
\textbf{concurrently with ingest}: writers keep adding edges while the
algorithm computes over a stable graph. Each is a compact Rust kernel;
this chapter gives the method, the pseudocode, and the complexity.

Throughout, \texttt{V} = number of nodes, \texttt{E} = number of
directed edges, \texttt{d(u)} = out -degree of \texttt{u}.

\section{Breadth-first search \& shortest
path}\label{breadth-first-search-shortest-path}

Unweighted shortest paths --- ``degrees of separation,'' reachability,
k-hop neighborhoods --- are BFS over the out-edges.

\begin{Shaded}
\begin{Highlighting}[]
\NormalTok{flowchart LR}
\NormalTok{    S((start)):::f {-}{-}\textgreater{} A:::f1 {-}{-}\textgreater{} B:::f2}
\NormalTok{    S {-}{-}\textgreater{} C:::f1}
\NormalTok{    C {-}{-}\textgreater{} B}
\NormalTok{    C {-}{-}\textgreater{} D:::f2}
\NormalTok{    classDef f fill:\#8ecae6; classDef f1 fill:\#bde0fe; classDef f2 fill:\#e7f0ff;}
\end{Highlighting}
\end{Shaded}

\begin{Shaded}
\begin{Highlighting}[]
\NormalTok{bfs(start):}
\NormalTok{  depth[*] = ∞;  depth[start] = 0;  queue = [start]}
\NormalTok{  while queue not empty:}
\NormalTok{    u = queue.pop\_front()}
\NormalTok{    for v in neighbors(u):}
\NormalTok{      if depth[v] == ∞:}
\NormalTok{        depth[v] = depth[u] + 1}
\NormalTok{        queue.push\_back(v)}
\NormalTok{  return depth}
\end{Highlighting}
\end{Shaded}

\texttt{shortest\_path(a,\ b)} is the same wave with a
\texttt{parent{[}{]}} array; on reaching \texttt{b}, it walks parents
back to \texttt{a}. \textbf{Complexity \texttt{O(V\ +\ E)}}, one queue,
one visited array.

\begin{Shaded}
\begin{Highlighting}[]
\PreprocessorTok{assert\_eq!}\NormalTok{(view}\OperatorTok{.}\NormalTok{shortest\_path(}\DecValTok{1}\OperatorTok{,} \DecValTok{4}\NormalTok{)}\OperatorTok{,} \ConstantTok{Some}\NormalTok{(}\PreprocessorTok{vec!}\NormalTok{[}\DecValTok{1}\OperatorTok{,} \DecValTok{3}\OperatorTok{,} \DecValTok{4}\NormalTok{]))}\OperatorTok{;}
\end{Highlighting}
\end{Shaded}

\section{PageRank}\label{pagerank}

The influence/importance score: a node is important if important nodes
point at it. ChakraDB computes it by \textbf{power iteration},
scattering each node's rank across its out-edges.

\begin{Shaded}
\begin{Highlighting}[]
\NormalTok{pagerank(iterations, d):            \# d = damping, \textasciitilde{}0.85}
\NormalTok{  rank[*] = 1/V}
\NormalTok{  repeat \textasciigrave{}iterations\textasciigrave{} times:}
\NormalTok{    dangling = Σ rank[u] for u with d(u) == 0}
\NormalTok{    base = (1{-}d)/V + d·dangling/V     \# teleport + redistributed dead{-}ends}
\NormalTok{    next[*] = base}
\NormalTok{    for u in 0..V:}
\NormalTok{      if d(u) \textgreater{} 0:}
\NormalTok{        share = d · rank[u] / d(u)}
\NormalTok{        for v in neighbors(u): next[v] += share}
\NormalTok{    rank = next}
\NormalTok{  return rank                         \# sums to \textasciitilde{}1}
\end{Highlighting}
\end{Shaded}

The \textbf{dangling-node} handling matters: nodes with no out-edges
would leak probability mass; their rank is collected and redistributed
uniformly through \texttt{base}, so the scores stay a proper
distribution. \textbf{Complexity \texttt{O(iterations\ ·\ E)}.}

\begin{Shaded}
\begin{Highlighting}[]
\CommentTok{// A star where 2,3,4,5 all point at 1 → node 1 ranks highest.}
\KeywordTok{let}\NormalTok{ pr }\OperatorTok{=}\NormalTok{ view}\OperatorTok{.}\NormalTok{pagerank(}\DecValTok{30}\OperatorTok{,} \DecValTok{0.85}\NormalTok{)}\OperatorTok{;}
\end{Highlighting}
\end{Shaded}

\emph{Personalized PageRank} (recommendation, ``similar to X'') is the
same iteration with the teleport vector concentrated on a seed set
instead of uniform --- on the roadmap.

\section{Connected components}\label{connected-components}

``Which nodes belong to the same island?'' ChakraDB computes
\textbf{weakly-connected components} (edges treated as undirected) with
a \textbf{union-find} (disjoint-set) structure using path-halving.

\begin{Shaded}
\begin{Highlighting}[]
\NormalTok{connected\_components():}
\NormalTok{  parent[u] = u for all u}
\NormalTok{  for each edge (u, v):}
\NormalTok{    union(u, v)          \# merge the two sets}
\NormalTok{  relabel roots to 0..k}
\NormalTok{  return node {-}\textgreater{} component id}
\end{Highlighting}
\end{Shaded}

Union-find with path compression runs in near-linear
\textbf{\texttt{O(E\ ·\ α(V))}} time, where \texttt{α} is the
inverse-Ackermann function (effectively constant). Strongly-connected
components (Tarjan) are a roadmap addition.

\begin{Shaded}
\begin{Highlighting}[]
\KeywordTok{let}\NormalTok{ c }\OperatorTok{=}\NormalTok{ view}\OperatorTok{.}\NormalTok{connected\_components()}\OperatorTok{;}
\PreprocessorTok{assert\_eq!}\NormalTok{(c[}\OperatorTok{\&}\DecValTok{1}\NormalTok{]}\OperatorTok{,}\NormalTok{ c[}\OperatorTok{\&}\DecValTok{4}\NormalTok{])}\OperatorTok{;}   \CommentTok{// same island}
\PreprocessorTok{assert\_ne!}\NormalTok{(c[}\OperatorTok{\&}\DecValTok{1}\NormalTok{]}\OperatorTok{,}\NormalTok{ c[}\OperatorTok{\&}\DecValTok{5}\NormalTok{])}\OperatorTok{;}   \CommentTok{// different island}
\end{Highlighting}
\end{Shaded}

\section{Triangle counting}\label{triangle-counting}

Triangles measure local density --- community strength, clustering
coefficient, spam detection. ChakraDB counts them on the
\textbf{undirected} graph by sorted-neighbor intersection with an
ordering trick that counts each triangle exactly once.

\begin{Shaded}
\begin{Highlighting}[]
\NormalTok{triangle\_count():}
\NormalTok{  build sorted, de{-}duplicated undirected adjacency}
\NormalTok{  total = 0}
\NormalTok{  for u in 0..V:}
\NormalTok{    for v in neighbors(u) with v \textgreater{} u:}
\NormalTok{      total += | \{ w in adj(u) ∩ adj(v) : w \textgreater{} v \} |   \# sorted two{-}pointer}
\NormalTok{  return total}
\end{Highlighting}
\end{Shaded}

Requiring \texttt{u\ \textless{}\ v\ \textless{}\ w} means the triangle
\texttt{\{u,v,w\}} is counted once, not six times. \textbf{Complexity
\texttt{O(E\^{}\{1.5\})}} in the worst case (the classic node-iterator
bound).

\begin{Shaded}
\begin{Highlighting}[]
\CommentTok{// Triangle 1{-}2{-}3, plus a dangling edge 3{-}4:}
\PreprocessorTok{assert\_eq!}\NormalTok{(view}\OperatorTok{.}\NormalTok{triangle\_count()}\OperatorTok{,} \DecValTok{1}\NormalTok{)}\OperatorTok{;}
\end{Highlighting}
\end{Shaded}

\section{Degree \& neighborhoods}\label{degree-neighborhoods}

The cheapest signals need no CSR at all:

\begin{itemize}
\tightlist
\item
  \texttt{out\_degree(node)} --- the length of the node's neighbor
  range.
\item
  \texttt{out\_neighbors(node)} --- a single \href{model.md}{pruned
  range scan}, live (no snapshot copy needed).
\item
  k-hop neighborhood --- BFS truncated at depth k.
\end{itemize}

\section{Systemic risk: the Eisenberg--Noe clearing
vector}\label{systemic-risk-the-eisenbergnoe-clearing-vector}

Reading each directed edge \texttt{i\ →\ j} of weight \texttt{w} as a
\emph{liability} --- \texttt{i} owes \texttt{j} the amount \texttt{w}
--- \texttt{eisenberg\_noe(external\_assets)} computes the clearing
payment vector that models \textbf{default contagion}: who can pay, who
defaults, and who a cascade drags under. It is the canonical
systemic-risk model, solved by a monotone Picard iteration to the
greatest clearing vector, and it powers the
\href{../case-studies/ccr.md}{Counterparty Risk case study}. Full
derivation there.

\section{Link prediction \&
recommendations}\label{link-prediction-recommendations}

Two primitives turn the graph into a recommender
(\href{../case-studies/reco.md}{Recommendations case study}):

\begin{itemize}
\tightlist
\item
  \textbf{\texttt{adamic\_adar(a,\ b)}} --- a link-prediction score
  summing \texttt{1/ln\ deg(z)} over shared neighbours \texttt{z}: rare,
  niche connections count for more than popular ones. On a user↔item
  graph it is ``you might also like.''
\item
  \textbf{\texttt{recommend(seed,\ k)}} --- personalized PageRank seeded
  at \texttt{seed}, with the seed's existing neighbours and zero-signal
  candidates removed, returning the top-\texttt{k} remainder.
  Collaborative filtering as a random-walk-with-restart, in one call.
\end{itemize}

\section{The algorithm menu}\label{the-algorithm-menu}

Every algorithm below is built into the core (\texttt{src/graph.rs}) and
available from Rust and Python, over a consistent
\href{snapshot.md}{\texttt{GraphView}} snapshot.

\begin{longtable}[]{@{}
  >{\raggedright\arraybackslash}p{(\linewidth - 2\tabcolsep) * \real{0.5000}}
  >{\raggedright\arraybackslash}p{(\linewidth - 2\tabcolsep) * \real{0.5000}}@{}}
\toprule\noalign{}
\begin{minipage}[b]{\linewidth}\raggedright
Family
\end{minipage} & \begin{minipage}[b]{\linewidth}\raggedright
Algorithms
\end{minipage} \\
\midrule\noalign{}
\endhead
\bottomrule\noalign{}
\endlastfoot
\textbf{Traversal \& paths} & \texttt{bfs}, \texttt{shortest\_path},
\texttt{dijkstra}, \texttt{weighted\_shortest\_path},
\texttt{topological\_order} \\
\textbf{Centrality} & \texttt{pagerank},
\texttt{personalized\_pagerank}, \texttt{degree\_centrality},
\texttt{closeness\_centrality}, \texttt{betweenness\_centrality}
(Brandes) \\
\textbf{Community \& structure} & \texttt{connected\_components} (WCC),
\texttt{strongly\_connected\_components} / \texttt{laundering\_cycles},
\texttt{label\_propagation}, \texttt{k\_core},
\texttt{triangle\_count} \\
\textbf{Similarity \& link prediction} & \texttt{common\_neighbors},
\texttt{jaccard\_similarity}, \texttt{adamic\_adar},
\texttt{recommend} \\
\textbf{Systemic risk} & \texttt{eisenberg\_noe} (clearing vector /
default cascade) \\
\textbf{Degrees} & \texttt{in\_degree}, \texttt{out\_degree},
\texttt{in\_neighbors}, \texttt{out\_neighbors} \\
\end{longtable}

Deliberately \textbf{not} yet in core (expensive or specialized ---
candidates for a later release): Louvain/Leiden modularity, exact
diameter, general subgraph isomorphism, Johnson's all-simple-cycles,
temporal motifs, and max-flow/min-cut. See the
\href{../../graph-exploration.md}{design exploration} for the roadmap.

\part{Reactive}

\chapter{Change Streams \& Materialized
Workers}\label{change-streams-materialized-workers}

\epigraph{You can never step into the same river twice, for new waters are always flowing on to you.}{--- Heraclitus}

A traditional database is passive: you ask a question, it answers. The
systems this book cares about --- real-time AML, live risk --- are
\emph{reactive}: the instant a transaction commits, something must
happen. ChakraDB closes that loop inside the engine, on one consistent
live dataset, without ever slowing the writer. This chapter is the
machinery: a committed-change stream, workers that maintain derived
state from it, and the transports that carry it across processes.

\section{Why not in-SQL triggers}\label{why-not-in-sql-triggers}

The obvious design --- \texttt{CREATE\ TRIGGER} running a stored
procedure inside the transaction --- is the wrong one for this workload.
A synchronous trigger runs user code (often Python, under the GIL) in
the \textbf{writer's critical section}, so it serializes the very ingest
that ChakraDB is built to keep fast. Worse, it couples the writer's
latency and failure to arbitrary application logic.

ChakraDB instead exposes the model an event-driven system actually
wants: a \textbf{post-commit change stream} (change-data-capture, CDC).
Changes are delivered \emph{after} they commit, on a separate thread ---
the writer never waits for a reader.

\section{The change stream}\label{the-change-stream}

A \texttt{CdcBackend} decorates the engine's backend and publishes a
\texttt{Change} for every committed \texttt{INSERT} / \texttt{UPDATE} /
\texttt{DELETE}, carrying the operation, the commit sequence number
(CSN), and the old and new row images:

\begin{Shaded}
\begin{Highlighting}[]
\KeywordTok{pub} \KeywordTok{struct}\NormalTok{ Change }\OperatorTok{\{}
    \KeywordTok{pub}\NormalTok{ table}\OperatorTok{:} \DataTypeTok{String}\OperatorTok{,}
    \KeywordTok{pub}\NormalTok{ op}\OperatorTok{:}\NormalTok{ ChangeOp}\OperatorTok{,}          \CommentTok{// Insert | Update | Delete}
    \KeywordTok{pub}\NormalTok{ csn}\OperatorTok{:}\NormalTok{ Csn}\OperatorTok{,}              \CommentTok{// total commit order}
    \KeywordTok{pub}\NormalTok{ columns}\OperatorTok{:}\NormalTok{ Arc}\OperatorTok{\textless{}}\DataTypeTok{Vec}\OperatorTok{\textless{}}\DataTypeTok{String}\OperatorTok{\textgreater{}\textgreater{},}
    \KeywordTok{pub}\NormalTok{ old}\OperatorTok{:} \DataTypeTok{Option}\OperatorTok{\textless{}}\DataTypeTok{Vec}\OperatorTok{\textless{}}\NormalTok{Value}\OperatorTok{\textgreater{}\textgreater{},}   \CommentTok{// pre{-}image; None for Insert}
    \KeywordTok{pub}\NormalTok{ new}\OperatorTok{:} \DataTypeTok{Option}\OperatorTok{\textless{}}\DataTypeTok{Vec}\OperatorTok{\textless{}}\NormalTok{Value}\OperatorTok{\textgreater{}\textgreater{},}   \CommentTok{// post{-}image; None for Delete}
\OperatorTok{\}}
\end{Highlighting}
\end{Shaded}

Because the wrap sits \emph{above} the write path, it adds nothing to
the engine's hot locks. An \texttt{INSERT} --- the ingest hot path ---
pays only a channel send; the old-row pre-image for
\texttt{UPDATE}/\texttt{DELETE} is read outside any core lock. This
reuses the two things ChakraDB already has and most engines lack: an
ordered commit log (the \href{../engine/durability.md}{WAL}, stamped
with CSN) and cheap non-blocking \href{../engine/mvcc.md}{snapshots}.

Wire it up and subscribe:

\begin{Shaded}
\begin{Highlighting}[]
\KeywordTok{let}\NormalTok{ cdc }\OperatorTok{=} \PreprocessorTok{Cdc::}\NormalTok{new()}\OperatorTok{;}
\KeywordTok{let}\NormalTok{ engine }\OperatorTok{=} \PreprocessorTok{SqlEngine::}\NormalTok{with\_backend(}\PreprocessorTok{CdcBackend::}\NormalTok{wrap(db}\OperatorTok{,}\NormalTok{ cdc}\OperatorTok{.}\NormalTok{clone()))}\OperatorTok{;}
\KeywordTok{let}\NormalTok{ stream }\OperatorTok{=}\NormalTok{ cdc}\OperatorTok{.}\NormalTok{subscribe(}\ConstantTok{Some}\NormalTok{(}\StringTok{"transactions"}\NormalTok{))}\OperatorTok{;}   \CommentTok{// one table, or None for all}

\ControlFlowTok{while} \KeywordTok{let} \ConstantTok{Some}\NormalTok{(batch) }\OperatorTok{=}\NormalTok{ stream}\OperatorTok{.}\NormalTok{recv() }\OperatorTok{\{}              \CommentTok{// one commit\textquotesingle{}s changes, in order}
    \ControlFlowTok{for}\NormalTok{ change }\KeywordTok{in} \OperatorTok{\&}\NormalTok{batch }\OperatorTok{\{}
        \ControlFlowTok{if}\NormalTok{ change}\OperatorTok{.}\NormalTok{op }\OperatorTok{==} \PreprocessorTok{ChangeOp::}\NormalTok{Insert }\OperatorTok{\{}
\NormalTok{            react(change}\OperatorTok{.}\NormalTok{new}\OperatorTok{.}\NormalTok{as\_ref()}\OperatorTok{.}\NormalTok{unwrap())}\OperatorTok{;}
        \OperatorTok{\}}
    \OperatorTok{\}}
\OperatorTok{\}}
\end{Highlighting}
\end{Shaded}

Delivery is \textbf{at-least-once} and \textbf{in commit order}; a
rolled-back transaction never appears. The stream is a \emph{pull}
primitive --- the consumer owns its thread, cadence, and backpressure.

\begin{quote}
\textbf{ALGORITHM --- publishing a committed change}

\begin{Shaded}
\begin{Highlighting}[]
\NormalTok{On a mutating call m (insert/upsert/update/delete) through CdcBackend:}
\NormalTok{1  if m is update/delete: old ← inner.get\_latest(key)      ▷ pre{-}image, no core lock}
\NormalTok{2  csn ← inner.apply(m)                                    ▷ the real write (+ WAL)}
\NormalTok{3  op  ← Insert | Update | Delete   (upsert: Update iff old present)}
\NormalTok{4  publish Change\{ table, op, csn, columns, old, new \}     ▷ after commit, off the write path}
\end{Highlighting}
\end{Shaded}
\end{quote}

\section{The Python hook}\label{the-python-hook}

In Python the stream is a callback, and the connection is CDC-wrapped by
default. \texttt{on\_change} returns a \texttt{Subscription} whose
lifecycle the caller owns:

\begin{Shaded}
\begin{Highlighting}[]
\KeywordTok{def}\NormalTok{ react(old, new):                     }\CommentTok{\# dicts (column → value), or None}
    \ControlFlowTok{if}\NormalTok{ new }\KeywordTok{and}\NormalTok{ new[}\StringTok{"amount"}\NormalTok{] }\OperatorTok{\textgreater{}=} \DecValTok{9000}\NormalTok{:}
\NormalTok{        alert(new[}\StringTok{"src"}\NormalTok{], new[}\StringTok{"dst"}\NormalTok{])}

\NormalTok{sub }\OperatorTok{=}\NormalTok{ conn.on\_change(}\StringTok{"transactions"}\NormalTok{, react)}
\NormalTok{...                                       }\CommentTok{\# worker runs on its own thread}
\NormalTok{sub.close()                              }\CommentTok{\# stop it — or use it as a context manager}
\end{Highlighting}
\end{Shaded}

The callback fires after commit, so it never blocks the writer --- the
whole engine, including the \href{../graph/algorithms.md}{graph
algorithms}, reachable from a hook in a few lines.

\section{Materialized workers}\label{materialized-workers}

Most reactions are not one-shot: they \emph{maintain state} --- a
running aggregate, a graph projection, a set of alerts. That is a
\textbf{materialized worker}: a named, incrementally-maintained
derivation of the data. You implement \emph{what} is maintained; the
runtime owns the loop, commit ordering, cursor tracking, and lifecycle.

\begin{Shaded}
\begin{Highlighting}[]
\KeywordTok{pub} \KeywordTok{trait}\NormalTok{ MaterializedWorker}\OperatorTok{:} \BuiltInTok{Send} \OperatorTok{+} \OtherTok{\textquotesingle{}static} \OperatorTok{\{}
    \KeywordTok{fn}\NormalTok{ apply(}\OperatorTok{\&}\KeywordTok{mut} \KeywordTok{self}\OperatorTok{,}\NormalTok{ change}\OperatorTok{:} \OperatorTok{\&}\NormalTok{Change)}\OperatorTok{;}      \CommentTok{// fold one change into state}
    \KeywordTok{fn}\NormalTok{ on\_commit(}\OperatorTok{\&}\KeywordTok{mut} \KeywordTok{self}\OperatorTok{,}\NormalTok{ \_csn}\OperatorTok{:}\NormalTok{ Csn) }\OperatorTok{\{\}}      \CommentTok{// periodic hook (heavier passes)}
\OperatorTok{\}}

\KeywordTok{let}\NormalTok{ m }\OperatorTok{=}\NormalTok{ cdc}\OperatorTok{.}\NormalTok{register(}\StringTok{"aml{-}detector"}\OperatorTok{,} \ConstantTok{Some}\NormalTok{(}\StringTok{"transactions"}\NormalTok{)}\OperatorTok{,} \PreprocessorTok{Worker::}\NormalTok{new())}\OperatorTok{;}
\NormalTok{m}\OperatorTok{.}\NormalTok{query(}\OperatorTok{|}\NormalTok{w}\OperatorTok{|}\NormalTok{ w}\OperatorTok{.}\NormalTok{current\_exposure}\OperatorTok{.}\NormalTok{clone())}\OperatorTok{;}       \CommentTok{// read the derived state any time}
\NormalTok{m}\OperatorTok{.}\NormalTok{cursor()}\OperatorTok{;}                                    \CommentTok{// how far it has consumed (a resume point)}
\NormalTok{m}\OperatorTok{.}\NormalTok{stop()}\OperatorTok{;}                                      \CommentTok{// client owns the lifecycle}
\end{Highlighting}
\end{Shaded}

This is the disciplined worker primitive --- \textbf{a function of the
data}, not an application host. A worker consumes changes, maintains
state, exposes it for query, and writes derived rows back; it does not
open sockets or run arbitrary services. That boundary is what keeps
ChakraDB a database and not an app server.

\subsection{The registry}\label{the-registry}

Named workers are tracked and observable --- the control surface for a
fleet:

\begin{Shaded}
\begin{Highlighting}[]
\ControlFlowTok{for}\NormalTok{ w }\KeywordTok{in}\NormalTok{ cdc}\OperatorTok{.}\NormalTok{workers() }\OperatorTok{\{}
    \CommentTok{// w.name, w.table, w.cursor (CSN), w.running}
\OperatorTok{\}}
\NormalTok{cdc}\OperatorTok{.}\NormalTok{stop\_worker(}\StringTok{"aml{-}detector"}\NormalTok{)}\OperatorTok{;}
\end{Highlighting}
\end{Shaded}

The CSN cursor is the key to durability: on restart, rebuild the
worker's state from a snapshot and resume the stream from the last
cursor --- no missed or double-counted change, a hard requirement for
compliance workloads.

\subsection{The three-tier pattern}\label{the-three-tier-pattern}

Firing per event rewards \emph{incremental} work. The case studies use a
tiered model so one embedded engine keeps up with millions of events per
hour:

\begin{longtable}[]{@{}
  >{\raggedright\arraybackslash}p{(\linewidth - 6\tabcolsep) * \real{0.2500}}
  >{\raggedright\arraybackslash}p{(\linewidth - 6\tabcolsep) * \real{0.2500}}
  >{\raggedright\arraybackslash}p{(\linewidth - 6\tabcolsep) * \real{0.2500}}
  >{\raggedright\arraybackslash}p{(\linewidth - 6\tabcolsep) * \real{0.2500}}@{}}
\toprule\noalign{}
\begin{minipage}[b]{\linewidth}\raggedright
Tier
\end{minipage} & \begin{minipage}[b]{\linewidth}\raggedright
Runs
\end{minipage} & \begin{minipage}[b]{\linewidth}\raggedright
Scope
\end{minipage} & \begin{minipage}[b]{\linewidth}\raggedright
Example
\end{minipage} \\
\midrule\noalign{}
\endhead
\bottomrule\noalign{}
\endlastfoot
\textbf{T0} & every committed change & O(1)--O(degree) & fan-in counter,
limit check \\
\textbf{T1} & micro-batch (seconds) & touched subgraph & bounded cycle
search \\
\textbf{T2} & periodic (minutes) & whole graph, on a \texttt{view()} &
PageRank, Eisenberg--Noe, VaR \\
\end{longtable}

T0 keeps up with the firehose; T2 is heavy but rare and runs on an MVCC
snapshot, so it never blocks ingest. That non-blocking split is the
difference between real-time reaction and an overnight batch.

\section{Transports: in-process, files, and
Kafka}\label{transports-in-process-files-and-kafka}

A \texttt{ChangeSink} carries the stream beyond the local thread:

\begin{Shaded}
\begin{Highlighting}[]
\KeywordTok{pub} \KeywordTok{trait}\NormalTok{ ChangeSink}\OperatorTok{:} \BuiltInTok{Send} \OperatorTok{+} \BuiltInTok{Sync} \OperatorTok{\{} \KeywordTok{fn}\NormalTok{ emit(}\OperatorTok{\&}\KeywordTok{self}\OperatorTok{,}\NormalTok{ batch}\OperatorTok{:} \OperatorTok{\&}\NormalTok{[Change])}\OperatorTok{;} \OperatorTok{\}}
\end{Highlighting}
\end{Shaded}

\begin{itemize}
\tightlist
\item
  \textbf{In-process channel} (default): lowest latency, single node ---
  a worker in the same process, as the streaming case studies run.
\item
  \textbf{\texttt{JsonlSink}}: appends each change as one JSON line to a
  file --- a broker-less \textbf{cross-process} change log. A separate
  worker process tails the file and consumes the stream, with no shared
  memory and no database lock.
\item
  \textbf{Kafka / Redpanda / NATS}: implement \texttt{emit} to produce
  each \texttt{change.to\_json()} keyed by a partition key (e.g.~account
  id). \emph{N} workers then consume partitions in parallel for
  horizontal scale-out --- the same worker code, a different transport.
\end{itemize}

\begin{Shaded}
\begin{Highlighting}[]
\NormalTok{flowchart LR}
\NormalTok{    W["committed writes"] {-}{-}\textgreater{} CDC["change stream (CSN{-}ordered)"]}
\NormalTok{    CDC {-}{-}\textgreater{} S\{"ChangeSink"\}}
\NormalTok{    S {-}{-}\textgreater{}|in{-}process| A["worker (same process)"]}
\NormalTok{    S {-}{-}\textgreater{}|JsonlSink| B["worker (separate process, tails the log)"]}
\NormalTok{    S {-}{-}\textgreater{}|Kafka: key=hash(entity)| C["workers ×N (sharded)"]}
\end{Highlighting}
\end{Shaded}

The scaling story is one line: \textbf{resident workers for single-node
hot state; Kafka-fanned workers for scale-out --- the same derivation
either way.}

\section{Where it leads}\label{where-it-leads}

Two complete systems are built on this machinery: the
\href{../case-studies/aml.md}{Real-Time AML} and
\href{../case-studies/ccr.md}{Counterparty Credit \& Market Risk} case
studies. Each generates its own synthetic feed, reacts to every
committed row through a registered materialized worker, and sustains
150+ million events per hour on a single node --- detection running
concurrently with ingest, because in ChakraDB the reader never blocks
the writer.

\part{Getting Started}

\chapter{Installation \& Building}\label{installation-building}

\epigraph{The beginning is the most important part of the work.}{--- Plato, \emph{The Republic}}

ChakraDB is an embedded database: there is no server to install, no
daemon to run. You add it to a Rust project as a crate, or install the
Python package and \texttt{import\ chakradb}. This chapter gets you from
a clean checkout to a running build.

\section{Prerequisites}\label{prerequisites}

\begin{itemize}
\tightlist
\item
  \textbf{Rust} 1.75 or newer (\texttt{rustup} recommended).
\item
  A C toolchain (for a handful of transitive native dependencies).
\item
  \textbf{Python} 3.9+ and
  \href{https://www.maturin.rs/}{\texttt{maturin}} \emph{only} if you
  want the Python bindings.
\end{itemize}

\section{Building the core (Rust)}\label{building-the-core-rust}

The core builds in two flavours, selected by Cargo features:

\begin{longtable}[]{@{}
  >{\raggedright\arraybackslash}p{(\linewidth - 4\tabcolsep) * \real{0.2333}}
  >{\raggedright\arraybackslash}p{(\linewidth - 4\tabcolsep) * \real{0.3000}}
  >{\raggedright\arraybackslash}p{(\linewidth - 4\tabcolsep) * \real{0.4667}}@{}}
\toprule\noalign{}
\begin{minipage}[b]{\linewidth}\raggedright
Build
\end{minipage} & \begin{minipage}[b]{\linewidth}\raggedright
Command
\end{minipage} & \begin{minipage}[b]{\linewidth}\raggedright
What you get
\end{minipage} \\
\midrule\noalign{}
\endhead
\bottomrule\noalign{}
\endlastfoot
Interpreter only & \texttt{cargo\ build\ -\/-no-default-features} & The
full storage engine, SQL interpreter, MVCC, and the graph library --- no
analytical vectorized engine. Small and fast to compile. \\
With DataFusion & \texttt{cargo\ build\ -\/-features\ datafusion} & Adds
the vectorized analytical engine and the HTAP query router. \\
\end{longtable}

Run the test suite the same way:

\begin{Shaded}
\begin{Highlighting}[]
\ExtensionTok{cargo}\NormalTok{ test }\AttributeTok{{-}{-}no{-}default{-}features}          \CommentTok{\# core + graph}
\ExtensionTok{cargo}\NormalTok{ test }\AttributeTok{{-}{-}features}\NormalTok{ datafusion          }\CommentTok{\# + analytical path}
\end{Highlighting}
\end{Shaded}

\begin{quote}
\textbf{Why two profiles?} The interpreter path is what powers the
transactional workload and the graph engine; it has no heavy
dependencies and compiles in seconds. DataFusion is bought in only when
you want the analytical half of HTAP. Every algorithm in
\href{../engine/overview.md}{Part IV --- Graph} is available in
\emph{both} profiles.
\end{quote}

\section{Adding ChakraDB to a Rust
project}\label{adding-chakradb-to-a-rust-project}

\begin{Shaded}
\begin{Highlighting}[]
\CommentTok{\# Cargo.toml}
\KeywordTok{[dependencies]}
\DataTypeTok{chakradb} \OperatorTok{=} \OperatorTok{\{ }\DataTypeTok{path}\OperatorTok{ =} \StringTok{"../chakradb"}\OperatorTok{, }\DataTypeTok{default{-}features}\OperatorTok{ =} \ConstantTok{false}\OperatorTok{ \}}
\end{Highlighting}
\end{Shaded}

Then, in code:

\begin{Shaded}
\begin{Highlighting}[]
\KeywordTok{use} \PreprocessorTok{chakradb::}\OperatorTok{\{}\NormalTok{Database}\OperatorTok{,}\NormalTok{ SqlEngine}\OperatorTok{\};}
\KeywordTok{use} \PreprocessorTok{std::sync::}\NormalTok{Arc}\OperatorTok{;}

\KeywordTok{let}\NormalTok{ db }\OperatorTok{=} \PreprocessorTok{Arc::}\NormalTok{new(}\PreprocessorTok{Database::}\NormalTok{new())}\OperatorTok{;}     \CommentTok{// in{-}memory}
\KeywordTok{let}\NormalTok{ sql }\OperatorTok{=} \PreprocessorTok{SqlEngine::}\NormalTok{new(db}\OperatorTok{.}\NormalTok{clone())}\OperatorTok{;}
\NormalTok{sql}\OperatorTok{.}\NormalTok{run(}\StringTok{"CREATE TABLE t (id INTEGER PRIMARY KEY, v INTEGER)"}\NormalTok{)}\OperatorTok{?;}
\end{Highlighting}
\end{Shaded}

For a durable, on-disk database, open a \texttt{Storage} over a
directory instead --- see \href{getting-started.md}{Your First Database
(Rust)}.

\section{Building the Python
bindings}\label{building-the-python-bindings}

The Python package is a thin PyO3 extension over the same core. From
\texttt{bindings/python}:

\begin{Shaded}
\begin{Highlighting}[]
\BuiltInTok{cd}\NormalTok{ bindings/python}
\ExtensionTok{maturin}\NormalTok{ develop }\AttributeTok{{-}{-}release}        \CommentTok{\# builds the extension into your venv}
\CommentTok{\# or, for a wheel you can install elsewhere:}
\ExtensionTok{maturin}\NormalTok{ build }\AttributeTok{{-}{-}release}
\end{Highlighting}
\end{Shaded}

Once built, the package is importable anywhere:

\begin{Shaded}
\begin{Highlighting}[]
\ImportTok{import}\NormalTok{ chakradb}
\NormalTok{conn }\OperatorTok{=}\NormalTok{ chakradb.}\ExtensionTok{connect}\NormalTok{(}\StringTok{":memory:"}\NormalTok{)}
\end{Highlighting}
\end{Shaded}

The Python driver implements \href{getting-started.md}{PEP 249 (DB-API
2.0)} and exposes the graph engine via \texttt{conn.graph(name)} --- see
\href{getting-started.md}{A Graph in Five Minutes}.

\section{Running the examples}\label{running-the-examples}

Two end-to-end examples ship in the repository:

\begin{Shaded}
\begin{Highlighting}[]
\CommentTok{\# Rust: a full real{-}time AML system over synthetic data}
\ExtensionTok{cargo}\NormalTok{ run }\AttributeTok{{-}{-}release} \AttributeTok{{-}{-}example}\NormalTok{ aml\_realtime }\AttributeTok{{-}{-}no{-}default{-}features}

\CommentTok{\# Python: the same application through the client bindings}
\ExtensionTok{python}\NormalTok{ examples/aml\_app.py}
\end{Highlighting}
\end{Shaded}

Both generate their own data, run the graph-analytics ensemble, and
verify that every planted laundering typology is detected. They are the
reference for the \href{../case-studies/aml.md}{Real-Time AML case
study}.

\section{Your First Database (Rust)}\label{your-first-database-rust}

This chapter walks through a complete Rust program: create a durable
database, define a table with real constraints and types, insert rows in
a transaction, and query them back. Every feature here is exercised by
the test suite.

\section{An in-memory database}\label{an-in-memory-database}

The smallest possible program:

\begin{Shaded}
\begin{Highlighting}[]
\KeywordTok{use} \PreprocessorTok{chakradb::}\OperatorTok{\{}\NormalTok{Database}\OperatorTok{,}\NormalTok{ SqlEngine}\OperatorTok{\};}
\KeywordTok{use} \PreprocessorTok{std::sync::}\NormalTok{Arc}\OperatorTok{;}

\KeywordTok{fn}\NormalTok{ main() }\OperatorTok{{-}\textgreater{}} \DataTypeTok{Result}\OperatorTok{\textless{}}\NormalTok{()}\OperatorTok{,} \DataTypeTok{Box}\OperatorTok{\textless{}}\KeywordTok{dyn} \PreprocessorTok{std::error::}\BuiltInTok{Error}\OperatorTok{\textgreater{}\textgreater{}} \OperatorTok{\{}
    \KeywordTok{let}\NormalTok{ db }\OperatorTok{=} \PreprocessorTok{Arc::}\NormalTok{new(}\PreprocessorTok{Database::}\NormalTok{new())}\OperatorTok{;}
    \KeywordTok{let}\NormalTok{ sql }\OperatorTok{=} \PreprocessorTok{SqlEngine::}\NormalTok{new(db}\OperatorTok{.}\NormalTok{clone())}\OperatorTok{;}

\NormalTok{    sql}\OperatorTok{.}\NormalTok{run(}\StringTok{"CREATE TABLE accounts (}
\StringTok{        id     INTEGER PRIMARY KEY,}
\StringTok{        owner  VARCHAR(64) NOT NULL,}
\StringTok{        balance DECIMAL(14,2) NOT NULL DEFAULT 0.00}
\StringTok{    )"}\NormalTok{)}\OperatorTok{?;}

\NormalTok{    sql}\OperatorTok{.}\NormalTok{run(}\StringTok{"INSERT INTO accounts VALUES (1, \textquotesingle{}ada\textquotesingle{}, 500.00)"}\NormalTok{)}\OperatorTok{?;}
\NormalTok{    sql}\OperatorTok{.}\NormalTok{run(}\StringTok{"INSERT INTO accounts VALUES (2, \textquotesingle{}grace\textquotesingle{}, 1250.50)"}\NormalTok{)}\OperatorTok{?;}

    \KeywordTok{let}\NormalTok{ rows }\OperatorTok{=}\NormalTok{ sql}\OperatorTok{.}\NormalTok{query(}\StringTok{"SELECT id, owner, balance FROM accounts ORDER BY id"}\NormalTok{)}\OperatorTok{?;}
    \ControlFlowTok{for}\NormalTok{ r }\KeywordTok{in}\NormalTok{ rows }\OperatorTok{\{}
        \PreprocessorTok{println!}\NormalTok{(}\StringTok{"\{r:?\}"}\NormalTok{)}\OperatorTok{;}
    \OperatorTok{\}}
    \ConstantTok{Ok}\NormalTok{(())}
\OperatorTok{\}}
\end{Highlighting}
\end{Shaded}

\texttt{SqlEngine::run} executes a statement (DDL or DML) and returns an
\href{getting-started.md}{\texttt{Outcome}}; \texttt{SqlEngine::query}
is a convenience that returns rows already rendered as
\texttt{Vec\textless{}Vec\textless{}String\textgreater{}\textgreater{}}.

\section{Types and constraints are
enforced}\label{types-and-constraints-are-enforced}

ChakraDB is not a stringly-typed store. The declared types and
constraints are checked at write time:

\begin{Shaded}
\begin{Highlighting}[]
\CommentTok{// Rejected: NOT NULL violated.}
\PreprocessorTok{assert!}\NormalTok{(sql}\OperatorTok{.}\NormalTok{run(}\StringTok{"INSERT INTO accounts (id, owner) VALUES (3, NULL)"}\NormalTok{)}\OperatorTok{.}\NormalTok{is\_err())}\OperatorTok{;}

\CommentTok{// Rejected: DECIMAL(14,2) — an over{-}precise literal is a constraint error, not}
\CommentTok{// a silent rounding.}
\PreprocessorTok{assert!}\NormalTok{(sql}\OperatorTok{.}\NormalTok{run(}\StringTok{"INSERT INTO accounts VALUES (3, \textquotesingle{}lin\textquotesingle{}, 1.234)"}\NormalTok{)}\OperatorTok{.}\NormalTok{is\_err())}\OperatorTok{;}

\CommentTok{// Rejected: VARCHAR(64) length is enforced.}
\CommentTok{// Rejected: duplicate PRIMARY KEY.}
\PreprocessorTok{assert!}\NormalTok{(sql}\OperatorTok{.}\NormalTok{run(}\StringTok{"INSERT INTO accounts VALUES (1, \textquotesingle{}dup\textquotesingle{}, 0.00)"}\NormalTok{)}\OperatorTok{.}\NormalTok{is\_err())}\OperatorTok{;}
\end{Highlighting}
\end{Shaded}

Money is stored as \textbf{exact} fixed-point \texttt{DECIMAL}, never
binary floating point --- see \href{../engine/data-types.md}{Exact
Decimal Arithmetic}.

\section{Transactions}\label{transactions-1}

By default each statement autocommits. Wrap several in a transaction for
all-or-nothing semantics under snapshot isolation:

\begin{Shaded}
\begin{Highlighting}[]
\NormalTok{sql}\OperatorTok{.}\NormalTok{begin()}\OperatorTok{?;}
\NormalTok{sql}\OperatorTok{.}\NormalTok{run(}\StringTok{"INSERT INTO accounts VALUES (10, \textquotesingle{}noether\textquotesingle{}, 300.00)"}\NormalTok{)}\OperatorTok{?;}
\NormalTok{sql}\OperatorTok{.}\NormalTok{run(}\StringTok{"INSERT INTO accounts VALUES (11, \textquotesingle{}hopper\textquotesingle{}, 300.00)"}\NormalTok{)}\OperatorTok{?;}
\NormalTok{sql}\OperatorTok{.}\NormalTok{commit()}\OperatorTok{?;}                       \CommentTok{// both visible, atomically}
\CommentTok{// ... or sql.rollback()? to discard the whole unit.}
\end{Highlighting}
\end{Shaded}

Readers never block writers and writers never block readers: a
transaction sees a consistent \href{../engine/mvcc.md}{MVCC snapshot}
taken at its start.

\section{A durable, on-disk database}\label{a-durable-on-disk-database}

Swap the in-memory \texttt{Database} for a \texttt{Storage} opened over
a directory. Nothing else in your code changes --- the SQL surface is
identical:

\begin{Shaded}
\begin{Highlighting}[]
\KeywordTok{use} \PreprocessorTok{chakradb::}\OperatorTok{\{}\NormalTok{PosixIo}\OperatorTok{,}\NormalTok{ SqlEngine}\OperatorTok{\};}
\KeywordTok{use} \PreprocessorTok{chakradb::storage::}\OperatorTok{\{}\NormalTok{Storage}\OperatorTok{,}\NormalTok{ StorageConfig}\OperatorTok{\};}
\KeywordTok{use} \PreprocessorTok{std::sync::}\NormalTok{Arc}\OperatorTok{;}

\KeywordTok{let}\NormalTok{ io }\OperatorTok{=} \PreprocessorTok{PosixIo::}\NormalTok{new(}\StringTok{"/var/lib/myapp"}\NormalTok{)}\OperatorTok{?;}
\KeywordTok{let}\NormalTok{ storage }\OperatorTok{=} \PreprocessorTok{Storage::}\NormalTok{open(}\PreprocessorTok{Arc::}\NormalTok{new(io)}\OperatorTok{,} \PreprocessorTok{StorageConfig::}\KeywordTok{default}\NormalTok{())}\OperatorTok{?;}
\KeywordTok{let}\NormalTok{ sql }\OperatorTok{=} \PreprocessorTok{SqlEngine::}\NormalTok{durable(}\PreprocessorTok{Arc::}\NormalTok{new(storage))}\OperatorTok{;}
\end{Highlighting}
\end{Shaded}

Writes are made durable through a write-ahead log with group commit, and
the engine recovers automatically on the next open --- see
\href{../engine/durability.md}{Durability} and
\href{../engine/durability.md}{Crash Recovery}.

\section{Where to go next}\label{where-to-go-next}

\begin{itemize}
\tightlist
\item
  \href{getting-started.md}{Python in Five Minutes} --- the same engine
  from Python.
\item
  \href{getting-started.md}{A Graph in Five Minutes} --- treat your
  tables as a graph.
\item
  \href{getting-started.md}{The SQL Reference} --- the full statement
  surface.
\end{itemize}

\section{Python in Five Minutes}\label{python-in-five-minutes}

The Python bindings wrap the exact same engine as the Rust core. The SQL
surface is exposed through a standard \href{getting-started.md}{PEP 249
(DB-API 2.0)} driver, so if you have used \texttt{sqlite3} you already
know the shape of it.

\section{Connect, execute, fetch}\label{connect-execute-fetch}

\begin{Shaded}
\begin{Highlighting}[]
\ImportTok{import}\NormalTok{ chakradb}

\NormalTok{conn }\OperatorTok{=}\NormalTok{ chakradb.}\ExtensionTok{connect}\NormalTok{(}\StringTok{":memory:"}\NormalTok{)          }\CommentTok{\# or a directory path for durability}
\NormalTok{cur }\OperatorTok{=}\NormalTok{ conn.execute(}\StringTok{"""}
\StringTok{    CREATE TABLE accounts (}
\StringTok{        id      INTEGER PRIMARY KEY,}
\StringTok{        owner   VARCHAR(64) NOT NULL,}
\StringTok{        balance DECIMAL(14,2) NOT NULL DEFAULT 0.00}
\StringTok{    )"""}\NormalTok{)}

\NormalTok{conn.execute(}\StringTok{"INSERT INTO accounts VALUES (1, \textquotesingle{}ada\textquotesingle{}, 500.00)"}\NormalTok{)}
\NormalTok{conn.execute(}\StringTok{"INSERT INTO accounts VALUES (2, \textquotesingle{}grace\textquotesingle{}, 1250.50)"}\NormalTok{)}
\NormalTok{conn.commit()}

\ControlFlowTok{for}\NormalTok{ row }\KeywordTok{in}\NormalTok{ conn.execute(}\StringTok{"SELECT id, owner, balance FROM accounts ORDER BY id"}\NormalTok{):}
    \BuiltInTok{print}\NormalTok{(row)          }\CommentTok{\# (1, \textquotesingle{}ada\textquotesingle{}, \textquotesingle{}500.00\textquotesingle{})  ...}
\end{Highlighting}
\end{Shaded}

\texttt{connect}, \texttt{Connection.execute},
\texttt{Cursor.fetchone/fetchmany/fetchall}, \texttt{commit},
\texttt{rollback}, and the full \href{getting-started.md}{exception
hierarchy} (\texttt{IntegrityError}, \texttt{ProgrammingError},
\texttt{OperationalError}, \ldots) all behave exactly as PEP 249
specifies.

\section{Parameters and transactions}\label{parameters-and-transactions}

Parameters use the \texttt{qmark} style; transactions are explicit
around a unit of work:

\begin{Shaded}
\begin{Highlighting}[]
\NormalTok{conn.execute(}\StringTok{"INSERT INTO accounts VALUES (?, ?, ?)"}\NormalTok{, (}\DecValTok{3}\NormalTok{, }\StringTok{\textquotesingle{}lin\textquotesingle{}}\NormalTok{, }\FloatTok{42.00}\NormalTok{))}

\ControlFlowTok{try}\NormalTok{:}
\NormalTok{    conn.execute(}\StringTok{"BEGIN"}\NormalTok{)  }\CommentTok{\# or rely on autocommit{-}off via .execute + .commit}
\NormalTok{    conn.execute(}\StringTok{"INSERT INTO accounts VALUES (10, \textquotesingle{}noether\textquotesingle{}, 300.00)"}\NormalTok{)}
\NormalTok{    conn.execute(}\StringTok{"INSERT INTO accounts VALUES (11, \textquotesingle{}hopper\textquotesingle{}, 300.00)"}\NormalTok{)}
\NormalTok{    conn.commit()}
\ControlFlowTok{except}\NormalTok{ chakradb.IntegrityError }\ImportTok{as}\NormalTok{ e:}
\NormalTok{    conn.rollback()}
    \BuiltInTok{print}\NormalTok{(}\StringTok{"rejected:"}\NormalTok{, e)}
\end{Highlighting}
\end{Shaded}

Type and constraint violations surface as \texttt{IntegrityError};
malformed SQL as \texttt{ProgrammingError}; a write--write conflict as
\texttt{OperationalError}.

\section{The graph engine, from
Python}\label{the-graph-engine-from-python}

Any table can be viewed as a graph. \texttt{conn.graph(name)} returns a
handle backed by table \texttt{name} in the same database; edges are
ordinary transactional rows:

\begin{Shaded}
\begin{Highlighting}[]
\NormalTok{g }\OperatorTok{=}\NormalTok{ conn.graph(}\StringTok{"payments"}\NormalTok{)}
\NormalTok{g.add\_edges([(}\DecValTok{1}\NormalTok{, }\DecValTok{2}\NormalTok{, }\FloatTok{100.0}\NormalTok{), (}\DecValTok{2}\NormalTok{, }\DecValTok{3}\NormalTok{, }\FloatTok{250.0}\NormalTok{), (}\DecValTok{3}\NormalTok{, }\DecValTok{1}\NormalTok{, }\FloatTok{90.0}\NormalTok{)])   }\CommentTok{\# (src, dst, weight)}

\NormalTok{view }\OperatorTok{=}\NormalTok{ g.view()                        }\CommentTok{\# one consistent snapshot for analytics}
\BuiltInTok{print}\NormalTok{(view.laundering\_cycles())        }\CommentTok{\# [[1, 2, 3]] — a round{-}trip ring}
\BuiltInTok{print}\NormalTok{(view.pagerank())                 }\CommentTok{\# \{1: 0.33, 2: 0.33, 3: 0.33\}}
\BuiltInTok{print}\NormalTok{(view.personalized\_pagerank(seeds}\OperatorTok{=}\NormalTok{[}\DecValTok{1}\NormalTok{]))}
\end{Highlighting}
\end{Shaded}

Every algorithm returns plain Python \texttt{dict}/\texttt{list} values
keyed by node id. The full catalogue is in \href{getting-started.md}{A
Graph in Five Minutes} and \href{../graph/algorithms.md}{Graph
Algorithms}.

\section{A complete application}\label{a-complete-application}

\texttt{examples/aml\_app.py} is a full real-time anti-money-laundering
system in \textasciitilde300 lines of Python: it builds a synthetic
payment network, runs an ensemble of graph detectors, and ranks
suspicious accounts. Walk through it in the
\href{../case-studies/aml.md}{Real-Time AML case study}.

\section{A Graph in Five Minutes}\label{a-graph-in-five-minutes}

ChakraDB has a graph engine built into the core --- not a bolt-on, not a
separate store. A graph \emph{is} a table of edges, clustered so that a
node's neighbours sit together on disk; analytics run over a consistent
\href{../engine/mvcc.md}{MVCC snapshot} of that table. This chapter is
the whirlwind tour.

\section{Open a graph and add edges}\label{open-a-graph-and-add-edges}

A graph is backed by a named table in an ordinary database. In Rust:

\begin{Shaded}
\begin{Highlighting}[]
\KeywordTok{use} \PreprocessorTok{chakradb::}\OperatorTok{\{}\NormalTok{Database}\OperatorTok{,}\NormalTok{ Graph}\OperatorTok{\};}
\KeywordTok{use} \PreprocessorTok{std::sync::}\NormalTok{Arc}\OperatorTok{;}

\KeywordTok{let}\NormalTok{ db }\OperatorTok{=} \PreprocessorTok{Arc::}\NormalTok{new(}\PreprocessorTok{Database::}\NormalTok{new())}\OperatorTok{;}
\KeywordTok{let}\NormalTok{ g }\OperatorTok{=} \PreprocessorTok{Graph::}\NormalTok{open(db}\OperatorTok{.}\NormalTok{clone()}\OperatorTok{,} \StringTok{"transfers"}\NormalTok{)}\OperatorTok{?;}

\CommentTok{// (src, dst, weight) — inserts are upserts, and fully transactional.}
\NormalTok{g}\OperatorTok{.}\NormalTok{add\_edges([(}\DecValTok{1}\OperatorTok{,} \DecValTok{2}\OperatorTok{,} \DecValTok{100.0}\NormalTok{)}\OperatorTok{,}\NormalTok{ (}\DecValTok{2}\OperatorTok{,} \DecValTok{3}\OperatorTok{,} \DecValTok{250.0}\NormalTok{)}\OperatorTok{,}\NormalTok{ (}\DecValTok{3}\OperatorTok{,} \DecValTok{1}\OperatorTok{,} \DecValTok{90.0}\NormalTok{)])}\OperatorTok{?;}
\end{Highlighting}
\end{Shaded}

In Python, through an existing connection:

\begin{Shaded}
\begin{Highlighting}[]
\NormalTok{g }\OperatorTok{=}\NormalTok{ conn.graph(}\StringTok{"transfers"}\NormalTok{)}
\NormalTok{g.add\_edges([(}\DecValTok{1}\NormalTok{, }\DecValTok{2}\NormalTok{, }\FloatTok{100.0}\NormalTok{), (}\DecValTok{2}\NormalTok{, }\DecValTok{3}\NormalTok{, }\FloatTok{250.0}\NormalTok{), (}\DecValTok{3}\NormalTok{, }\DecValTok{1}\NormalTok{, }\FloatTok{90.0}\NormalTok{)])}
\end{Highlighting}
\end{Shaded}

Because edges are just rows, the \emph{same} data is visible to SQL
(\texttt{SELECT\ ...\ FROM\ transfers}) and to the graph engine
simultaneously --- that is the HTAP promise applied to graphs.

\section{Freeze a snapshot, then
compute}\label{freeze-a-snapshot-then-compute}

All algorithms run against a \texttt{view()} --- an immutable in-memory
\href{../graph/snapshot.md}{CSR} snapshot. Take it once and a whole
pipeline sees one coherent graph, even while writers keep appending
edges:

\begin{Shaded}
\begin{Highlighting}[]
\KeywordTok{let}\NormalTok{ view }\OperatorTok{=}\NormalTok{ g}\OperatorTok{.}\NormalTok{view()}\OperatorTok{?;}
\PreprocessorTok{println!}\NormalTok{(}\StringTok{"\{\} nodes, \{\} edges"}\OperatorTok{,}\NormalTok{ view}\OperatorTok{.}\NormalTok{node\_count()}\OperatorTok{,}\NormalTok{ view}\OperatorTok{.}\NormalTok{edge\_count())}\OperatorTok{;}

\KeywordTok{let}\NormalTok{ dist   }\OperatorTok{=}\NormalTok{ view}\OperatorTok{.}\NormalTok{bfs(}\DecValTok{1}\NormalTok{)}\OperatorTok{;}                       \CommentTok{// hops from node 1}
\KeywordTok{let}\NormalTok{ path   }\OperatorTok{=}\NormalTok{ view}\OperatorTok{.}\NormalTok{shortest\_path(}\DecValTok{1}\OperatorTok{,} \DecValTok{3}\NormalTok{)}\OperatorTok{;}          \CommentTok{// Some([1, 2, 3])}
\KeywordTok{let}\NormalTok{ cost   }\OperatorTok{=}\NormalTok{ view}\OperatorTok{.}\NormalTok{weighted\_shortest\_path(}\DecValTok{1}\OperatorTok{,} \DecValTok{3}\NormalTok{)}\OperatorTok{;} \CommentTok{// Dijkstra with weights}
\KeywordTok{let}\NormalTok{ rank   }\OperatorTok{=}\NormalTok{ view}\OperatorTok{.}\NormalTok{pagerank(}\DecValTok{20}\OperatorTok{,} \DecValTok{0.85}\NormalTok{)}\OperatorTok{;}           \CommentTok{// \{node: score\}}
\KeywordTok{let}\NormalTok{ rings  }\OperatorTok{=}\NormalTok{ view}\OperatorTok{.}\NormalTok{laundering\_cycles()}\OperatorTok{;}          \CommentTok{// non{-}trivial SCCs}
\end{Highlighting}
\end{Shaded}

The Python surface is identical, returning dicts and lists:

\begin{Shaded}
\begin{Highlighting}[]
\NormalTok{view }\OperatorTok{=}\NormalTok{ g.view()}
\NormalTok{view.bfs(}\DecValTok{1}\NormalTok{)                       }\CommentTok{\# \{1: 0, 2: 1, 3: 1\}}
\NormalTok{view.shortest\_path(}\DecValTok{1}\NormalTok{, }\DecValTok{3}\NormalTok{)          }\CommentTok{\# [1, 2, 3]}
\NormalTok{view.pagerank()                   }\CommentTok{\# \{1: 0.33, 2: 0.33, 3: 0.33\}}
\NormalTok{view.laundering\_cycles()          }\CommentTok{\# [[1, 2, 3]]}
\NormalTok{view.personalized\_pagerank([}\DecValTok{1}\NormalTok{])   }\CommentTok{\# risk spread from a seed set}
\end{Highlighting}
\end{Shaded}

\section{The algorithm catalogue}\label{the-algorithm-catalogue}

Every one of these is built into the core and available from both Rust
and Python. See \href{../graph/algorithms.md}{Graph Algorithms} for the
derivations.

\begin{longtable}[]{@{}
  >{\raggedright\arraybackslash}p{(\linewidth - 2\tabcolsep) * \real{0.4762}}
  >{\raggedright\arraybackslash}p{(\linewidth - 2\tabcolsep) * \real{0.5238}}@{}}
\toprule\noalign{}
\begin{minipage}[b]{\linewidth}\raggedright
Category
\end{minipage} & \begin{minipage}[b]{\linewidth}\raggedright
Algorithms
\end{minipage} \\
\midrule\noalign{}
\endhead
\bottomrule\noalign{}
\endlastfoot
\textbf{Traversal \& paths} & \texttt{bfs}, \texttt{shortest\_path},
\texttt{dijkstra}, \texttt{weighted\_shortest\_path},
\texttt{topological\_order} \\
\textbf{Centrality} & \texttt{pagerank},
\texttt{personalized\_pagerank}, \texttt{degree\_centrality},
\texttt{closeness\_centrality}, \texttt{betweenness\_centrality} \\
\textbf{Community \& structure} & \texttt{connected\_components},
\texttt{strongly\_connected\_components}, \texttt{laundering\_cycles},
\texttt{label\_propagation}, \texttt{k\_core},
\texttt{triangle\_count} \\
\textbf{Systemic risk} & \texttt{eisenberg\_noe} (clearing vector /
default cascade) \\
\textbf{Similarity \& link prediction} & \texttt{common\_neighbors},
\texttt{jaccard\_similarity}, \texttt{adamic\_adar},
\texttt{recommend} \\
\textbf{Degrees} & \texttt{in\_degree}, \texttt{out\_degree},
\texttt{in\_neighbors}, \texttt{out\_neighbors} \\
\end{longtable}

\section{Snapshots are stable under
writes}\label{snapshots-are-stable-under-writes}

The defining property --- an analysis is reproducible because its view
does not move under it:

\begin{Shaded}
\begin{Highlighting}[]
\KeywordTok{let}\NormalTok{ v }\OperatorTok{=}\NormalTok{ g}\OperatorTok{.}\NormalTok{view()}\OperatorTok{?;}
\KeywordTok{let}\NormalTok{ before }\OperatorTok{=}\NormalTok{ v}\OperatorTok{.}\NormalTok{edge\_count()}\OperatorTok{;}
\NormalTok{g}\OperatorTok{.}\NormalTok{add\_edges([(}\DecValTok{9}\OperatorTok{,} \DecValTok{10}\OperatorTok{,} \DecValTok{1.0}\NormalTok{)])}\OperatorTok{?;}           \CommentTok{// keep writing after the view is taken}
\PreprocessorTok{assert\_eq!}\NormalTok{(v}\OperatorTok{.}\NormalTok{edge\_count()}\OperatorTok{,}\NormalTok{ before)}\OperatorTok{;}     \CommentTok{// the snapshot did not change}
\PreprocessorTok{assert\_eq!}\NormalTok{(g}\OperatorTok{.}\NormalTok{view()}\OperatorTok{?.}\NormalTok{edge\_count()}\OperatorTok{,}\NormalTok{ before }\OperatorTok{+} \DecValTok{1}\NormalTok{)}\OperatorTok{;}  \CommentTok{// a fresh view sees the write}
\end{Highlighting}
\end{Shaded}

\section{A real application}\label{a-real-application}

The \href{../case-studies/aml.md}{Real-Time AML case study} composes
these primitives into a working fraud-detection system. The runnable
code is \texttt{examples/aml\_realtime.rs} (Rust) and
\texttt{examples/aml\_app.py} (Python).

\part{Case Studies}

\chapter{Case Study: A Real-Time Anti-Money-Laundering
System}\label{case-study-a-real-time-anti-money-laundering-system}

\epigraph{Follow the money.}{--- Deep Throat, \textit{All the President's Men}}

This is the chapter where every part of ChakraDB is used at once, on a
workload that genuinely needs all of it: \textbf{anti-money-laundering
(AML)}. Money laundering is a \emph{graph} problem wearing a
\emph{transactional} disguise --- a stream of individually innocuous
transfers whose \emph{shape}, taken together, is the crime. Detecting it
in real time means ingesting transactions transactionally, scoring them
analytically, and traversing the entity graph they form --- over one
consistent view, as the data arrives. That is exactly the
HTAP-plus-graph workload ChakraDB is built for.

We build a working detection engine: the schema, the ingestion path, a
detector for each of the major laundering typologies (mapped onto a
ChakraDB primitive), the scoring pipeline that combines them, and the
case-management output --- all in one embedded process, over live data.

\section{The problem}\label{the-problem}

Laundering has three classic stages, and each leaves a different
footprint in the data:

\begin{itemize}
\tightlist
\item
  \textbf{Placement} --- dirty cash enters the system (many small
  deposits). Footprint: \emph{structuring / smurfing} --- many small
  credits, kept under reporting thresholds, fanning \textbf{into} one
  account.
\item
  \textbf{Layering} --- the money is moved through a maze of transfers
  to obscure its origin. Footprint: \emph{round-tripping and chains} ---
  directed \textbf{cycles} and long paths through many accounts, often
  rapidly.
\item
  \textbf{Integration} --- the laundered money returns to the criminal
  as apparently legitimate funds. Footprint: \emph{mule networks} ---
  dense clusters of accounts that concentrate flow toward a few
  beneficiaries.
\end{itemize}

\begin{Shaded}
\begin{Highlighting}[]
\NormalTok{flowchart LR}
\NormalTok{    subgraph P["Placement (structuring)"]}
\NormalTok{      D1(( )) {-}{-}\textgreater{} M[Mule A]}
\NormalTok{      D2(( )) {-}{-}\textgreater{} M}
\NormalTok{      D3(( )) {-}{-}\textgreater{} M}
\NormalTok{      D4(( )) {-}{-}\textgreater{} M}
\NormalTok{    end}
\NormalTok{    subgraph L["Layering (round{-}trip)"]}
\NormalTok{      M {-}{-}\textgreater{} B[Acct B] {-}{-}\textgreater{} C[Acct C] {-}{-}\textgreater{} M}
\NormalTok{    end}
\NormalTok{    subgraph I["Integration"]}
\NormalTok{      C {-}{-}\textgreater{} Z[Beneficiary]}
\NormalTok{    end}
\NormalTok{    classDef s fill:\#f4c430; class M,Z s;}
\end{Highlighting}
\end{Shaded}

No single transfer is suspicious. The \emph{fan-in} at the mule, the
\emph{cycle} B→C→A→B, and the \emph{concentration} toward the
beneficiary are --- and each is a graph query.

\section{Why one engine}\label{why-one-engine}

The traditional AML stack fragments this: an OLTP core banking system
for the transactions, a nightly batch to a warehouse for the rules, and
a separate graph database for network analysis --- reconciled by ETL.
The consequences are exactly the ones that matter for AML: the graph is
\textbf{hours stale}, so layering is caught after the money is gone; the
systems disagree on \emph{which} transactions are in scope; and running
the network analysis contends with ingestion.

ChakraDB collapses the stack. Transactions, rule analytics, and graph
traversal all read the \textbf{same MVCC snapshot} of the \textbf{same
live data}, in one process, with writers never blocked by the scorer.
The detector below runs \emph{as transactions arrive}.

\section{The schema}\label{the-schema}

Two logical layers: the transactional records, and the entity/flow graph
derived from them.

\begin{Shaded}
\begin{Highlighting}[]
\CommentTok{{-}{-} Customers and their risk attributes.}
\KeywordTok{CREATE} \KeywordTok{TABLE}\NormalTok{ customers (}
  \KeywordTok{id}          \DataTypeTok{INT} \KeywordTok{PRIMARY} \KeywordTok{KEY}\NormalTok{,}
\NormalTok{  name        }\DataTypeTok{VARCHAR}\NormalTok{(}\DecValTok{128}\NormalTok{) }\KeywordTok{NOT} \KeywordTok{NULL}\NormalTok{,}
\NormalTok{  country     }\DataTypeTok{VARCHAR}\NormalTok{(}\DecValTok{2}\NormalTok{),}
\NormalTok{  risk\_rating }\DataTypeTok{VARCHAR}\NormalTok{(}\DecValTok{8}\NormalTok{) }\KeywordTok{DEFAULT} \StringTok{\textquotesingle{}low\textquotesingle{}}\NormalTok{,     }\CommentTok{{-}{-} low | medium | high}
\NormalTok{  is\_pep      }\DataTypeTok{BOOLEAN} \KeywordTok{DEFAULT} \KeywordTok{false}\NormalTok{,        }\CommentTok{{-}{-} politically{-}exposed person}
\NormalTok{  onboarded   }\DataTypeTok{DATE}
\NormalTok{);}

\CommentTok{{-}{-} Accounts, each owned by a customer.}
\KeywordTok{CREATE} \KeywordTok{TABLE}\NormalTok{ accounts (}
  \KeywordTok{id}          \DataTypeTok{INT} \KeywordTok{PRIMARY} \KeywordTok{KEY}\NormalTok{,}
\NormalTok{  customer\_id }\DataTypeTok{INT} \KeywordTok{NOT} \KeywordTok{NULL}\NormalTok{,}
\NormalTok{  opened      }\DataTypeTok{DATE}\NormalTok{,}
\NormalTok{  status      }\DataTypeTok{VARCHAR}\NormalTok{(}\DecValTok{8}\NormalTok{) }\KeywordTok{DEFAULT} \StringTok{\textquotesingle{}active\textquotesingle{}}
\NormalTok{);}

\CommentTok{{-}{-} The transaction ledger — exact money, never rounded.}
\KeywordTok{CREATE} \KeywordTok{TABLE}\NormalTok{ txns (}
  \KeywordTok{id}       \DataTypeTok{INT} \KeywordTok{PRIMARY} \KeywordTok{KEY}\NormalTok{,}
\NormalTok{  src\_acct }\DataTypeTok{INT} \KeywordTok{NOT} \KeywordTok{NULL}\NormalTok{,}
\NormalTok{  dst\_acct }\DataTypeTok{INT} \KeywordTok{NOT} \KeywordTok{NULL}\NormalTok{,}
\NormalTok{  amount   }\DataTypeTok{DECIMAL}\NormalTok{(}\DecValTok{14}\NormalTok{,}\DecValTok{2}\NormalTok{) }\KeywordTok{NOT} \KeywordTok{NULL} \KeywordTok{CHECK}\NormalTok{ (amount }\OperatorTok{\textgreater{}} \DecValTok{0}\NormalTok{),}
\NormalTok{  ts       }\DataTypeTok{TIMESTAMP} \KeywordTok{NOT} \KeywordTok{NULL}\NormalTok{,}
\NormalTok{  channel  }\DataTypeTok{VARCHAR}\NormalTok{(}\DecValTok{12}\NormalTok{),                     }\CommentTok{{-}{-} wire | ach | card | crypto}
  \KeywordTok{CHECK}\NormalTok{ (src\_acct }\OperatorTok{\textless{}\textgreater{}}\NormalTok{ dst\_acct)}
\NormalTok{);}

\CommentTok{{-}{-} A watchlist of accounts already known to be bad (SARs filed, sanctions, ...).}
\KeywordTok{CREATE} \KeywordTok{TABLE}\NormalTok{ known\_bad (}
\NormalTok{  acct     }\DataTypeTok{INT} \KeywordTok{PRIMARY} \KeywordTok{KEY}\NormalTok{,}
\NormalTok{  reason   }\DataTypeTok{VARCHAR}\NormalTok{(}\DecValTok{32}\NormalTok{),}
\NormalTok{  added    }\DataTypeTok{DATE}
\NormalTok{);}
\end{Highlighting}
\end{Shaded}

The \textbf{flow graph} is derived: a node is an account, and a directed
edge \texttt{src\ →\ dst} means ``money moved from src to dst,''
weighted by amount. ChakraDB's \texttt{Graph} handle maintains it as an
edges table keyed \texttt{(src,\ dst)} so adjacency is clustered
(\href{../graph/model.md}{The Graph Model}).

\section{Ingestion}\label{ingestion}

Each incoming transaction is one ACID transaction: the ledger row and
the flow edge commit together, so the graph is never half-updated
relative to the ledger.

\begin{Shaded}
\begin{Highlighting}[]
\KeywordTok{use} \PreprocessorTok{chakradb::}\OperatorTok{\{}\NormalTok{SqlEngine}\OperatorTok{,}\NormalTok{ Graph}\OperatorTok{\};}

\KeywordTok{fn}\NormalTok{ ingest(sql}\OperatorTok{:} \OperatorTok{\&}\NormalTok{SqlEngine}\OperatorTok{,}\NormalTok{ flow}\OperatorTok{:} \OperatorTok{\&}\NormalTok{Graph}\OperatorTok{,}\NormalTok{ t}\OperatorTok{:} \OperatorTok{\&}\NormalTok{Txn) }\OperatorTok{{-}\textgreater{}} \PreprocessorTok{anyhow::}\DataTypeTok{Result}\OperatorTok{\textless{}}\NormalTok{()}\OperatorTok{\textgreater{}} \OperatorTok{\{}
\NormalTok{    sql}\OperatorTok{.}\NormalTok{run(}\OperatorTok{\&}\PreprocessorTok{format!}\NormalTok{(}
        \StringTok{"BEGIN;}
\StringTok{         INSERT INTO txns VALUES (\{\}, \{\}, \{\}, \{:.2\}, \textquotesingle{}\{\}\textquotesingle{}, \textquotesingle{}\{\}\textquotesingle{});}
\StringTok{         COMMIT;"}\OperatorTok{,}
\NormalTok{        t}\OperatorTok{.}\NormalTok{id}\OperatorTok{,}\NormalTok{ t}\OperatorTok{.}\NormalTok{src\_acct}\OperatorTok{,}\NormalTok{ t}\OperatorTok{.}\NormalTok{dst\_acct}\OperatorTok{,}\NormalTok{ t}\OperatorTok{.}\NormalTok{amount}\OperatorTok{,}\NormalTok{ t}\OperatorTok{.}\NormalTok{ts}\OperatorTok{,}\NormalTok{ t}\OperatorTok{.}\NormalTok{channel}
\NormalTok{    ))}\OperatorTok{?;}
    \CommentTok{// The flow edge, weighted by amount (idempotent upsert on (src,dst)).}
\NormalTok{    flow}\OperatorTok{.}\NormalTok{add\_edge(t}\OperatorTok{.}\NormalTok{src\_acct}\OperatorTok{,}\NormalTok{ t}\OperatorTok{.}\NormalTok{dst\_acct}\OperatorTok{,}\NormalTok{ t}\OperatorTok{.}\NormalTok{amount)}\OperatorTok{?;}
    \ConstantTok{Ok}\NormalTok{(())}
\OperatorTok{\}}
\end{Highlighting}
\end{Shaded}

\texttt{DECIMAL(14,2)} keeps amounts exact; the \texttt{CHECK}s reject
nonsense before it reaches the log. Ingestion never pauses for the
detector --- the wedge.

\section{The detection engine}\label{the-detection-engine}

Each laundering typology maps onto a ChakraDB primitive. We take
\textbf{one snapshot} --- a SQL view of the current state and a
\texttt{GraphView} (CSR) built from the same instant
(\href{../graph/snapshot.md}{Live Graph Analytics}) --- so every
detector sees a consistent graph.

\begin{Shaded}
\begin{Highlighting}[]
\KeywordTok{let}\NormalTok{ view }\OperatorTok{=}\NormalTok{ flow}\OperatorTok{.}\NormalTok{view()}\OperatorTok{?;}   \CommentTok{// consistent CSR snapshot; ingestion continues}
\end{Highlighting}
\end{Shaded}

\subsection{1. Structuring / smurfing --- graph fan-in + SQL
velocity}\label{structuring-smurfing-graph-fan-in-sql-velocity}

Many small credits fanning \textbf{into} one account, kept under a
reporting threshold, in a short window. Fan-in is \texttt{in\_degree};
the threshold and velocity are SQL.

\begin{Shaded}
\begin{Highlighting}[]
\CommentTok{// Candidate mules: high fan{-}in of small, sub{-}threshold credits in 24h.}
\KeywordTok{fn}\NormalTok{ structuring\_suspects(sql}\OperatorTok{:} \OperatorTok{\&}\NormalTok{SqlEngine}\OperatorTok{,}\NormalTok{ view}\OperatorTok{:} \OperatorTok{\&}\NormalTok{GraphView}\OperatorTok{,}\NormalTok{ acct}\OperatorTok{:}\NormalTok{ NodeId) }\OperatorTok{{-}\textgreater{}} \DataTypeTok{bool} \OperatorTok{\{}
    \KeywordTok{let}\NormalTok{ fan\_in }\OperatorTok{=}\NormalTok{ view}\OperatorTok{.}\NormalTok{in\_degree(acct)}\OperatorTok{;}                 \CommentTok{// graph: many senders}
    \KeywordTok{let}\NormalTok{ small\_credits}\OperatorTok{:} \DataTypeTok{i64} \OperatorTok{=}\NormalTok{ sql}\OperatorTok{.}\NormalTok{query(}\OperatorTok{\&}\PreprocessorTok{format!}\NormalTok{(}
        \StringTok{"SELECT COUNT(*) FROM txns}
\StringTok{         WHERE dst\_acct = \{acct\} AND amount \textless{} 10000.00}
\StringTok{           AND ts \textgreater{} NOW\_MINUS\_24H"}\NormalTok{))}\OperatorTok{?}\NormalTok{[}\DecValTok{0}\NormalTok{][}\DecValTok{0}\NormalTok{]}\OperatorTok{.}\NormalTok{parse()}\OperatorTok{.}\NormalTok{unwrap\_or(}\DecValTok{0}\NormalTok{)}\OperatorTok{;}
\NormalTok{    fan\_in }\OperatorTok{\textgreater{}=} \DecValTok{8} \OperatorTok{\&\&}\NormalTok{ small\_credits }\OperatorTok{\textgreater{}=} \DecValTok{8}                  \CommentTok{// tune per institution}
\OperatorTok{\}}
\end{Highlighting}
\end{Shaded}

\subsection{2. Round-tripping / layering --- laundering cycles
(SCC)}\label{round-tripping-layering-laundering-cycles-scc}

Funds that can return to their origin through a chain of transfers form
a \textbf{directed cycle} --- a strongly-connected component of size ≥
2. \texttt{laundering\_cycles} finds them directly
(\href{../graph/algorithms.md}{Graph Algorithms}).

\begin{Shaded}
\begin{Highlighting}[]
\CommentTok{// Every set of accounts among which money can circulate.}
\KeywordTok{let}\NormalTok{ cycles}\OperatorTok{:} \DataTypeTok{Vec}\OperatorTok{\textless{}}\DataTypeTok{Vec}\OperatorTok{\textless{}}\NormalTok{NodeId}\OperatorTok{\textgreater{}\textgreater{}} \OperatorTok{=}\NormalTok{ view}\OperatorTok{.}\NormalTok{laundering\_cycles()}\OperatorTok{;}
\ControlFlowTok{for}\NormalTok{ ring }\KeywordTok{in} \OperatorTok{\&}\NormalTok{cycles }\OperatorTok{\{}
\NormalTok{    alert(}\PreprocessorTok{Alert::}\NormalTok{layering(ring}\OperatorTok{,} \StringTok{"funds can round{-}trip among these accounts"}\NormalTok{))}\OperatorTok{;}
\OperatorTok{\}}
\end{Highlighting}
\end{Shaded}

A single \texttt{laundering\_cycles()} call replaces the recursive
path-search that a relational-only system cannot express (SQL has no
recursion in ChakraDB's surface --- which is \emph{why} the algorithm
lives in the engine).

\subsection{3. Mule networks --- connected components +
centrality}\label{mule-networks-connected-components-centrality}

A laundering operation is a \textbf{dense cluster} that concentrates
flow toward a few beneficiaries. Weakly-connected components isolate the
clusters; PageRank finds the beneficiaries within one.

\begin{Shaded}
\begin{Highlighting}[]
\KeywordTok{let}\NormalTok{ comps }\OperatorTok{=}\NormalTok{ view}\OperatorTok{.}\NormalTok{connected\_components()}\OperatorTok{;}               \CommentTok{// isolate clusters}
\KeywordTok{let}\NormalTok{ ranks }\OperatorTok{=}\NormalTok{ view}\OperatorTok{.}\NormalTok{pagerank(}\DecValTok{30}\OperatorTok{,} \DecValTok{0.85}\NormalTok{)}\OperatorTok{;}                   \CommentTok{// flow concentration}
\CommentTok{// A component with an unusually high{-}PageRank sink is a candidate mule network.}
\end{Highlighting}
\end{Shaded}

\subsection{4. Risk propagation from known-bad --- personalized
PageRank}\label{risk-propagation-from-known-bad-personalized-pagerank}

The most powerful signal: \textbf{proximity to accounts already known to
be bad}. Seed personalized PageRank with the watchlist, and every
account gets a risk score that is its exposure to the seeds --- decaying
with graph distance.

\begin{Shaded}
\begin{Highlighting}[]
\KeywordTok{let}\NormalTok{ seeds}\OperatorTok{:} \DataTypeTok{Vec}\OperatorTok{\textless{}}\NormalTok{NodeId}\OperatorTok{\textgreater{}} \OperatorTok{=}\NormalTok{ sql}\OperatorTok{.}\NormalTok{query(}\StringTok{"SELECT acct FROM known\_bad"}\NormalTok{)}\OperatorTok{?}
    \OperatorTok{.}\NormalTok{iter()}\OperatorTok{.}\NormalTok{map(}\OperatorTok{|}\NormalTok{r}\OperatorTok{|}\NormalTok{ r[}\DecValTok{0}\NormalTok{]}\OperatorTok{.}\NormalTok{parse()}\OperatorTok{.}\NormalTok{unwrap())}\OperatorTok{.}\NormalTok{collect()}\OperatorTok{;}
\KeywordTok{let}\NormalTok{ risk }\OperatorTok{=}\NormalTok{ view}\OperatorTok{.}\NormalTok{personalized\_pagerank(}\OperatorTok{\&}\NormalTok{seeds}\OperatorTok{,} \DecValTok{40}\OperatorTok{,} \DecValTok{0.85}\NormalTok{)}\OperatorTok{;}
\CommentTok{// risk[\&acct] is high for accounts transacting near the watchlist — even}
\CommentTok{// several hops away, which flat rules miss entirely.}
\end{Highlighting}
\end{Shaded}

\subsection{5. Rapid movement \& fan-out --- SQL + graph
out-degree}\label{rapid-movement-fan-out-sql-graph-out-degree}

Layering also shows as money that arrives and leaves almost immediately,
and as \emph{fan-out} (one account distributing to many). Velocity is a
temporal SQL query; fan-out is \texttt{out\_degree}.

\begin{Shaded}
\begin{Highlighting}[]
\KeywordTok{let}\NormalTok{ fan\_out }\OperatorTok{=}\NormalTok{ view}\OperatorTok{.}\NormalTok{out\_degree(acct)}\OperatorTok{;}                   \CommentTok{// distribution}
\KeywordTok{let}\NormalTok{ dwell\_ms}\OperatorTok{:} \DataTypeTok{i64} \OperatorTok{=}\NormalTok{ sql}\OperatorTok{.}\NormalTok{query(}\OperatorTok{\&}\PreprocessorTok{format!}\NormalTok{(}
    \StringTok{"SELECT MIN(out.ts {-} inn.ts) FROM txns inn, txns out}
\StringTok{     WHERE inn.dst\_acct = \{acct\} AND out.src\_acct = \{acct\}"}\NormalTok{))}\OperatorTok{?}\NormalTok{[}\DecValTok{0}\NormalTok{][}\DecValTok{0}\NormalTok{]}\OperatorTok{.}\NormalTok{parse()}\OperatorTok{?;}
\KeywordTok{let}\NormalTok{ rapid }\OperatorTok{=}\NormalTok{ dwell\_ms }\OperatorTok{\textless{}} \DecValTok{60\_000}\OperatorTok{;}                          \CommentTok{// in{-}and{-}out under a minute}
\end{Highlighting}
\end{Shaded}

\section{The scoring pipeline}\label{the-scoring-pipeline}

The detectors combine into a single risk score per account, computed
over one snapshot, and above a threshold they open a case:

\begin{Shaded}
\begin{Highlighting}[]
\KeywordTok{fn}\NormalTok{ score\_accounts(sql}\OperatorTok{:} \OperatorTok{\&}\NormalTok{SqlEngine}\OperatorTok{,}\NormalTok{ flow}\OperatorTok{:} \OperatorTok{\&}\NormalTok{Graph) }\OperatorTok{{-}\textgreater{}} \PreprocessorTok{anyhow::}\DataTypeTok{Result}\OperatorTok{\textless{}}\DataTypeTok{Vec}\OperatorTok{\textless{}}\NormalTok{Case}\OperatorTok{\textgreater{}\textgreater{}} \OperatorTok{\{}
    \KeywordTok{let}\NormalTok{ view }\OperatorTok{=}\NormalTok{ flow}\OperatorTok{.}\NormalTok{view()}\OperatorTok{?;}                            \CommentTok{// one consistent instant}
    \KeywordTok{let}\NormalTok{ seeds }\OperatorTok{=}\NormalTok{ known\_bad\_accounts(sql)}\OperatorTok{?;}
    \KeywordTok{let}\NormalTok{ risk  }\OperatorTok{=}\NormalTok{ view}\OperatorTok{.}\NormalTok{personalized\_pagerank(}\OperatorTok{\&}\NormalTok{seeds}\OperatorTok{,} \DecValTok{40}\OperatorTok{,} \DecValTok{0.85}\NormalTok{)}\OperatorTok{;}
    \KeywordTok{let}\NormalTok{ cycles }\OperatorTok{=}\NormalTok{ view}\OperatorTok{.}\NormalTok{laundering\_cycles()}\OperatorTok{;}
    \KeywordTok{let}\NormalTok{ in\_cycle}\OperatorTok{:}\NormalTok{ HashSet}\OperatorTok{\textless{}}\NormalTok{NodeId}\OperatorTok{\textgreater{}} \OperatorTok{=}\NormalTok{ cycles}\OperatorTok{.}\NormalTok{iter()}\OperatorTok{.}\NormalTok{flatten()}\OperatorTok{.}\NormalTok{copied()}\OperatorTok{.}\NormalTok{collect()}\OperatorTok{;}

    \KeywordTok{let} \KeywordTok{mut}\NormalTok{ cases }\OperatorTok{=} \DataTypeTok{Vec}\PreprocessorTok{::}\NormalTok{new()}\OperatorTok{;}
    \ControlFlowTok{for} \OperatorTok{\&}\NormalTok{acct }\KeywordTok{in}\NormalTok{ view\_accounts(}\OperatorTok{\&}\NormalTok{view) }\OperatorTok{\{}
        \KeywordTok{let} \KeywordTok{mut}\NormalTok{ score }\OperatorTok{=} \DecValTok{0.0}\OperatorTok{;}
        \KeywordTok{let} \KeywordTok{mut}\NormalTok{ reasons }\OperatorTok{=} \DataTypeTok{Vec}\PreprocessorTok{::}\NormalTok{new()}\OperatorTok{;}

        \CommentTok{// (1) proximity to known{-}bad (the dominant term)}
        \ControlFlowTok{if} \KeywordTok{let} \ConstantTok{Some}\NormalTok{(}\OperatorTok{\&}\NormalTok{r) }\OperatorTok{=}\NormalTok{ risk}\OperatorTok{.}\NormalTok{get(}\OperatorTok{\&}\NormalTok{acct) }\OperatorTok{\{}
\NormalTok{            score }\OperatorTok{+=} \DecValTok{100.0} \OperatorTok{*}\NormalTok{ r}\OperatorTok{;}
            \ControlFlowTok{if}\NormalTok{ r }\OperatorTok{\textgreater{}}\NormalTok{ SEED\_PROXIMITY }\OperatorTok{\{}\NormalTok{ reasons}\OperatorTok{.}\NormalTok{push(}\StringTok{"near a watchlisted account"}\NormalTok{)}\OperatorTok{;} \OperatorTok{\}}
        \OperatorTok{\}}
        \CommentTok{// (2) in a laundering cycle}
        \ControlFlowTok{if}\NormalTok{ in\_cycle}\OperatorTok{.}\NormalTok{contains(}\OperatorTok{\&}\NormalTok{acct) }\OperatorTok{\{}
\NormalTok{            score }\OperatorTok{+=} \DecValTok{40.0}\OperatorTok{;}\NormalTok{ reasons}\OperatorTok{.}\NormalTok{push(}\StringTok{"member of a round{-}trip cycle"}\NormalTok{)}\OperatorTok{;}
        \OperatorTok{\}}
        \CommentTok{// (3) structuring fan{-}in}
        \ControlFlowTok{if}\NormalTok{ structuring\_suspects(sql}\OperatorTok{,} \OperatorTok{\&}\NormalTok{view}\OperatorTok{,}\NormalTok{ acct) }\OperatorTok{\{}
\NormalTok{            score }\OperatorTok{+=} \DecValTok{30.0}\OperatorTok{;}\NormalTok{ reasons}\OperatorTok{.}\NormalTok{push(}\StringTok{"structuring: high small{-}credit fan{-}in"}\NormalTok{)}\OperatorTok{;}
        \OperatorTok{\}}
        \CommentTok{// (4) fan{-}out distribution}
        \ControlFlowTok{if}\NormalTok{ view}\OperatorTok{.}\NormalTok{out\_degree(acct) }\OperatorTok{\textgreater{}=} \DecValTok{20} \OperatorTok{\{}
\NormalTok{            score }\OperatorTok{+=} \DecValTok{15.0}\OperatorTok{;}\NormalTok{ reasons}\OperatorTok{.}\NormalTok{push(}\StringTok{"fan{-}out distribution"}\NormalTok{)}\OperatorTok{;}
        \OperatorTok{\}}

        \ControlFlowTok{if}\NormalTok{ score }\OperatorTok{\textgreater{}=}\NormalTok{ CASE\_THRESHOLD }\OperatorTok{\{}
\NormalTok{            cases}\OperatorTok{.}\NormalTok{push(Case }\OperatorTok{\{}\NormalTok{ acct}\OperatorTok{,}\NormalTok{ score}\OperatorTok{,}\NormalTok{ reasons}\OperatorTok{,}\NormalTok{ snapshot}\OperatorTok{:}\NormalTok{ view}\OperatorTok{.}\NormalTok{node\_count() }\OperatorTok{\}}\NormalTok{)}\OperatorTok{;}
        \OperatorTok{\}}
    \OperatorTok{\}}
\NormalTok{    cases}\OperatorTok{.}\NormalTok{sort\_by(}\OperatorTok{|}\NormalTok{a}\OperatorTok{,}\NormalTok{ b}\OperatorTok{|}\NormalTok{ b}\OperatorTok{.}\NormalTok{score}\OperatorTok{.}\NormalTok{partial\_cmp(}\OperatorTok{\&}\NormalTok{a}\OperatorTok{.}\NormalTok{score)}\OperatorTok{.}\NormalTok{unwrap())}\OperatorTok{;}
    \ConstantTok{Ok}\NormalTok{(cases)}
\OperatorTok{\}}
\end{Highlighting}
\end{Shaded}

Because all four signals read the \textbf{same} \texttt{view} and the
\textbf{same} SQL snapshot, the score is internally consistent --- no
detector is looking at a different instant of the graph, which in a
fragmented stack is a constant source of false alerts.

\section{Case management and SAR
output}\label{case-management-and-sar-output}

A case above threshold is written back as a queryable table --- so
investigators triage with plain SQL, joined to the customer records,
over the same live data:

\begin{Shaded}
\begin{Highlighting}[]
\CommentTok{{-}{-} The investigator\textquotesingle{}s queue: highest{-}risk cases with the human context.}
\KeywordTok{SELECT}\NormalTok{ c.name, c.country, c.is\_pep, a.score, a.reasons}
\KeywordTok{FROM}\NormalTok{   alerts a}
\KeywordTok{JOIN}\NormalTok{   accounts   acct }\KeywordTok{ON}\NormalTok{ acct.}\KeywordTok{id} \OperatorTok{=}\NormalTok{ a.acct}
\KeywordTok{JOIN}\NormalTok{   customers  c    }\KeywordTok{ON}\NormalTok{ c.}\KeywordTok{id} \OperatorTok{=}\NormalTok{ acct.customer\_id}
\KeywordTok{WHERE}\NormalTok{  a.score }\OperatorTok{\textgreater{}=} \DecValTok{80}
\KeywordTok{ORDER}  \KeywordTok{BY}\NormalTok{ a.score }\KeywordTok{DESC}\NormalTok{;}
\end{Highlighting}
\end{Shaded}

Confirmed cases feed back into \texttt{known\_bad}, which
\textbf{re-seeds} the personalized PageRank on the next pass --- the
system learns: each confirmed mule sharpens the risk propagation for the
accounts around it.

\section{Operating it live}\label{operating-it-live}

\begin{Shaded}
\begin{Highlighting}[]
\NormalTok{sequenceDiagram}
\NormalTok{    participant TX as Transaction stream}
\NormalTok{    participant DB as ChakraDB}
\NormalTok{    participant SC as Scorer (every N seconds)}
\NormalTok{    participant INV as Investigator}
\NormalTok{    loop continuously}
\NormalTok{      TX{-}\textgreater{}\textgreater{}DB: ingest (txn + flow edge, one commit)}
\NormalTok{    end}
\NormalTok{    loop on a cadence}
\NormalTok{      SC{-}\textgreater{}\textgreater{}DB: view() — consistent snapshot; ingestion never pauses}
\NormalTok{      SC{-}\textgreater{}\textgreater{}SC: PPR + cycles + fan{-}in/out + velocity}
\NormalTok{      SC{-}\textgreater{}\textgreater{}DB: write alerts table}
\NormalTok{    end}
\NormalTok{    INV{-}\textgreater{}\textgreater{}DB: SQL over alerts JOIN customers (same live data)}
\end{Highlighting}
\end{Shaded}

\begin{itemize}
\tightlist
\item
  The \textbf{inline} checks (fan-in, velocity, watchlist membership)
  run per transaction using live adjacency (\texttt{in\_degree},
  \texttt{out\_neighbors}) --- cheap, no snapshot copy.
\item
  The \textbf{network} analysis (cycles, PageRank, personalized
  PageRank) runs on a cadence over a \texttt{view()}; each pass is a
  fresh consistent snapshot while ingestion continues at full rate.
\item
  Watch \texttt{Storage::stats()} for checkpoint lag and backpressure
  under peak load
  (\href{../getting-started/getting-started.md}{Observability}); back up
  the store on a schedule
  (\href{../getting-started/getting-started.md}{Backup \& Restore}).
\end{itemize}

\section{Scaling to millions per hour --- the event-driven
pipeline}\label{scaling-to-millions-per-hour-the-event-driven-pipeline}

The batch scorer above is correct but polls. A production AML system is
\emph{event-driven}: the instant a transaction commits, detection should
react --- and it must do so without ever slowing the ingest that feeds
it. This is where ChakraDB's design pays off, and it is worth building
out in full because it is the workload that sets ChakraDB apart.

\subsection{The trigger: a committed-change
stream}\label{the-trigger-a-committed-change-stream}

ChakraDB deliberately has \textbf{no in-SQL triggers}. Running user code
(especially Python under the GIL) inside the writer's critical section
would throttle the one thing that has to stay fast. Instead it exposes
the model an event-driven system actually wants: a \textbf{stream of
committed row changes} (change-data-capture).

A \texttt{CdcBackend} decorates the engine's backend and publishes a
\texttt{Change} \texttt{\{op,\ csn,\ old,\ new\}} after each
INSERT/UPDATE/DELETE commits --- after the write is applied and, on the
durable backend, after it is WAL-logged. An INSERT (the ingest hot path)
pays only a channel send; the old-row image for UPDATE/DELETE is read
outside any engine lock. Nothing touches the writer's hot locks.

In Rust, a subscription is a pull-based stream:

\begin{Shaded}
\begin{Highlighting}[]
\KeywordTok{let}\NormalTok{ cdc }\OperatorTok{=} \PreprocessorTok{Cdc::}\NormalTok{new()}\OperatorTok{;}
\KeywordTok{let}\NormalTok{ engine }\OperatorTok{=} \PreprocessorTok{SqlEngine::}\NormalTok{with\_backend(}\PreprocessorTok{CdcBackend::}\NormalTok{wrap(db}\OperatorTok{,}\NormalTok{ cdc}\OperatorTok{.}\NormalTok{clone()))}\OperatorTok{;}
\KeywordTok{let}\NormalTok{ stream }\OperatorTok{=}\NormalTok{ cdc}\OperatorTok{.}\NormalTok{subscribe(}\ConstantTok{Some}\NormalTok{(}\StringTok{"transactions"}\NormalTok{))}\OperatorTok{;}   \CommentTok{// one table, or None for all}

\CommentTok{// on the worker thread:}
\ControlFlowTok{while} \KeywordTok{let} \ConstantTok{Some}\NormalTok{(batch) }\OperatorTok{=}\NormalTok{ stream}\OperatorTok{.}\NormalTok{recv() }\OperatorTok{\{}              \CommentTok{// one commit\textquotesingle{}s changes}
    \ControlFlowTok{for}\NormalTok{ change }\KeywordTok{in} \OperatorTok{\&}\NormalTok{batch }\OperatorTok{\{}
        \ControlFlowTok{if}\NormalTok{ change}\OperatorTok{.}\NormalTok{op }\OperatorTok{==} \PreprocessorTok{ChangeOp::}\NormalTok{Insert }\OperatorTok{\{}
\NormalTok{            react(change}\OperatorTok{.}\NormalTok{new}\OperatorTok{.}\NormalTok{as\_ref()}\OperatorTok{.}\NormalTok{unwrap())}\OperatorTok{;}     \CommentTok{// fire detectors}
        \OperatorTok{\}}
    \OperatorTok{\}}
\OperatorTok{\}}
\end{Highlighting}
\end{Shaded}

In Python it is a callback --- the entire engine, including the graph
algorithms, in a few lines:

\begin{Shaded}
\begin{Highlighting}[]
\KeywordTok{def}\NormalTok{ react(old, new):}
    \ControlFlowTok{if}\NormalTok{ new:}
\NormalTok{        worker.on\_transaction(new)          }\CommentTok{\# dict: \{"src":…, "dst":…, "amount":…\}}

\NormalTok{conn.on\_change(}\StringTok{"transactions"}\NormalTok{, react)       }\CommentTok{\# fires after every committed write}
\end{Highlighting}
\end{Shaded}

Delivery is at-least-once and in commit order; a rolled-back transaction
never appears. A \texttt{ChangeSink} trait leaves room for a Kafka
transport (below) without changing a line of detector code.

\subsection{The three-tier detection
model}\label{the-three-tier-detection-model}

The trap every ``fraud on a graph'' system falls into is recomputing
global graph algorithms per event. At millions of transactions per hour
you cannot run PageRank on every insert. The pipeline splits detection
into three tiers by \emph{how much of the graph each needs} and
\emph{how fresh it must be}:

\begin{longtable}[]{@{}
  >{\raggedright\arraybackslash}p{(\linewidth - 8\tabcolsep) * \real{0.2000}}
  >{\raggedright\arraybackslash}p{(\linewidth - 8\tabcolsep) * \real{0.2000}}
  >{\raggedright\arraybackslash}p{(\linewidth - 8\tabcolsep) * \real{0.2000}}
  >{\raggedright\arraybackslash}p{(\linewidth - 8\tabcolsep) * \real{0.2000}}
  >{\raggedright\arraybackslash}p{(\linewidth - 8\tabcolsep) * \real{0.2000}}@{}}
\toprule\noalign{}
\begin{minipage}[b]{\linewidth}\raggedright
Tier
\end{minipage} & \begin{minipage}[b]{\linewidth}\raggedright
Runs
\end{minipage} & \begin{minipage}[b]{\linewidth}\raggedright
Scope
\end{minipage} & \begin{minipage}[b]{\linewidth}\raggedright
Cost
\end{minipage} & \begin{minipage}[b]{\linewidth}\raggedright
Catches
\end{minipage} \\
\midrule\noalign{}
\endhead
\bottomrule\noalign{}
\endlastfoot
\textbf{T0 --- per-event} & on every committed txn, at ingest rate &
O(1)--O(degree) & counter bump, set lookup & fan-in threshold crossed,
transacts-with-known-bad, fan-out \\
\textbf{T1 --- micro-batch} & every few seconds, on the touched subgraph
& k-hop neighbourhood of new edges & bounded & short laundering cycles
(bounded length + Δt), local dense mule clusters \\
\textbf{T2 --- periodic} & every few minutes, over a \texttt{view()}
snapshot & whole graph & seconds, \textbf{off the write path} & global
\texttt{personalized\_pagerank}, full
\texttt{strongly\_connected\_components}, communities \\
\end{longtable}

T0 keeps up with the firehose because it never touches more than one
node's neighbourhood. T2 is expensive but rare, and runs on an
\textbf{MVCC snapshot} --- so it never blocks ingest. That single fact
is the difference between ChakraDB and a bolt-on stack.

The worker's T0 and T2 in pseudocode:

\begin{quote}
\textbf{ALGORITHM --- T0 (per committed transaction)}

\begin{Shaded}
\begin{Highlighting}[]
\NormalTok{Input: change (src, dst, amount)}
\NormalTok{1  graph.add\_edge(src, dst, amount)                 ▷ mirror into payment graph}
\NormalTok{2  in\_sources[dst].add(src);  fan\_in ← |in\_sources[dst]|}
\NormalTok{3  if 9000 ≤ amount \textless{} 10000: near\_threshold[dst] += 1}
\NormalTok{4  out\_targets[src].add(dst); fan\_out ← |out\_targets[src]|}
\NormalTok{5  if fan\_in ≥ Θ\_in  and  near\_threshold[dst] ≥ Θ\_near:  alert(dst, STRUCTURING)}
\NormalTok{6  if fan\_out ≥ Θ\_out:                                    alert(src, FAN\_OUT)}
\NormalTok{7  if src ∈ known\_bad:                                    alert(dst, KNOWN\_BAD)}
\end{Highlighting}
\end{Shaded}
\end{quote}

\begin{quote}
\textbf{ALGORITHM --- T2 (every N events, over one snapshot)}

\begin{Shaded}
\begin{Highlighting}[]
\NormalTok{1  view ← graph.view()                               ▷ consistent MVCC snapshot}
\NormalTok{2  for ring in view.laundering\_cycles():             ▷ non{-}trivial SCCs}
\NormalTok{3      for m in ring: alert(m, LAYERING\_CYCLE)}
\NormalTok{4  risk ← view.personalized\_pagerank(known\_bad)       ▷ risk propagation}
\NormalTok{5  for acct in top\_k(risk): alert(acct, HIGH\_RISK\_EXPOSURE)}
\end{Highlighting}
\end{Shaded}
\end{quote}

\subsection{The architecture}\label{the-architecture}

\begin{center}\includegraphics[width=0.92\linewidth,height=0.80\textheight,keepaspectratio]{diagrams/diagram-33.pdf}\end{center}

\begin{itemize}
\tightlist
\item
  \textbf{Ingest} rides ChakraDB's group commit; the writer is never the
  bottleneck.
\item
  \textbf{The change stream} emits one batch per commit, CSN-stamped ---
  ordered, replayable, resumable by cursor.
\item
  \textbf{Transport is pluggable}: the in-process channel (default, µs
  latency, single node) or Kafka (partition key =
  \texttt{hash(account)}) so \emph{N} workers consume in parallel and
  survive restarts --- near-linear horizontal scale-out.
\item
  \textbf{T2 runs on a \texttt{view()}} --- a consistent snapshot --- so
  heavy analytics never block the writer. A Neo4j + Postgres + Flink
  stack must ETL OLTP → graph → OLAP and can never give you one
  transactionally-consistent view; ChakraDB does it in one process, no
  copy.
\end{itemize}

\subsection{Capacity: it is not close}\label{capacity-it-is-not-close}

The shipped in-process pipeline (\texttt{examples/aml\_pipeline.rs})
runs a generator thread and the reacting worker concurrently over one
engine. On a single machine it sustains, end to end --- ingest
\emph{plus} per-event detection \emph{plus} periodic global graph
passes:

\begin{Shaded}
\begin{Highlighting}[]
\NormalTok{Ingest: 60,037 transactions in 1.39s}
\NormalTok{        = 43,196 txn/s  ≈  155 million txn/hour}
\NormalTok{Worker: reacted to 60,037 committed transactions via the change stream}
\NormalTok{All typologies detected live — pipeline verified.}
\end{Highlighting}
\end{Shaded}

The standalone generator (\texttt{examples/aml\_gen.rs}) writes at
\textasciitilde190 million txn/hour. Against a target of ``millions per
hour,'' the engine has roughly two orders of magnitude of headroom ---
because readers never block writers, so the detection load and the
ingest load do not contend. This is the axis the project competes on,
made concrete.

\subsection{Horizontal scale-out}\label{horizontal-scale-out}

Beyond a single node, the Kafka sink partitions the change stream by
\texttt{hash(account)}. Each worker owns a shard of the payment graph
and its incremental T0 state; the known-bad seed set is replicated to
every shard, and cross-shard rings are escalated to a coordinator. The
CSN cursor per partition makes each worker resume exactly where it left
off after a restart --- no missed or double-counted transaction, which
is a hard compliance requirement.

\subsection{The two implementations}\label{the-two-implementations}

The complete pipeline ships in both languages:

\begin{itemize}
\tightlist
\item
  \textbf{\texttt{examples/aml\_pipeline.rs}} --- the scalable
  reference: a generator thread and a tiered worker over one
  \texttt{SqlEngine}, reacting through the Rust \texttt{ChangeStream}.
\item
  \textbf{\texttt{examples/aml\_stream.py}} --- the ergonomic mirror:
  the same tiers driven by \texttt{conn.on\_change}. (While a T2 pass
  holds the GIL, ingest pauses briefly --- which is precisely why the
  scalable worker is the Rust one.)
\item
  \textbf{\texttt{examples/aml\_gen.rs}} --- the standalone producer CLI
  (\texttt{-\/-db}, \texttt{-\/-count}, \texttt{-\/-seed}), the feed for
  the Kafka-partitioned topology.
\end{itemize}

\section{What ChakraDB made possible}\label{what-chakradb-made-possible}

\begin{longtable}[]{@{}ll@{}}
\toprule\noalign{}
Laundering typology & ChakraDB primitive \\
\midrule\noalign{}
\endhead
\bottomrule\noalign{}
\endlastfoot
Structuring / smurfing & \texttt{in\_degree} (fan-in) + SQL velocity \\
Round-tripping / layering & \texttt{laundering\_cycles} (SCC) \\
Mule networks & \texttt{connected\_components} + \texttt{pagerank} \\
Proximity to known-bad & \texttt{personalized\_pagerank(seeds)} \\
Rapid movement & temporal SQL \\
Fan-out distribution & \texttt{out\_degree} \\
\end{longtable}

Every one of these ran over the \textbf{same live snapshot}, in
\textbf{one embedded process}, while transactions kept arriving --- the
combination the fragmented AML stack cannot offer. The graph algorithms
the system needed (reverse adjacency for fan-in, SCC for cycles,
personalized PageRank for risk propagation) are built into ChakraDB's
core (\texttt{src/graph.rs}), so the application is a few hundred lines
of scoring logic, not a three-system integration project. \textbf{Follow
the money --- in real time, over one consistent view of it.}

\section{The reference
implementation}\label{the-reference-implementation}

This case study ships as two runnable programs --- the same system in
each language, generating their own synthetic data (legitimate traffic
with laundering rings deliberately injected) and \emph{asserting} that
every planted typology is caught:

\begin{Shaded}
\begin{Highlighting}[]
\ExtensionTok{cargo}\NormalTok{ run }\AttributeTok{{-}{-}release} \AttributeTok{{-}{-}example}\NormalTok{ aml\_realtime }\AttributeTok{{-}{-}no{-}default{-}features}   \CommentTok{\# Rust}
\ExtensionTok{python}\NormalTok{ examples/aml\_app.py                                          }\CommentTok{\# Python}
\end{Highlighting}
\end{Shaded}

Both build a \textasciitilde230-account payment network, run the five
detectors over one \texttt{view()}, and emit a ranked SAR list. The
output speaks for itself --- signal cleanly separated from a sea of
legitimate activity, with zero false positives:

\begin{Shaded}
\begin{Highlighting}[]
\NormalTok{── Detector 1 — Structuring (fan{-}in of near{-}threshold deposits) ──}
\NormalTok{  account  500: 12 distinct sources, 12 just{-}under{-}threshold deposits, $114,669 received}

\NormalTok{── Detector 2 — Layering (round{-}trip cycles / SCCs) ──}
\NormalTok{  ring \#1: [700, 701, 702, 703]  — funds return to origin}

\NormalTok{── Detector 3 — Mule fan{-}out (distribution hubs) ──}
\NormalTok{  account  800: pays out to 20 distinct counterparties}

\NormalTok{── Detector 4 — Risk propagation (personalized PageRank from known{-}bad) ──}
\NormalTok{  seed(s) [800] → risk reached 66 downstream accounts.}

\NormalTok{── SAR candidates — fused risk ranking ──}
\NormalTok{   acct   score  reasons}
\NormalTok{    800    6.00  known{-}bad, mule{-}fan{-}out}
\NormalTok{    703    3.00  laundering{-}cycle}
\NormalTok{    700    3.00  laundering{-}cycle}
\NormalTok{    702    3.00  laundering{-}cycle}
\NormalTok{    701    3.00  laundering{-}cycle}
\NormalTok{    500    2.50  structuring{-}collector}
\end{Highlighting}
\end{Shaded}

The Rust program is the canonical reference; the Python program is the
same logic through \texttt{conn.graph(...)}, showing that the
\emph{entire} engine --- SQL, exact money, and the graph algorithms ---
is available to a client in a few hundred lines. \textbf{Follow the
money --- in real time, over one consistent view of it.}

\chapter{Case Study: Real-Time Counterparty Credit \& Market
Risk}\label{case-study-real-time-counterparty-credit-market-risk}

\epigraph{The revolutionary idea that defines the boundary between modern times and the past is the mastery of risk.}{--- Peter L. Bernstein, \emph{Against the Gods}}

Banks still compute much of their risk \textbf{overnight}, in batch: the
book is frozen, market data is snapshotted, and a grid churns until
morning. By the time the numbers land, the market has moved. The prize
everyone chases is \textbf{intraday, real-time risk} --- exposure,
limits, VaR, and systemic contagion recomputed as prices tick and trades
book. That is precisely the shape of workload ChakraDB is built for, and
this chapter builds it end to end.

It is the risk analogue of the \href{aml.md}{AML case study}: a live
portfolio and a streaming market-data feed drive a \textbf{materialized
worker} that reacts to every price tick through the change stream, over
one consistent dataset, never blocking ingest. The runnable code is
\texttt{examples/ccr\_pipeline.rs} (Rust) and
\texttt{examples/ccr\_stream.py} (Python).

\section{Two risk domains, one
engine}\label{two-risk-domains-one-engine}

\begin{longtable}[]{@{}
  >{\raggedright\arraybackslash}p{(\linewidth - 4\tabcolsep) * \real{0.3333}}
  >{\raggedright\arraybackslash}p{(\linewidth - 4\tabcolsep) * \real{0.3333}}
  >{\raggedright\arraybackslash}p{(\linewidth - 4\tabcolsep) * \real{0.3333}}@{}}
\toprule\noalign{}
\begin{minipage}[b]{\linewidth}\raggedright
Domain
\end{minipage} & \begin{minipage}[b]{\linewidth}\raggedright
Question
\end{minipage} & \begin{minipage}[b]{\linewidth}\raggedright
Cadence
\end{minipage} \\
\midrule\noalign{}
\endhead
\bottomrule\noalign{}
\endlastfoot
\textbf{Market risk} & How much could the book lose from price moves?
(MtM, VaR) & per tick + periodic \\
\textbf{Counterparty credit risk (CCR)} & If a counterparty defaults,
what do we lose --- and who else falls? & per trade + periodic \\
\end{longtable}

Both are functions of the \emph{same} live data --- the trades, the
prices, and the web of who owes whom. Stitching them across an OLTP
store, a market-data bus, a graph database, and an overnight risk grid
is the status quo; doing them together, in one embedded engine, on one
consistent snapshot, is the ChakraDB thesis.

\section{The exposure network is a
graph}\label{the-exposure-network-is-a-graph}

CCR is intrinsically a \textbf{network} problem, and this is where the
built-in graph engine becomes the risk engine. Model the interbank
system as a directed, weighted graph: a node is a legal entity; an edge
\texttt{i\ →\ j} of weight \texttt{w} means \emph{\texttt{i} owes
\texttt{j} the amount \texttt{w}}. On that one structure:

\begin{longtable}[]{@{}
  >{\raggedright\arraybackslash}p{(\linewidth - 2\tabcolsep) * \real{0.5000}}
  >{\raggedright\arraybackslash}p{(\linewidth - 2\tabcolsep) * \real{0.5000}}@{}}
\toprule\noalign{}
\begin{minipage}[b]{\linewidth}\raggedright
Risk question
\end{minipage} & \begin{minipage}[b]{\linewidth}\raggedright
Built-in primitive
\end{minipage} \\
\midrule\noalign{}
\endhead
\bottomrule\noalign{}
\endlastfoot
Who is systemically important (too-connected-to-fail)? &
\texttt{pagerank} on the liability network \\
If X defaults, who is dragged under, and by how much? &
\textbf{\texttt{eisenberg\_noe}} (clearing vector / default cascade) \\
Hidden circular exposures = netting opportunities &
\texttt{laundering\_cycles} (SCCs) \\
Netting sets / collateral pools & \texttt{connected\_components} \\
Concentration / Cover-2 & top-k weighted degree \\
\end{longtable}

\begin{center}\includegraphics[width=0.92\linewidth,height=0.80\textheight,keepaspectratio]{diagrams/diagram-34.pdf}\end{center}

\section{The schema}\label{the-schema-1}

Trades and market data are ordinary tables; the exposure network is a
graph over the same database.

\begin{Shaded}
\begin{Highlighting}[]
\KeywordTok{CREATE} \KeywordTok{TABLE}\NormalTok{ counterparties (}\KeywordTok{id} \DataTypeTok{INTEGER} \KeywordTok{PRIMARY} \KeywordTok{KEY}\NormalTok{, name }\DataTypeTok{VARCHAR}\NormalTok{(}\DecValTok{32}\NormalTok{),}
\NormalTok{    rating }\DataTypeTok{VARCHAR}\NormalTok{(}\DecValTok{4}\NormalTok{), external\_assets }\DataTypeTok{DECIMAL}\NormalTok{(}\DecValTok{18}\NormalTok{,}\DecValTok{2}\NormalTok{));}
\KeywordTok{CREATE} \KeywordTok{TABLE}\NormalTok{ trades (trade\_id }\DataTypeTok{INTEGER} \KeywordTok{PRIMARY} \KeywordTok{KEY}\NormalTok{, counterparty }\DataTypeTok{INTEGER}\NormalTok{,}
\NormalTok{    instrument }\DataTypeTok{INTEGER}\NormalTok{, notional }\DataTypeTok{DECIMAL}\NormalTok{(}\DecValTok{18}\NormalTok{,}\DecValTok{2}\NormalTok{), trade\_price }\DataTypeTok{DECIMAL}\NormalTok{(}\DecValTok{12}\NormalTok{,}\DecValTok{4}\NormalTok{));}
\KeywordTok{CREATE} \KeywordTok{TABLE}\NormalTok{ market\_data (tick\_id }\DataTypeTok{INTEGER} \KeywordTok{PRIMARY} \KeywordTok{KEY}\NormalTok{, instrument }\DataTypeTok{INTEGER}\NormalTok{,}
\NormalTok{    price }\DataTypeTok{DECIMAL}\NormalTok{(}\DecValTok{12}\NormalTok{,}\DecValTok{4}\NormalTok{), ts }\DataTypeTok{TIMESTAMP}\NormalTok{);              }\CommentTok{{-}{-} the streaming feed}
\KeywordTok{CREATE} \KeywordTok{TABLE}\NormalTok{ limits (counterparty }\DataTypeTok{INTEGER} \KeywordTok{PRIMARY} \KeywordTok{KEY}\NormalTok{, limit\_amount }\DataTypeTok{DECIMAL}\NormalTok{(}\DecValTok{18}\NormalTok{,}\DecValTok{2}\NormalTok{));}
\KeywordTok{CREATE} \KeywordTok{TABLE}\NormalTok{ interbank (}\KeywordTok{id} \DataTypeTok{INTEGER} \KeywordTok{PRIMARY} \KeywordTok{KEY}\NormalTok{, debtor }\DataTypeTok{INTEGER}\NormalTok{, creditor }\DataTypeTok{INTEGER}\NormalTok{,}
\NormalTok{    amount }\DataTypeTok{DECIMAL}\NormalTok{(}\DecValTok{18}\NormalTok{,}\DecValTok{2}\NormalTok{));                            }\CommentTok{{-}{-} edges of the exposure graph}
\KeywordTok{CREATE} \KeywordTok{TABLE}\NormalTok{ risk\_alerts (}\KeywordTok{id} \DataTypeTok{INTEGER} \KeywordTok{PRIMARY} \KeywordTok{KEY}\NormalTok{, entity }\DataTypeTok{INTEGER}\NormalTok{, kind }\DataTypeTok{VARCHAR}\NormalTok{(}\DecValTok{40}\NormalTok{));}
\end{Highlighting}
\end{Shaded}

Money is exact \texttt{DECIMAL}, never binary float --- a hard
requirement when the number is a P\&L or an exposure
(\href{../engine/data-types.md}{Exact Decimal Arithmetic}).

\section{Real-time exposure: the per-tick
calculation}\label{real-time-exposure-the-per-tick-calculation}

\textbf{Current exposure} to a counterparty is the netted, non-negative
mark-to-market of the trades facing it:
\texttt{CE\_c\ =\ max(Σ\_\{t∈c\}\ MtM\_t,\ 0)}. For a linear instrument
\texttt{MtM\_t\ =\ notional\_t\ ·\ (price\ −\ trade\_price\_t)}. The
naïve approach reprices the whole book on every tick; the scalable
approach reprices \textbf{only the trades on the instrument that
ticked}, and folds the delta into a running per-counterparty exposure
--- O(trades-on-that-instrument), not O(book).

\begin{quote}
\textbf{ALGORITHM --- T0: on each committed price tick
\texttt{(instrument\ I,\ price\ p)}}

\begin{Shaded}
\begin{Highlighting}[]
\NormalTok{1  Δ ← p − last\_price[I]}
\NormalTok{2  pnl\_window.push( Δ · Σ\_\{t on I\} notional\_t )        ▷ portfolio P\&L for VaR}
\NormalTok{3  for each trade t on I:}
\NormalTok{4      mtm′ ← notional\_t · (p − trade\_price\_t)}
\NormalTok{5      exposure[t.cp] += mtm′ − mtm[t];  mtm[t] ← mtm′  ▷ incremental netting}
\NormalTok{6  for each affected counterparty c:}
\NormalTok{7      if max(exposure[c], 0) \textgreater{} limit[c]:  alert(c, LIMIT\_BREACH)}
\NormalTok{8  last\_price[I] ← p}
\end{Highlighting}
\end{Shaded}
\end{quote}

This runs on every committed tick, at ingest speed --- the change stream
delivers the tick the instant it commits.

\section{Portfolio \& systemic risk: the periodic
pass}\label{portfolio-systemic-risk-the-periodic-pass}

The heavy analytics run on a cadence over a consistent
\texttt{Graph::view()} snapshot, off the write path.

\textbf{Value-at-Risk.} A 99\% one-tick VaR by historical simulation:
the worst 1\% loss in the rolling P\&L window. A breach of the VaR limit
raises a portfolio alert.

\textbf{Systemic importance.} \texttt{pagerank} over the liability
network. Because edges run debtor → creditor, PageRank mass accumulates
at the entities everyone owes --- the hubs whose failure would ripple
furthest.

\textbf{Default contagion --- the Eisenberg--Noe clearing vector.} The
centerpiece. Given external assets \texttt{e\_i} (each entity's outside
cash) and the liability edges, the clearing payment vector \texttt{p}
solves

\[
  p_i = \min\!\Big( \bar p_i,\; e_i + \sum_j \Pi_{ji}\, p_j \Big),
  \qquad \bar p_i = \sum_j L_{ij},\quad \Pi_{ji} = L_{ji}/\bar p_j .
\]

Each node pays the smaller of what it owes and what it actually has once
upstream payments arrive. ChakraDB solves it by the monotone Picard
iteration from \texttt{p\ =\ p̄} downward --- which converges to the
greatest clearing vector --- and reports each entity's payment,
surviving equity, and, crucially, the set of \textbf{defaulters}:
exactly who a cascade sinks. Feed it a stressed
\texttt{external\_assets} (one entity's cash set to zero) and it
simulates the shock's contagion.

\begin{quote}
\textbf{ALGORITHM --- T2: periodic systemic pass over one snapshot}

\begin{Shaded}
\begin{Highlighting}[]
\NormalTok{1  view ← exposures.view()}
\NormalTok{2  if VaR₉₉(pnl\_window) \textgreater{} var\_limit:  alert(PORTFOLIO, VAR\_BREACH)}
\NormalTok{3  hub ← argmax view.pagerank();       alert(hub, SYSTEMICALLY\_IMPORTANT)}
\NormalTok{4  for d in view.eisenberg\_noe(external\_assets).defaulted:}
\NormalTok{5      alert(d, DEFAULT\_CASCADE)}
\NormalTok{6  for ring in view.laundering\_cycles():}
\NormalTok{7      for m in ring: alert(m, CIRCULAR\_EXPOSURE)}
\end{Highlighting}
\end{Shaded}
\end{quote}

\section{The worker is a materialized
worker}\label{the-worker-is-a-materialized-worker}

The whole risk engine is one \href{../engine/overview.md}{materialized
worker} --- a named, incrementally-maintained derivation of the data. In
Rust it implements \texttt{MaterializedWorker} and is registered on the
change stream:

\begin{Shaded}
\begin{Highlighting}[]
\KeywordTok{let}\NormalTok{ ccr }\OperatorTok{=}\NormalTok{ cdc}\OperatorTok{.}\NormalTok{materialize(}\ConstantTok{Some}\NormalTok{(}\StringTok{"market\_data"}\NormalTok{)}\OperatorTok{,} \PreprocessorTok{RiskWorker::}\NormalTok{new(engine}\OperatorTok{,}\NormalTok{ exposures}\OperatorTok{,}\NormalTok{ trades}\OperatorTok{,}\NormalTok{ assets))}\OperatorTok{;}
\CommentTok{// ... the feed streams; the worker maintains exposure and fires alerts ...}
\NormalTok{ccr}\OperatorTok{.}\NormalTok{query(}\OperatorTok{|}\NormalTok{w}\OperatorTok{|}\NormalTok{ w}\OperatorTok{.}\NormalTok{cp\_exposure}\OperatorTok{.}\NormalTok{clone())}\OperatorTok{;}   \CommentTok{// real{-}time exposure, any time}
\NormalTok{ccr}\OperatorTok{.}\NormalTok{stop()}\OperatorTok{;}                             \CommentTok{// client owns the lifecycle}
\end{Highlighting}
\end{Shaded}

In Python it is a class plus \texttt{conn.on\_change} --- the same
design, the full engine (SQL, exact money, the graph algorithms,
\texttt{eisenberg\_noe}) in a few hundred lines:

\begin{Shaded}
\begin{Highlighting}[]
\KeywordTok{def}\NormalTok{ react(old, new):}
\NormalTok{    worker.on\_change(old, new)          }\CommentTok{\# dict: \{"instrument":…, "price":…\}}
\NormalTok{conn.on\_change(}\StringTok{"market\_data"}\NormalTok{, react)}
\end{Highlighting}
\end{Shaded}

\section{Capacity}\label{capacity}

The shipped pipeline (\texttt{examples/ccr\_pipeline.rs}) runs the feed
and the reacting risk worker concurrently over one engine:

\begin{Shaded}
\begin{Highlighting}[]
\NormalTok{Ingest: 40,000 market{-}data ticks in 0.84s}
\NormalTok{        = 47,790 ticks/s  ≈  172 million/hour}
\NormalTok{All injected risks detected live — CCR pipeline verified.}
\end{Highlighting}
\end{Shaded}

Real-time exposure, single-name limits, VaR, systemic-importance,
default-cascade contagion, and circular-exposure detection --- all live,
over one consistent dataset, while the feed keeps ticking. The periodic
Eisenberg--Noe and PageRank passes run on an MVCC snapshot, so they
never stall ingest. That non-blocking split is the difference between
real-time risk and an overnight batch.

\section{What ChakraDB made
possible}\label{what-chakradb-made-possible-1}

\begin{longtable}[]{@{}
  >{\raggedright\arraybackslash}p{(\linewidth - 2\tabcolsep) * \real{0.5000}}
  >{\raggedright\arraybackslash}p{(\linewidth - 2\tabcolsep) * \real{0.5000}}@{}}
\toprule\noalign{}
\begin{minipage}[b]{\linewidth}\raggedright
Risk measure
\end{minipage} & \begin{minipage}[b]{\linewidth}\raggedright
ChakraDB primitive
\end{minipage} \\
\midrule\noalign{}
\endhead
\bottomrule\noalign{}
\endlastfoot
Real-time current exposure \& netting & incremental per-tick repricing
(T0) \\
Single-name limit monitoring & live exposure vs \texttt{limits} \\
Portfolio VaR / Expected Shortfall & historical simulation over the P\&L
window \\
Systemic importance & \texttt{pagerank} on the exposure network \\
Default-cascade contagion & \textbf{\texttt{eisenberg\_noe}} clearing
vector \\
Circular exposures / netting & \texttt{laundering\_cycles} (SCC) \\
Netting sets & \texttt{connected\_components} \\
\end{longtable}

The injected risks --- a concentrated single-name breach, a portfolio
VaR breach, a systemic hub, a thinly-capitalised default cascade, and a
ring of circular exposures --- are every one detected \textbf{live}, in
one embedded process, while the market-data feed streams at 170 million
ticks/hour. \textbf{Risk, in real time, over one consistent view of the
book.}

\chapter{Case Study: Real-Time
Recommendations}\label{case-study-real-time-recommendations}

\epigraph{People who liked this also liked \ldots}{--- every store, everywhere}

The first two case studies catch \emph{bad} behaviour. This one is the
opposite: it helps, and it shows ChakraDB's reach beyond fraud and risk.
A stream of user interactions --- views, clicks, purchases --- drives a
\textbf{live recommendation engine}: as engagement arrives, ``what
should this user see next?'' is answered in real time, over one
consistent graph, never blocking the ingest that feeds it.

The workload usually needs a separate feature store, a nightly
model-training job, and a serving tier reconciled by pipelines. Here it
is one embedded process reacting to the interaction stream. The runnable
code is \texttt{examples/reco\_pipeline.rs} (Rust) and
\texttt{examples/reco\_stream.py} (Python).

\section{Recommendation is link prediction on a bipartite
graph}\label{recommendation-is-link-prediction-on-a-bipartite-graph}

Model engagement as a bipartite graph: users on one side, items on the
other, an edge for each interaction. ``What should user \emph{u} see?''
becomes \emph{link prediction} --- which edges are likely to form next
--- and ChakraDB has the primitives built in:

\begin{longtable}[]{@{}
  >{\raggedright\arraybackslash}p{(\linewidth - 2\tabcolsep) * \real{0.5000}}
  >{\raggedright\arraybackslash}p{(\linewidth - 2\tabcolsep) * \real{0.5000}}@{}}
\toprule\noalign{}
\begin{minipage}[b]{\linewidth}\raggedright
Question
\end{minipage} & \begin{minipage}[b]{\linewidth}\raggedright
Built-in primitive
\end{minipage} \\
\midrule\noalign{}
\endhead
\bottomrule\noalign{}
\endlastfoot
What should this user see next? & \texttt{recommend(user,\ k)} ---
random-walk-with-restart \\
You might also like \ldots{} & \texttt{adamic\_adar(a,\ b)} ---
rarity-weighted item similarity \\
What's trending? & \texttt{out\_degree} / \texttt{pagerank} \\
Users like this one & \texttt{jaccard\_similarity},
\texttt{common\_neighbors} \\
\end{longtable}

Edges are added \textbf{both ways} (\texttt{u\ →\ i} and
\texttt{i\ →\ u}) so a random walk can bounce user → item → user → item,
which is exactly collaborative filtering: the walk reaches items engaged
by users similar to you.

\begin{center}\includegraphics[width=0.92\linewidth,height=0.80\textheight,keepaspectratio]{diagrams/diagram-35.pdf}\end{center}

The target user engaged A and B; a cohort engaged A, B \textbf{and C};
the walk from the target reaches C through the cohort --- so C is
recommended.

\section{The schema}\label{the-schema-2}

\begin{Shaded}
\begin{Highlighting}[]
\KeywordTok{CREATE} \KeywordTok{TABLE}\NormalTok{ interactions (}\KeywordTok{id} \DataTypeTok{INTEGER} \KeywordTok{PRIMARY} \KeywordTok{KEY}\NormalTok{, user\_id }\DataTypeTok{INTEGER}\NormalTok{,}
\NormalTok{    item\_id }\DataTypeTok{INTEGER}\NormalTok{, kind }\DataTypeTok{VARCHAR}\NormalTok{(}\DecValTok{8}\NormalTok{), ts }\DataTypeTok{TIMESTAMP}\NormalTok{);   }\CommentTok{{-}{-} the engagement stream}
\KeywordTok{CREATE} \KeywordTok{TABLE}\NormalTok{ recommendations (}\KeywordTok{id} \DataTypeTok{INTEGER} \KeywordTok{PRIMARY} \KeywordTok{KEY}\NormalTok{, user\_id }\DataTypeTok{INTEGER}\NormalTok{, item\_id }\DataTypeTok{INTEGER}\NormalTok{);}
\end{Highlighting}
\end{Shaded}

The engagement \textbf{graph} is derived from the stream; item nodes
live above the user-id range so the two never collide.

\section{The recommender is a materialized
worker}\label{the-recommender-is-a-materialized-worker}

The engine is one \href{../reactive/change-streams.md}{materialized
worker} over the interaction stream. Each interaction folds into the
graph (T0); recommendations refresh periodically over a snapshot (T2).

\begin{quote}
\textbf{ALGORITHM --- T0: on each committed interaction
\texttt{(user\ u,\ item\ i)}}

\begin{Shaded}
\begin{Highlighting}[]
\NormalTok{1  graph.add\_edge(u, item\_node(i));  graph.add\_edge(item\_node(i), u)   ▷ both ways}
\end{Highlighting}
\end{Shaded}
\end{quote}

\begin{quote}
\textbf{ALGORITHM --- T2: refresh recommendations over one snapshot}

\begin{Shaded}
\begin{Highlighting}[]
\NormalTok{1  view ← graph.view()}
\NormalTok{2  recs    ← view.recommend(u, k)            ▷ items reached via similar users, unseen}
\NormalTok{3  trending← items sorted by view.out\_degree ▷ most distinct users}
\NormalTok{4  similar ← view.adamic\_adar(anchor, ·)     ▷ "you might also like"}
\end{Highlighting}
\end{Shaded}
\end{quote}

\texttt{recommend} runs personalized PageRank seeded at the user, drops
the user's already-engaged items and any zero-signal candidate, and
returns the top-\texttt{k} remainder --- collaborative filtering
expressed as a graph walk, in one call.

\section{In code}\label{in-code}

Rust --- register the recommender on the interaction stream:

\begin{Shaded}
\begin{Highlighting}[]
\KeywordTok{let}\NormalTok{ engagement }\OperatorTok{=} \PreprocessorTok{Graph::}\NormalTok{open(db}\OperatorTok{.}\NormalTok{clone()}\OperatorTok{,} \StringTok{"engagement"}\NormalTok{)}\OperatorTok{?;}
\KeywordTok{let}\NormalTok{ reco }\OperatorTok{=}\NormalTok{ cdc}\OperatorTok{.}\NormalTok{register(}\StringTok{"recommender"}\OperatorTok{,} \ConstantTok{Some}\NormalTok{(}\StringTok{"interactions"}\NormalTok{)}\OperatorTok{,}
                        \PreprocessorTok{Recommender::}\NormalTok{new(engine}\OperatorTok{.}\NormalTok{clone()}\OperatorTok{,}\NormalTok{ engagement))}\OperatorTok{;}
\CommentTok{// … the feed streams; recommendations refresh incrementally …}
\NormalTok{reco}\OperatorTok{.}\NormalTok{query(}\OperatorTok{|}\NormalTok{w}\OperatorTok{|}\NormalTok{ w}\OperatorTok{.}\NormalTok{last\_recommendations}\OperatorTok{.}\NormalTok{clone())}\OperatorTok{;}
\end{Highlighting}
\end{Shaded}

Python --- the same worker as a class plus \texttt{conn.on\_change}:

\begin{Shaded}
\begin{Highlighting}[]
\NormalTok{view }\OperatorTok{=}\NormalTok{ graph.view()}
\NormalTok{recs   }\OperatorTok{=}\NormalTok{ view.recommend(target\_user, }\DecValTok{10}\NormalTok{)      }\CommentTok{\# [(item\_node, score), …]}
\NormalTok{trend  }\OperatorTok{=} \BuiltInTok{sorted}\NormalTok{(items, key}\OperatorTok{=}\NormalTok{view.out\_degree, reverse}\OperatorTok{=}\VariableTok{True}\NormalTok{)}
\NormalTok{also   }\OperatorTok{=}\NormalTok{ view.adamic\_adar(item\_a, item\_b)}
\end{Highlighting}
\end{Shaded}

\section{Capacity}\label{capacity-1}

\begin{Shaded}
\begin{Highlighting}[]
\NormalTok{Ingest: 50,000 interactions in 1.1s}
\NormalTok{        = 45,899 interactions/s  ≈  165 million/hour}
\NormalTok{Recommendations verified — the cohort\textquotesingle{}s item surfaced live.}
\end{Highlighting}
\end{Shaded}

Recommendations stay fresh as engagement streams, because the heavy walk
runs on an MVCC snapshot off the write path. A cohort's co-engaged item
is surfaced to the target user; the trending item ranks first by
engagement; Adamic--Adar ranks the bundle items as most similar ---
every check passing on the live stream.

\section{The pattern, three times
over}\label{the-pattern-three-times-over}

This is the third system built on the same machinery --- the change
stream, a materialized worker, and the graph library --- as
\href{aml.md}{AML} and \href{ccr.md}{Counterparty Risk}. Catch
laundering, measure systemic risk, or personalize a feed: the engine is
the same, the workload reacts to live data in one process, and the
reader never blocks the writer. That is what makes ChakraDB more than a
query engine --- it is a substrate for real-time systems.

\chapter{Case Study: Real-Time Fraud Detection (HTAP +
Graph)}\label{case-study-real-time-fraud-detection-htap-graph}

\epigraph{Trust, but verify.}{--- Russian proverb}

This case study exercises everything at once: transactional ingest,
analytical scoring, and graph traversal --- over \textbf{one consistent
snapshot}, in one embedded process. It is the workload ChakraDB is built
for.

\section{The problem}\label{the-problem-1}

A payments service must, for each incoming transaction, decide
\emph{in-line} whether it looks fraudulent. Fraud signals are both
\textbf{statistical} (this card's spend rate just spiked) and
\textbf{relational} (this device has touched five accounts that share a
shipping address with a known mule). Traditionally that means two or
three systems: an OLTP store for the writes, a warehouse for the stats,
and a graph database for the relationships --- kept in sync by ETL,
always a little stale.

ChakraDB collapses them.

\begin{center}\includegraphics[width=0.92\linewidth,height=0.80\textheight,keepaspectratio]{diagrams/diagram-36.pdf}\end{center}

\section{The schema}\label{the-schema-3}

\begin{Shaded}
\begin{Highlighting}[]
\CommentTok{{-}{-} The transactional ledger (exact money — no rounding).}
\KeywordTok{CREATE} \KeywordTok{TABLE}\NormalTok{ txns (}
  \KeywordTok{id}        \DataTypeTok{INT} \KeywordTok{PRIMARY} \KeywordTok{KEY}\NormalTok{,}
\NormalTok{  card\_id   }\DataTypeTok{INT} \KeywordTok{NOT} \KeywordTok{NULL}\NormalTok{,}
\NormalTok{  device\_id }\DataTypeTok{INT} \KeywordTok{NOT} \KeywordTok{NULL}\NormalTok{,}
\NormalTok{  amount    }\DataTypeTok{DECIMAL}\NormalTok{(}\DecValTok{12}\NormalTok{,}\DecValTok{2}\NormalTok{) }\KeywordTok{NOT} \KeywordTok{NULL} \KeywordTok{CHECK}\NormalTok{ (amount }\OperatorTok{\textgreater{}=} \DecValTok{0}\NormalTok{),}
\NormalTok{  ts        }\DataTypeTok{TIMESTAMP} \KeywordTok{NOT} \KeywordTok{NULL}\NormalTok{,}
\NormalTok{  status    }\DataTypeTok{VARCHAR}\NormalTok{(}\DecValTok{8}\NormalTok{) }\KeywordTok{DEFAULT} \StringTok{\textquotesingle{}pending\textquotesingle{}}
\NormalTok{);}

\CommentTok{{-}{-} The entity graph: cards, devices, addresses linked by shared use.}
\CommentTok{{-}{-} (Managed by the Graph handle; edges keyed (src,dst) src{-}major.)}
\end{Highlighting}
\end{Shaded}

Nodes are entities (cards, devices, addresses, accounts); an edge means
``these two entities were seen together.'' A fraud \emph{ring} is a
dense cluster in this graph.

\section{Ingest: the transactional
write}\label{ingest-the-transactional-write}

Each transaction is one ACID transaction --- the ledger row and the
entity edges commit together, so no reader ever sees half of it:

\begin{Shaded}
\begin{Highlighting}[]
\ControlFlowTok{BEGIN}\NormalTok{;}
  \KeywordTok{INSERT} \KeywordTok{INTO}\NormalTok{ txns }\KeywordTok{VALUES}\NormalTok{ (}\DecValTok{9001}\NormalTok{, }\DecValTok{55}\NormalTok{, }\DecValTok{88}\NormalTok{, }\FloatTok{249.99}\NormalTok{, }\StringTok{\textquotesingle{}2026{-}07{-}22 14:03:01\textquotesingle{}}\NormalTok{, }\StringTok{\textquotesingle{}pending\textquotesingle{}}\NormalTok{);}
  \CommentTok{{-}{-} link card 55 \textless{}{-}\textgreater{} device 88 in the entity graph (via the Graph handle)}
\KeywordTok{COMMIT}\NormalTok{;                       }\CommentTok{{-}{-} first{-}committer{-}wins; crash{-}atomic}
\end{Highlighting}
\end{Shaded}

\texttt{DECIMAL(12,2)} keeps money exact;
\texttt{CHECK}/\texttt{NOT\ NULL} reject bad rows before they touch the
log.

\section{Scoring: analytical + graph, one
snapshot}\label{scoring-analytical-graph-one-snapshot}

The scorer takes a single snapshot and reads all three signals from it
--- no cross-system consistency to reason about:

\begin{Shaded}
\begin{Highlighting}[]
\KeywordTok{let}\NormalTok{ snap\_engine }\OperatorTok{=} \OperatorTok{\&}\NormalTok{engine}\OperatorTok{;}                 \CommentTok{// SQL over the current snapshot}
\KeywordTok{let}\NormalTok{ graph\_view  }\OperatorTok{=}\NormalTok{ fraud\_graph}\OperatorTok{.}\NormalTok{view()}\OperatorTok{?;}     \CommentTok{// CSR over the SAME instant}

\CommentTok{// (A) Statistical: this card\textquotesingle{}s velocity in the last minute.}
\KeywordTok{let}\NormalTok{ velocity}\OperatorTok{:} \DataTypeTok{i64} \OperatorTok{=}\NormalTok{ snap\_engine}
    \OperatorTok{.}\NormalTok{query(}\StringTok{"SELECT COUNT(*) FROM txns }\SpecialCharTok{\textbackslash{}}
\StringTok{            WHERE card\_id = 55 AND ts \textgreater{} \textquotesingle{}2026{-}07{-}22 14:02:01\textquotesingle{}"}\NormalTok{)}\OperatorTok{?}
\NormalTok{    [}\DecValTok{0}\NormalTok{][}\DecValTok{0}\NormalTok{]}\OperatorTok{.}\NormalTok{parse()}\OperatorTok{?;}

\CommentTok{// (G) Relational: how tightly is this device tied into a cluster?}
\KeywordTok{let}\NormalTok{ ring\_size }\OperatorTok{=}\NormalTok{ graph\_view}
    \OperatorTok{.}\NormalTok{connected\_components\_of(device\_node)   }\CommentTok{// its component}
    \OperatorTok{.}\NormalTok{len()}\OperatorTok{;}
\KeywordTok{let}\NormalTok{ local\_density }\OperatorTok{=}\NormalTok{ graph\_view}\OperatorTok{.}\NormalTok{triangle\_count\_around(device\_node)}\OperatorTok{;}

\KeywordTok{let}\NormalTok{ score }\OperatorTok{=}\NormalTok{ weight\_v }\OperatorTok{*}\NormalTok{ velocity }\KeywordTok{as} \DataTypeTok{f64}
          \OperatorTok{+}\NormalTok{ weight\_r }\OperatorTok{*}\NormalTok{ (ring\_size }\KeywordTok{as} \DataTypeTok{f64}\NormalTok{) }\OperatorTok{*}\NormalTok{ (}\DecValTok{1.0} \OperatorTok{+}\NormalTok{ local\_density }\KeywordTok{as} \DataTypeTok{f64}\NormalTok{)}\OperatorTok{;}
\end{Highlighting}
\end{Shaded}

The analytical \texttt{COUNT(*)} runs on the interpreter with zonemap
pruning; the graph traversal runs over the CSR view. \textbf{Both see
the transaction that was just committed} --- because they read the same
MVCC snapshot as the ledger.

\section{Why the other architectures struggle
here}\label{why-the-other-architectures-struggle-here}

\begin{itemize}
\tightlist
\item
  \textbf{Warehouse + ETL:} the graph and stats lag the ledger by
  minutes --- useless for in-line scoring.
\item
  \textbf{Separate graph DB:} the scorer must join across two systems
  and reconcile two consistency models; the graph mutation contends with
  the traversal.
\item
  \textbf{Single-writer OLAP (DuckDB):} cannot ingest concurrently with
  the scan; the write stream would have to pause for every analytical
  read.
\end{itemize}

ChakraDB serves all three from one snapshot, in one process, with
writers never blocked by the scorer.

\section{Operating it}\label{operating-it}

\begin{itemize}
\tightlist
\item
  Run the \textbf{scorer} in-line (per transaction) using live adjacency
  (\texttt{out\_neighbors}, \texttt{out\_degree}) for cheap signals, and
  a periodic \texttt{view()} +
  \texttt{pagerank}/\texttt{connected\_components} pass for the
  expensive ring analysis (see \href{../graph/snapshot.md}{Live Graph
  Analytics}).
\item
  Watch \texttt{Storage::stats()} for checkpoint lag and backpressure
  under peak ingest (see
  \href{../getting-started/getting-started.md}{Observability}).
\item
  Back up the store with \texttt{backup\_to} on a cadence (see
  \href{../getting-started/getting-started.md}{Backup \& Restore}).
\end{itemize}

\section{The takeaway}\label{the-takeaway}

Fraud detection is the archetypal \textbf{HTAP + graph} workload: fast
writes, live statistics, and relationship traversal, all needing to
agree on the \emph{same instant}. ChakraDB's contribution is that the
agreement is free --- it is one snapshot clock --- so the application is
a few queries and a graph view, not a three-system integration project.

\part{Perspective}

\chapter{ChakraDB vs.~DuckDB}\label{chakradb-vs.-duckdb}

\epigraph{If you know the enemy and know yourself, you need not fear the result of a hundred battles.}{--- Sun Tzu}

DuckDB is the reference point --- the best-in-class embedded analytical
engine. This chapter is honest about where DuckDB wins, where ChakraDB
wins, and the one axis where the comparison is \emph{structural} rather
than a matter of tuning.

\section{The structural difference: concurrent
writers}\label{the-structural-difference-concurrent-writers}

DuckDB holds a \textbf{single-writer file lock}. A second process
opening the same database for writing is refused at the OS level
(\texttt{IO\ Error:\ Could\ not\ set\ lock}). That is not a tuning
parameter --- it is the concurrency model.

ChakraDB permits \textbf{continuous concurrent writers with non-blocking
snapshot reads}. This is the wedge: ChakraDB serves a workload DuckDB
\emph{structurally cannot} --- a live write stream with analytics
running over it at the same time.

\begin{center}\includegraphics[width=0.92\linewidth,height=0.80\textheight,keepaspectratio]{diagrams/diagram-37.pdf}\end{center}

\textbf{Measured:} under 4 threads issuing durable, WAL-logged writes,
ChakraDB's analytical query p50 degrades just \textbf{1.49×} (2.2 → 3.3
ms) while \textasciitilde8,900 writes commit \emph{during} the
measurement, readers never block, and every query sees a stable snapshot
(\texttt{wedge-bench}). DuckDB cannot run this benchmark --- the second
writer is refused.

\section{Cold-scan analytics: DuckDB's home
turf}\label{cold-scan-analytics-duckdbs-home-turf}

On a static dataset scanned cold, DuckDB is excellent and mature.
ChakraDB does not try to out-scan it; the goal is to be
\emph{competitive} while adding concurrency. On a 105-column
ClickBench-shaped table, identical results to DuckDB, p50 ms:

\begin{longtable}[]{@{}lrrl@{}}
\toprule\noalign{}
Query & 10M --- ChakraDB & 10M --- DuckDB & Winner \\
\midrule\noalign{}
\endhead
\bottomrule\noalign{}
\endlastfoot
\texttt{COUNT(*)} & 0.0 & 0.0 & tie \\
filtered count (\texttt{\textless{}\textgreater{}}) & 4.8 & 1.0 &
DuckDB \\
sum/avg & 5.3 & 3.0 & DuckDB \\
\texttt{COUNT(DISTINCT\ UserID)} & \textbf{40.6} & 60.0 &
\textbf{ChakraDB 1.5×} \\
\texttt{COUNT(DISTINCT\ SearchPhrase)} & \textbf{12.0} & 36.0 &
\textbf{ChakraDB 3×} \\
min/max & 0.1 & 0.0 & tie (both instant) \\
group-by adv engine & 4.9 & 2.0 & DuckDB \\
top users (group+sort) & \textbf{59.7} & 71.0 & \textbf{ChakraDB} \\
top phrases / by time & 55.7 / 52.1 & 24.0 / 15.0 & DuckDB \\
pk range \textasciitilde20 rows & 1.1 & 1.0 & tie (both prune) \\
\end{longtable}

\textbf{Reading it:} - \textbf{ChakraDB wins the big
\texttt{COUNT(DISTINCT)}s at scale} (up to 3×) and top-users --- and
those margins \emph{widen} as rows grow. - \textbf{DuckDB wins simple
\texttt{GROUP\ BY} and top-K} --- its vectorized hash-agg and top-K
operators are more optimized, and zonemaps don't help those shapes. -
\textbf{Selective range scans are a tie} --- both prune (DuckDB via
rowgroup stats, ChakraDB via zonemaps); a needle range stays
\textasciitilde1 ms as the table grows 100×.

The full harness (both engines, identical CSV) is
\texttt{src/bin/clickbench.rs} + \texttt{scripts/clickbench\_duckdb.sh};
the numbers and method are in \href{comparisons.md}{Benchmark
Methodology}.

\section{Feature and model
comparison}\label{feature-and-model-comparison}

\begin{longtable}[]{@{}
  >{\raggedright\arraybackslash}p{(\linewidth - 4\tabcolsep) * \real{0.3333}}
  >{\raggedright\arraybackslash}p{(\linewidth - 4\tabcolsep) * \real{0.3333}}
  >{\raggedright\arraybackslash}p{(\linewidth - 4\tabcolsep) * \real{0.3333}}@{}}
\toprule\noalign{}
\begin{minipage}[b]{\linewidth}\raggedright
\end{minipage} & \begin{minipage}[b]{\linewidth}\raggedright
DuckDB
\end{minipage} & \begin{minipage}[b]{\linewidth}\raggedright
ChakraDB
\end{minipage} \\
\midrule\noalign{}
\endhead
\bottomrule\noalign{}
\endlastfoot
Model & embedded OLAP & embedded \textbf{HTAP + graph} \\
Concurrent writers & ❌ single-writer lock & ✅ MVCC, non-blocking
reads \\
Transactions & ✅ & ✅ (BEGIN/COMMIT/ROLLBACK, first-committer-wins) \\
Analytical engine & native vectorized & DataFusion (vectorized) \\
Graph & ❌ & ✅ built-in (adjacency + algorithms) \\
On-disk format & native (open via extensions) & Arrow IPC parts
(open) \\
Exact \texttt{DECIMAL} & ✅ & ✅ (i128 / Arrow Decimal128) \\
Backup/restore & file copy & \texttt{backup\_to} / \texttt{restore}
(consistent) \\
Maturity & very high & young, but crash-tested \\
\end{longtable}

\section{When to pick which}\label{when-to-pick-which}

\begin{itemize}
\tightlist
\item
  \textbf{Pick DuckDB} for the fastest cold-scan analytics over data you
  loaded earlier, on a laptop-class machine, with a single writer.
\item
  \textbf{Pick ChakraDB} when writes and analytics happen \emph{at the
  same time} on live data, when you also need \textbf{graph} traversals
  over that data, or when concurrent ingest must never pause for a
  reader.
\end{itemize}

They are not the same product. DuckDB optimizes the
\emph{static-analytics} axis; ChakraDB optimizes the \emph{live,
concurrent, multi-model} axis.

\backmatter
\end{document}
